%% ============================================================
%% PAPER 9 (of 11) -- cod_certification -- Part I: CONSTRUCTIVE (E2)
%% v31 (2026-09-29). Derived from paper09_cod_certification_v30.tex.
%%
%% WHY THIS FILE WAS SPLIT: v30 was TWO complete LaTeX documents concatenated.
%%   It contained 2x \documentclass, 2x \begin{document}, 2x \end{document},
%%   2x \maketitle. \part{Arv} sat at char 116612, AFTER the first
%%   \end{document} at char 116595 -- so LaTeX ignored everything past it and
%%   COMPILING v30 PRODUCED PART I ONLY. 46,127 chars (a third of the file) and
%%   38 labels were dead text, including all 14 labelled results.
%%
%% WHY SPLIT RATHER THAN MERGE (content and merits, per the user's criterion):
%%   Part I  : Schaefer 11, LRP 87, surplus production 3, exact rational 1,
%%             harvest-free 1, breach 2, labelled results 0.
%%   Part II : Schaefer 0,  LRP 0,  surplus production 0, exact rational 10,
%%             harvest-free 17, breach 27, labelled results 14 (ALL of them).
%%   Shared bibliography = data sources only (Regular 2025; Schijns 2021), not
%%   literature. Part II never cites Part I (0 mentions). Part I is a FITTED-
%%   MODEL floating-point empirical study; Part II is an EXACT-RATIONAL study
%%   with NO fitted model. Not one paper.
%%
%% SIBLING: Part II is paper09b_arv_certification_v1.tex (source: arv).
%% NOTE: 'arv' is a separate unit in the authoritative eight-paper scope.
%%
%% THIS HALF IS UNCHANGED IN SUBSTANCE except: (a) new \ref{priorart} section
%%   inserted before subsection 2, (b) three verified bibliography entries added
%%   (Butterworth 2007; Punt et al. 2016; Regnier and De Lara 2015).
%% ============================================================

\documentclass[11pt]{article}
\usepackage[a4paper,margin=1in]{geometry}
\usepackage{amsmath,amssymb}
\usepackage{graphicx}
\usepackage{booktabs,longtable,array,calc}
\usepackage{xcolor}
\usepackage[colorlinks=true,allcolors=blue!45!black]{hyperref}
\usepackage[font=small,labelfont=bf]{caption}
\usepackage{etoolbox}
\AtBeginEnvironment{longtable}{\small}  % dense data tables
\setlength{\tabcolsep}{4pt}
\providecommand{\tightlist}{\setlength{\itemsep}{0pt}\setlength{\parskip}{0pt}}
\setcounter{secnumdepth}{-1}
\emergencystretch=3em
\graphicspath{{../}}
\usepackage{amsthm}
\newtheorem{theorem}{Theorem}
\newtheorem{lemma}[theorem]{Lemma}
\newtheorem{proposition}[theorem]{Proposition}
\newtheorem{corollary}[theorem]{Corollary}
\newtheorem{remark}[theorem]{Remark}

\begin{document}
\title{Certification, not simulation: the viability of harvest rules on two real resource systems, and the horizon of what can be certified}
\begin{abstract}
\noindent\textbf{The general claim.} Reference points and management rules are usually defended by
\emph{simulation}: a rule is run, the trajectory is inspected, and the rule is recommended. This
paper defends them by \emph{certification} instead, on two real resource systems, and argues that
the difference is not stylistic. A simulation shows what happened on the runs that were tried; a
certification states \emph{why}, locates the verdict in a constant, and --- most importantly ---
makes explicit \emph{the horizon over which the verdict holds}. That last point is the paper's
central finding, and it recurs across two systems and independent pipelines: \emph{certification has
a horizon, and the horizon is set by the fitted map, not by the method.} On the Northern cod stock
the certified layer has a finite horizon --- six years under the two harsher floors, seven under the
informative one --- because the map is expansive at the reference point for every admissible carrying
capacity. On the Edwards Aquifer, every positive-pumping certified kernel is empty beyond three
years, an optimistic bound on a defect the out-of-sample audit exceeds. Two systems, different
hydrologies, different institutions, the same shape of limitation. A certified verdict is therefore
not a stronger version of a simulated one; it is a different kind of object, whose validity is
\emph{local in time} and whose locality can be quantified.

\noindent\textbf{Why three parts, two systems.} Certification is not one operation, and the three
parts do three different ones on two real assessment records. Part I \emph{characterises}: the
robustness of a single-threshold reference point to persistent productivity shocks reduces to three
constants, and the kernel computation collapses to algebra. Part II \emph{certifies a historical
event}: it certifies, in exact rational arithmetic, the obstruction structure that the Northern cod
collapse exhibited. Part III \emph{scores a policy family}: it computes robust viability kernels for
competing governance rules on the Edwards Aquifer and retains a rule only if it meets a fixed
criterion. Parts I and II are on Northern cod; Part III is on the Edwards Aquifer. The two systems
have no hydrology, no biology and no institutional history in common --- which is precisely why the
recurrence of the horizon limitation is a finding rather than a property of one fit.

\noindent\textbf{What is found.} Part I: on the increasing branch of the surplus curve, the largest
catch holding the reference point from itself is $C^* = g(\mathrm{LRP}) - |e| = 91.59$ kt, and any
rule's protection margin is exactly $C^*$ minus its catch there --- at that point harvest and
protection are one budget. A second constant, $g_{\max} - |e| = 215.2$ kt, separates viable-but-raised
kernels from empty ones, and the two reproduce the kernel table in closed form. That identity makes
the no-dominance verdict \emph{structural, not empirical}: a reactive rule cannot out-supply a flat
cap it matches in protection. Part II: the certified object is a harvest-free multiplier bracket,
$[\rho,\ \max\{\rho+r,\ \rho/(1-r)\}]$; both collapse steps certify as harvest-free contractions
(upper bounds $0.754$ and $0.372$) under every timing convention, while the reference window's worst
step and the two prelude steps do not certify --- a typology, not a blanket verdict. Part III: at the
618-ft threshold under the drought-of-record floor, training-mean pumping fails within roughly
thirteen years while an interpolated $7.2\%$ mean cut secures the threshold; the reactive rules are
retained only \emph{nominally}, on a hybrid criterion whose wet-year margin is exactly what the
robust class excludes; and nothing is retained at 660 ft, where the rules are invisible to the
kernel of their own trigger.

\noindent\textbf{Structure.} Part I is the harvest--protection budget, Part II the regime-viability
certification, and Part III the Edwards policy-family scoring. Each is self-contained and reproduced
in full from its companion paper without condensation, opening with that paper's own abstract. A
cross-system conclusion follows, stating what certification buys, what it costs, and where its
horizon binds.
\end{abstract}
\maketitle

\tableofcontents

\newpage

\part*{Part I --- The harvest--protection budget at the limit reference point}
\label{part:budget}
\addcontentsline{toc}{part}{Part I --- The harvest--protection budget}

\begin{abstract}
Reference points are usually defended by simulation. Here the defence is a
characterisation: on the increasing branch of the surplus curve, the
robustness of a single-threshold reference point to persistent productivity
shocks reduces to three constants, and the kernel computation collapses to
algebra. Applied to Northern cod (NAFO 2J3KL; Schaefer fit 1983--2007; LRP
\(884.6\) kt) under a protocol frozen before any score was computed: (1) the
largest catch holding the reference point from itself is
\(C^* = g(\mathrm{LRP}) - |e| = 91.59\) kt, and any rule's protection margin
is exactly \(C^*\) minus its catch there --- at that point harvest and
protection are one budget; a second constant, \(g_{\max} - |e| = 215.2\) kt,
separates viable-but-raised kernels from empty ones, and the two reproduce the
kernel table in closed form. (2) That identity makes the no-dominance verdict
structural, not empirical: a reactive rule cannot out-supply a flat cap it
matches in protection. (3) The map is expansive at the reference point for
every admissible \(K \ge 2K^*\), so the certified layer has a finite horizon
--- six years under the two harsher floors, seven under the informative one
--- which is the crossing at which a geometrically growing erosion margin
overtakes the worst-case trajectory, and which no catch reduction
extends: it is the same for every admissible catch, including none at all.
(4) Carrying capacity is not identified
from above, but the functionals the results depend on are: across the profile
set \(C^*\) lies in \([67.9, 95.2]\) kt --- pinned to within
\(27\) kt while \(K\) itself is not identified from above ---
whereas resampling the \(24\) transitions and refitting both
parameters freely disperses the estimate over a far wider range
(median \(73.7\) kt, \(90\%\) interval \([-89.4, 125.7]\), or
\(88.1\) kt (\([-5.6, 130.6]\)) restricted to the expansive regime on which the
results are conditional, \(73\%\) of replicates). The two answer
different questions and the tight one is an identification result, not a
precision one; 20-year survival from
the reference point falls from \(0.91\) to \(0.65\) between moratorium-level
and \(120\) kt removals. (5) Self-viability is regime-dependent: \(171\) kt
on the post-moratorium window, \(0 \pm 8\) kt on the modern series. On this
map the reference point is protected by good years, and reactive design
cannot buy protection with supply.

\end{abstract}

\textbf{Keywords:} northern cod; NAFO 2J3KL; surplus production;
viability; robust viability kernel; harvest control rules; limit
reference point

\subsection{1. Introduction}\label{budget-introduction}

After the 1992 moratorium, Northern cod (NAFO divisions 2J3KL) posed a
clean governance question. Can any catch policy hold a collapsed stock's
spawning biomass above its limit reference point when productivity is
depressed? Or is the reference point instead protected by good years
rather than by demand management? Viability analysis supplies the
natural instrument for such questions. It asks which states admit
controls that keep every constraint satisfied over time, rather than
which policy maximizes an objective (Aubin, 1991). It has been applied
to fisheries in exactly this register: bioeconomic viability constraints
(B\'en\'e, Doyen, and Gabay, 2001); viable recovery paths after collapse
(Martinet, Th\'ebaud, and Doyen, 2007); stochastic co-viability for
ecosystem-based management (Doyen et al., 2012); and the numerical
computation of viability kernels for fishery case studies (Krawczyk,
Pharo, Serea, and Sinclair, 2013; Krawczyk and Pharo, 2013). The
reference point is a floor on the productive stock itself, namely the
spawning biomass whose regeneration underwrites future yield. The
question is therefore whether adjustments to the draw on that stock can
hold the base that regenerates the yield, or whether the reference point
is protected only by favourable productivity. This study answers the
question by scoring rather than assertion.

Four questions organize the results. (Q1) Under which productivity
regimes does any catch policy --- zero catch included --- hold the LRP,
and where is the question vacuous because the floor exceeds the map's
maximum surplus? (Q2) Under the informative regime, what is the largest
constant catch whose worst-case equilibrium remains at or above the LRP,
and how do reactive rules fare against it under the declared dominance
rule? (Q3) How fast does model error of the declared magnitude erode the
certified safety margin when the governed map is expansive at the
reference point? (Q4) Which answers are geometric properties of the rule
and the threshold, which are identified properties of the fitted map,
and which are sensitive to the model form or the carrying capacity?

The object is the governed surplus-production model fitted in a
companion forecast-evaluation study (under separate review), which ended
in a negative certificate with last-value persistence unbeaten on
out-of-sample RMSE. Here the same fitted model, on the primary
assessment specification (the 1983--2015 NCAM M-shift SSB series of DFO,
2016, Table A2, with LRP \(884.6\) kt), is scored as a closed-loop
governance object. The deliverables are robust viability kernels of the
LRP for a declared catch-policy family under persistent
productivity-shock floors, plus a certified layer that converts the
declared model defect into a safety-margin erosion. The intervention
protocol --- object, defect declaration, disturbance classes, policy
family, and the selection rule restated in Definition 2.6 as a dominance
partial order --- was frozen (dated 2026-08-26) before any kernel,
boundary, replay, or selection score was computed. The frozen protocol
document is archived with the analysis code. Section 2 states the
methods. Section 3 reports the kernels, the no-dominance verdict and the
vacuous-class identity, the constructive boundary, the certified layer,
and the stress replay, followed by the form comparison (Section 3.6),
the carrying-capacity, stochastic, finite-duration, and uncertainty
layers (Sections 3.7-3.10), and the second-specification row (Section
3.11). Section 4 discusses limitations.

\subsection{Related work}
\label{budget-priorart}

Three literatures meet here: the viability analysis of managed renewable resources, the
robust-viability treatment of uncertainty, and the evaluation of harvest control rules.

\textit{Viability in fisheries.} That the governing question for a managed renewable resource
is not which policy optimises but which states admit controls that keep every constraint
satisfied over time is the founding move of viability theory (Aubin, 1991). It has been taken
up in fisheries repeatedly: bio-economic viability, with the time-of-crisis function as an
overexploitation indicator and irreversible configurations related to resource extinction
(B\'en\'e, Doyen and Gabay, 2001); viable recovery paths toward sustainable fisheries
(Martinet, Th\'ebaud and Doyen, 2007); stochastic viability for ecosystem-based management
under multiple conflicting objectives (Doyen et al., 2012); and viability as a
``satisficing'' alternative to optimal control, whose admissible policies are generically
multiple and depend on fewer parameters than optimal or equilibrium strategies (Krawczyk and
Pharo, 2013; Krawczyk et al., 2013). This paper sits in that lineage and uses its vocabulary
throughout.

\textit{Robust viability kernels.} The set whose computation is collapsed here is the robust
viability kernel: the set of initial biomasses from which at least one harvesting strategy
guarantees the required production and preservation levels for all times, \emph{whatever} the
uncertainties, these being disturbances taking values in a known given set (Regnier and De
Lara, 2015). Their central finding --- that accounting for uncertainty sensibly shrinks the
deterministic viability kernel --- is the phenomenon this paper is organised around, and
nothing here contradicts it.

What differs is how the set is obtained. In that literature the kernel is computed
numerically, by a descending set iteration on a discretised state space. On the increasing
branch of the surplus curve the computation here \textbf{collapses to algebra} instead: two
constants, $C^{*} = g(\mathrm{LRP}) - |e| = 91.59$~kt and $g_{\max} - |e| = 215.2$~kt,
reproduce the kernel table in closed form, and every rule's protection margin is exactly
$C^{*}$ minus its catch at the reference point. At that point harvest and protection are one
budget. A closed-form collapse is not a faster numerical method; it is a different kind of
result, and it is what makes the no-dominance verdict below structural rather than empirical.

\textit{Harvest control rules and management strategy evaluation.} The standard instrument for
choosing among harvest control rules is management strategy evaluation: simulate the entire
management system --- operating model, observation, assessment and rule --- under specified
uncertainties, and compare candidate rules on pre-specified performance metrics (Butterworth,
2007; Punt et al., 2016). It is a large, mature and practically indispensable machinery. The
opening remark of this paper, that reference points are usually defended by simulation, is a
description of that literature rather than a criticism of it.

The distinction is in what the two deliver, and it can be stated without rhetoric. An MSE
verdict is \emph{conditional on the operating model and on the ensemble simulated}: a
no-dominance finding holds for the rules, the uncertainties and the replicates that were
actually run. The no-dominance verdict here falls out of an identity --- protection margin
equals $C^{*}$ minus catch --- and therefore holds for \emph{every} rule in the class at
once, including rules that were never simulated. That is a genuine difference in kind.

It is also a narrower one, and the narrowing should be conceded. MSE routinely handles
multiple conflicting objectives, implementation error, observation error and institutional
constraints simultaneously; the scalar reduction here buys its closed form by giving most of
that up. The two are complements rather than rivals: the characterisation says what cannot be
bought at any price within the class, while the simulation says which rules do best on the
metrics that were specified.

\subsection{2. Methods}\label{budget-methods}

\subsubsection{2.1 Operating model and
data}\label{budget-operating-model-and-data}

All inputs are public and the analysis is reproducible: every stochastic
layer uses a fixed seed and a fixed draw budget (Sections 3.8 and 3.10),
and every deterministic layer is a closed-form or fixed-point
computation. The stock
series is the Northern cod (NAFO 2J3KL) spawning-stock biomass of DFO
(2016), Table A2. Annual landings are Schijns et al.~(2021), Table 1.
The fit window is 1983--2007 (24 one-step transitions; the out-of-sample
audit runs 2008--2015), and the protocol was frozen on 2026-08-26,
before any kernel, boundary, replay, or selection score was computed.
The training-window residuals of the fitted map have mean \(-10.9\) kt,
SD \(114.9\) kt, range \([-329.0, +206.6]\) kt, and lag-1
autocorrelation \(0.55\). The 1992 transition residual is the minimum
(\(-329.0\) kt). The training-window maximum stock is \(940.75\) kt,
which is why the fitted carrying capacity, pinned at its \(5000\) kt
optimization bound, sits above the observed range (the declared
optimization box runs from just above the training maximum to \(5000\)
kt). A companion forecast-evaluation study scored the same
point-estimate map against persistence on the same series and issued a
negative certificate (persistence unbeaten out of sample). That
certificate is scored there, not here, and the closed-loop object below
is the same fitted map by construction.

All analyses beyond the frozen kernel tables --- the carrying-capacity
grid, the stochastic viability layers, the finite-duration floors, and
the bootstrap bands of Sections 3.7--3.10 --- are additional scored
objects executed after the freeze on the declared machinery. None of
them replaces or alters the frozen family, the frozen floor classes, or
the dominance rule of Definition 2.6.

\subsubsection{2.2 The governed object and the scoring
rule}\label{budget-the-governed-object-and-the-scoring-rule}

The governed object is the companion forecast-evaluation study's own
stock-flow class.

\[S_{t+1}=\bigl[S_t+rS_t(1-S_t/K)-C_t+e_t\bigr]_+\]

The Allee term is off. The map is fitted by one-step least squares on
1983--2007 with annual landings (Schijns et al., 2021), giving
\(r = 0.2369\) and \(K = 5000\) kt (pinned at its optimization bound;
the series never approaches carrying capacity, a declared defect; the
LRP-boundary results depend chiefly on the identified \(r\)). Fit
residual SD is \(114.9\) kt. The defect declaration is
\(\varepsilon = 329.0\) kt yr\textsuperscript{-1} (the 1992 collapse transition). The
out-of-sample audit on 2008--2015 gives a maximum of \(47.1\) kt, which
does not exceed the declared defect (by contrast, the groundwater object
of the companion intervention study exceeds its declared defect out of
sample).

The declared formal objects are stated as definitions and propositions
so that the empirical content can be read off cleanly.

\textbf{Definition 2.1 (Governed surplus-production map).} The governed
object is the one-step map above with parameters
\((r, K) = (0.2369, 5000\text{ kt})\) fitted by one-step least squares
on 1983--2007, Allee term off, and the declared defect declaration
\(\varepsilon = 329.0\) kt yr\textsuperscript{-1}.

\textbf{Definition 2.2 (Safe set).} The single declared threshold is
\(K^* = \mathrm{LRP} = 884.6\) kt (the 1983--1989 mean of Table A2). The
safe set is \([K^*, \infty)\). No row on the second assessment
specification (the 1954--2024 extended xteNCAM series, LRP \(276\) kt)
enters the frozen primary protocol. The labelled post-freeze row of
Section 3.11 reports that specification separately, unpooled, against
its own reference point, and no verdict transfers between the two.

\textbf{Definition 2.3 (Disturbance classes).} The disturbance classes
are persistent additive productivity floors from the fit-window residual
distribution: (i) the perpetual worst observed one-step shock
(\(-329.0\) kt yr\textsuperscript{-1}), (ii) the 5th-percentile class (\(-287.4\) kt
yr\textsuperscript{-1}), and (iii) the 10th-percentile class (\(-80.9\) kt yr\textsuperscript{-1}). Because
this object has no independent input channel (unlike the recharge series
of the companion groundwater study), the disturbance classes and the
defect declaration are the same measured quantity in two roles.

\textbf{Definition 2.4 (Governance family).} The governance family
consists of the following declared policies. BAU has \(C \equiv 5\) kt
(moratorium-level inshore removals, the declared implementable use
post-1992). Flat caps are \(\rho \cdot 240\) kt with
\(\rho \in \{1.0, 0.75, 0.5, 0.25, 0.0\}\)
(i.e.~\(240, 180, 120, 60, 0\) kt; every member is scored). S1 is the
DFO-2009 critical-zone rule (DFO, 2009) at a declared \(60\) kt cap
(\(60\) above the LRP, \(0\) below). The cascade is \(60/30/5/0\) kt at
LRP/\(0.75\cdot\)LRP/\(0.5\cdot\)LRP/below; the sub-LRP stages are
declared scenarios, not verified institutions. Two new families are
reactive in the strict sense --- their catch is a function of the
current stock, switching above the LRP and vanishing at or below it.
Family A is surplus-proportional: \(C(S) = \phi\, g(S)\) for \(S \ge\)
LRP and 0 below, with \(\phi \in \{0.25, 0.50, 0.75\}\); its catch
scales with the harvested surplus \(g(S) = rS(1 - S/K)\) and therefore
feeds back on the stock. Family B is graded above the LRP: graded2 is 0
/ LRP--\(1.25\cdot\)LRP: 60 / \textgreater{}\(1.25\cdot\)LRP: 90 kt, and
graded3 is 0 / LRP--\(1.15\cdot\)LRP: 30 /
\(1.15\cdot\)LRP--\(1.35\cdot\)LRP: 60 / \textgreater{}\(1.35\cdot\)LRP:
90 kt. Families A and B are post-freeze scored objects added to address
the reactive-management question directly; they are not declared
institutions.

Two further reactive designs were scored and are not reported as a
family because they are dominated. The surplus-harvest fraction
\(C(S) = \alpha\cdot\max\bigl(0, g(S) - g(K^*)\bigr)\) removes only the
surplus above the LRP's own regeneration, and the reactive ramp
\(C(S) = \min(C_{\max}, q\,(S - K^*))\) ramps linearly above the LRP.
Under the declared source-year classes, every member of both was empty
at the 5th-percentile \(T=\infty\) reading (the clause-H1 failure of
Result 3.1 applies to them as to every positive-catch rule) and, on the
training-supply replay, harvested at or below the moratorium level
(\(0.18\)--\(0.71\) kt mean catch for the surplus-harvest fractions,
\(0.64\) kt for the ramp; \(5\) kt for BAU), so neither family achieves
the ``harvest more than the moratorium'' property that lets Family A
register a genuine trade-off. They are recorded as rejected designs and
omitted from the kernel tables below.

\textbf{Definition 2.5 (Robust viability kernel, closed-loop reading).}
For each policy and disturbance class, the robust viability kernel is
the set of initial SSB values from which the closed-loop path keeps the
stock at or above the LRP under the persistent floor. Table 1 reports
kernel lower boundaries at \(T=1\) and \(T=\infty\), with ``empty''
denoting that no state is robustly viable. The term is used in its
closed-loop reading --- the robust positively invariant set of a fixed
declared policy --- not in the classical viability reading of Aubin
(1991), which is existential over controls. No control choice enters the
kernel computation; only the evaluation of declared rules under the
disturbance floor enters. The word is retained because the object
answers the same question (``from which states can the constraint be
held'') in the policy-fixed form.

\textbf{Definition 2.6 (Dominance partial order).} The rule the frozen
protocol calls its retention rule is restated here, clauses verbatim, as
a partial order on policies rather than a filter that selects an adopted
policy: a non-BAU policy \emph{dominates} BAU only if all three of the
following hold.

\begin{itemize}
\tightlist
\item
  (H1) Its kernel is at least as protective as BAU's at every reading
  (compared on the kernel lower boundary; empty = worst).
\item
  (H2) It improves on BAU somewhere.
\item
  (H3) At some reading where it improves, its mean allowed catch exceeds
  that of every at-least-as-protective flat cap.
\end{itemize}

\emph{Why the filter vocabulary is not used.} Clause (H1), read at every
reading with empty scored as worst, is a structural block rather than a
comparison that could go either way: under the 5th-percentile class at
\(T=\infty\) every declared positive-catch rule has an empty kernel
while BAU's is nonempty (\(2219.6\) kt), so within the declared family
no positive-catch rule can satisfy (H1) at all --- the rule's outcome is
fixed by the declared classes and BAU's marginal nonemptiness there.
Presenting that outcome as a ``retention'' decision would overstate it;
the vocabulary here is dominance, the outcome is stated as a selection
property of the declared partial order (Result 3.1), and the two
identity-level statements of the construction (the vacuous-class
identity and the structural block itself) are stated as definitional
notes.

\textbf{Definition 2.7 (Certified layer conversion).} The certified
layer applies the defect-to-margin conversion in the form the fitted map
admits. If the closed loop contracts with rate \(a < 1\) on the safe
domain, the erosion margin is \(r_T = \varepsilon(1-a^T)/(1-a)\).
Otherwise the expansive form
\(r_T = \varepsilon(a_{\max}^T - 1)/(a_{\max} - 1)\) applies with
\(a_{\max} = \sup|F'|\) over the safe domain, and the certified kernel
is the nominal kernel of \(K^* + r_T\).

\subsubsection{2.3 The catch-timing convention}\label{budget-the-catch-timing-convention}

One modelling choice has to be declared explicitly, because it moves the
reported numbers. Write the one-step residual of the governed map as

\[e_t = S_{t+1} - \bigl(S_t + g(S_t) - C_{t+\delta}\bigr),\qquad
\delta \in \{0, 1\},\]

so that \(\delta = 0\) pairs the transition \(t \to t+1\) with the catch
of the year the transition leaves (\emph{source-year}) and \(\delta = 1\)
pairs it with the catch of the year it enters
(\emph{destination-year}). The two are different residual distributions
on one and the same fitted map: at the committed \((r, K)\) the
source-year pool has mean \(-10.9\) kt, SD \(114.9\) kt and lag-1
autocorrelation \(0.554\), against \(-20.4\) kt, \(135.0\) kt and
\(0.652\) for the destination-year pool, and the classes of Definition 2.3
move with them (\(-329.0\)/\(-287.4\)/\(-80.9\) kt against
\(-460.0\)/\(-318.8\)/\(-114.9\) kt).

\textbf{Convention (adopted).} Every kernel, boundary, floor, replay,
figure and table of this paper is computed in the source-year convention
(\(\delta = 0\)).

It is adopted on three grounds, none of them provenance. \emph{(i)
Coherence.} The committed parameters are the minimisers of the
source-year loss. The estimation routine regresses the increment
\(S_{t+1} - S_t\) on the catch of year \(t\), so \(r = 0.2369\) and
\(K = 5000\) kt minimise the mean squared source-year residual:
\(12{,}772\) kt\(^2\), against \(17{,}873\) kt\(^2\) for the
destination-year pool at those same parameters (\(28.5\%\) lower).
Scoring those parameters against destination-year residuals would report
classes drawn from a loss the fit never minimised. \emph{(ii) Fit.} The
source-year pool is the better-fitting and the less serially correlated
of the two, as tabulated above. \emph{(iii) No interior alternative.}
Refitting \((r, K)\) freely against the blended catch
\(w\,C_t + (1-w)\,C_{t+1}\) gives a mean squared error monotone
decreasing in \(w\) over the whole unit interval, with the optimum at the
boundary \(w = 1\): \(17{,}713\) kt\(^2\) at \(w = 0\) (where the refit
gives \(r = 0.2084\)), \(15{,}026\) kt\(^2\) at \(w = 0.5\), and
\(12{,}772\) kt\(^2\) at \(w = 1\); the out-of-sample error over
2008--2015 moves the same way (\(834.9\) to \(783.1\) kt\(^2\), eight
one-step transitions). There is no interior optimum, so a half-and-half
compromise is not supported by the data.

\textbf{Limitation, stated because a referee will ask.} Nothing in the
archived source documentation records whether the spawning-stock column of
DFO (2016) Table A2 is measured before or after the calendar year's
removals, and the biology does not settle it either: Northern cod spawn in
March--May, so an SSB observation sits mid-year while landings are
calendar-year totals, and the interval between two SSB observations spans
the back half of one year's catch and the front half of the next. On that
reading neither pure convention is mechanistically exact and the truth is
a blend --- which is exactly why ground (iii) is needed. The convention is
therefore declared on statistical and coherence grounds rather than
derived from provenance, and it is stated here rather than left implicit.
The sensitivity is not suppressed: the primary kernels and their boundary
tables were also computed under \(\delta = 1\), and those artifacts are
archived with the superseded runner rather than overwritten, so every
affected number can be recomputed from the archive instead of taken on
trust. Four quantities move materially ---
the constructive bound (\(91.59\) kt against \(57.61\) kt), the number of
vacuous floor classes (one against two), the certified horizon (\(6\)/7
years against \(5\)/6), and the 60-kt rules' \(T=\infty\) boundary
(\(884.6\) kt against \(900.3\) kt). No parameter estimate, no form
comparison and no certificate direction changes.

\subsubsection{2.4 What the kernel problem reduces
to}\label{budget-what-the-kernel-problem-reduces-to}

The kernels of Section 3 are computed by exact backward recursion on the
declared interval engine, but they are not merely numerical. On the increasing
branch of the surplus curve --- the only branch the safe set occupies here,
since \(K^* < K/2\) --- the whole table reduces to three constants. Write
\(g\) for the surplus, \(|e|\) for the persistent floor, \(C(S)\) for the
policy's catch and \(F(S) = S + g(S) - C(S) - |e|\) for the worst-case closed
loop. Every statement below is verified cell by cell against the committed
kernels by \texttt{campaign\_e2\_structure\_v3.py} (23 checks).

\textbf{Proposition 2.1 (Constructive bound and the protection--supply
budget).} For any policy whose catch at the reference point is \(C(K^*)\),
the safe set is held from itself if and only if
\(C(K^*) \le C^* := g(K^*) - |e|\), and the margin by which it is held is
exactly \(C^* - C(K^*)\).

\emph{Proof.} One step from \(K^*\) gives
\(K^* + g(K^*) - C(K^*) - |e|\), which is at least \(K^*\) under exactly the
stated inequality, and the margin is the difference of the two sides. It
remains only that the reference point is the binding state, i.e. that
\(S \mapsto S + g(S) - C(S)\) attains its minimum over the safe set at
\(K^*\): this holds on the increasing branch for every declared policy, the
graded rules included, because their catch is piecewise constant and each
upward jump (\(30\) kt at most) is smaller than the surplus gained across the
jump. The campaign verifies it directly for all twelve policies.
\ensuremath{\square}

The largest catch that holds the reference point from itself is therefore
\(C^* = 172.46 - 80.87 = 91.59\) kt (Section 3.3). The identity is the
substantive part: at the reference point harvest and protection are one
budget, and every kilotonne taken there is a kilotonne of margin surrendered.
Section 3.2 shows that this, and not the tabulation, is why no declared rule
dominates the moratorium.

\textbf{Proposition 2.2 (Perpetual bound and the infinite-horizon
boundary).} For a constant catch \(C\), the \(T=\infty\) kernel is
\([b_\infty, \infty)\) with \(b_\infty = \max(K^*, s_-(C))\), where
\(s_-(C)\) is the smaller root of \(g(S) = C + |e|\); it is empty if and only
if \(C > C_{\mathrm{vac}} := g_{\max} - |e|\).

\emph{Proof.} \(F(S) - S = g(S) - C - |e|\) is positive exactly between the
two roots of \(g(S) = C + |e|\). From any state above the smaller root the
trajectory increases towards the larger and so never falls below the smaller;
from any state below it the trajectory declines. If \(C + |e| > g_{\max}\)
there is no root and the trajectory declines everywhere. \ensuremath{\square}

Two constants therefore organise the last column of Table 1: below
\(C^* = 91.59\) kt the kernel is the whole safe set; between \(C^*\) and
\(C_{\mathrm{vac}} = 296.09 - 80.87 = 215.2\) kt it is a proper half-line
strictly above the reference point; above \(C_{\mathrm{vac}}\) it is empty.
The column is reproduced in closed form to \(0.1\) kt for all six flat caps.

\textbf{Proposition 2.3 (Reactive families).} For \(C(S) = \phi\, g(S)\)
above the LRP and 0 below, the closed loop is
\(F(S) = S + (1-\phi)\,g(S) - |e|\) and there are three regimes: the kernel
is the whole safe set iff \(\phi \le 1 - |e|/g(K^*)\); it is a proper
half-line iff that fails but \(\phi < 1 - |e|/g_{\max}\); and it is empty iff
\(\phi \ge 1 - |e|/g_{\max}\).

\emph{Proof.} Proposition 2.1 with \(C(K^*) = \phi\, g(K^*)\), and
Proposition 2.2 with the effective surplus \((1-\phi)g\). \ensuremath{\square}

Under the declared 10th-percentile floor the two thresholds are \(0.531\) and
\(0.727\). The members \(\phi = 0.25\) and \(0.50\) hold the entire safe
set --- the second by \(5.4\) kt, not by the wide margin an earlier version
of this paper reported (Section 3.4) --- \(\phi = 0.75\) is empty, and the
band \(0.531 < \phi < 0.727\) is the regime in which the rule is viable but
only from above the reference point. No member of the originally declared
family fell in that band, which is why the erroneous criterion went unnoticed;
the family has been extended by \(\phi = 0.60\) so that Table 1 exhibits it.


\begin{figure}[htbp]
\centering
\includegraphics[width=\linewidth]{figs_e2_v3/fig8_frontier.png}
\caption{The two structural results of Section 2.4. (a) The protection--supply
budget of Proposition 2.1: every policy lies on the line, so a rule's harvest
at the reference point and the margin by which it holds that point sum to
\(C^* = 91.59\) kt; the flat 60-kt cap, BAU and \(\phi \le 0.50\) sit on the
protective segment, the 120- and 180-kt caps and \(\phi = 0.75\) below it, and
the dotted line at \(C_{\mathrm{vac}} = 215.2\) kt is where the kernel becomes
empty. (b) The three regimes of Proposition 2.3: the safe-set margin
\((1-\phi)g(K^*) - |e|\), zero at \(\phi = 0.531\), and the nonemptiness
margin \((1-\phi)g_{\max} - |e|\), zero at \(\phi = 0.727\), with the four
declared members marked.}
\end{figure}

\textbf{Proposition 2.4 (The certified horizon is a crossing).} For a
constant catch --- the case the certified layer is computed on --- let
\(S_{\mathrm{hi}}\) be the domain ceiling and
\(r_T = \varepsilon(a_{\max}^T-1)/(a_{\max}-1)\) the erosion margin of
Definition 2.7. The certified kernel at horizon \(T\) is nonempty iff
\(F^T(S_{\mathrm{hi}}) \ge K^* + r_T\), and since \(F^T(S_{\mathrm{hi}})\)
is nonincreasing in \(T\) while \(r_T\) increases, the certified horizon is
the unique crossing
\(T^* = \max\{T : F^T(S_{\mathrm{hi}}) \ge K^* + r_T\}\). It is finite
whenever the loop is expansive.

\emph{Proof.} For a constant catch \(F'(S) = 1 + g'(S) > 0\) throughout
\([K^*, S_{\mathrm{hi}}]\) --- the surplus slope reaches \(-1\) only at
\(S = \frac{K}{2}(1 + 1/r) \approx 13{,}000\) kt, beyond the domain --- so
\(F\) is increasing and the supremum over initial states of the \(T\)-step
minimum is attained at the ceiling; the rest is the definition of the certified
kernel. Finiteness follows because \(r_T\) grows
geometrically in \(T\) while \(F^T(S_{\mathrm{hi}})\) is bounded below (it
converges to the attracting fixed point when one exists, and declines
otherwise). \ensuremath{\square}

Nothing in this subsection uses the Schaefer form. The only properties
invoked are that the surplus is unimodal with its peak above the reference
point and that the safe set lies on the increasing branch, so the constants
travel to any surplus function with those two properties. They do: on the
Fox row of Section 3.6 the formula \(C^* = g(K^*) - |e|\) gives
\(79.1\) kt against the tabulated \(79.05\), with
\(C_{\mathrm{vac}} = 111.2\) kt and \(g_{\max} = 192.0\) kt; on the
two Allee rows it gives \(123.27\) kt (declared \(s_0\)) and
\(115.19\) kt (data-preferred), both reproducing the tabulated constructives
exactly. The Schaefer numbers of Section 3.3 are one instance of a statement
that holds across the forms the paper fits.

The crossing reproduces the computed horizons exactly --- \(T^* = 6\) under
the perpetual-worst and 5th-percentile floors and \(T^* = 7\) under the
informative one, for both BAU and zero catch (Section 3.4). Where the loop has
an attracting fixed point \(S^\dagger\) the crossing is bounded below by
\(\lfloor\log(1 + (S^\dagger-K^*)(a_{\max}-1)/\varepsilon)/\log
a_{\max}\rfloor\), which is attained once the ceiling trajectory has
converged: it gives 7 against 7 under the informative floor, and 4 against 6
under the 5th-percentile floor, where at six years the ceiling trajectory is
still \(1{,}500\) kt above the attractor.

\subsection{3. Results}\label{budget-results}

\subsubsection{3.1 Robust kernels
(nominal)}\label{budget-robust-kernels-nominal}

\textbf{Table 1.} Lower boundaries of the robust viability kernels of
the LRP (kt) under the declared catch-policy family (frozen rules plus
the two post-freeze reactive families) and disturbance classes, in the
source-year convention. ``empty'' denotes that no state is robustly
viable at that horizon. The infinite-horizon emptiness of Section 3.2 is
recorded in the table, not only in prose.

\begin{longtable}[]{@{}
  >{\raggedright\arraybackslash}p{(\columnwidth - 12\tabcolsep) * \real{0.1111}}
  >{\raggedleft\arraybackslash}p{(\columnwidth - 12\tabcolsep) * \real{0.1481}}
  >{\raggedleft\arraybackslash}p{(\columnwidth - 12\tabcolsep) * \real{0.1481}}
  >{\raggedleft\arraybackslash}p{(\columnwidth - 12\tabcolsep) * \real{0.1481}}
  >{\raggedleft\arraybackslash}p{(\columnwidth - 12\tabcolsep) * \real{0.1481}}
  >{\raggedleft\arraybackslash}p{(\columnwidth - 12\tabcolsep) * \real{0.1481}}
  >{\raggedleft\arraybackslash}p{(\columnwidth - 12\tabcolsep) * \real{0.1481}}@{}}
\toprule\noalign{}
\begin{minipage}[b]{\linewidth}\raggedright
Policy
\end{minipage} & \begin{minipage}[b]{\linewidth}\raggedleft
Worst, \(T=1\)
\end{minipage} & \begin{minipage}[b]{\linewidth}\raggedleft
\(T=\infty\)
\end{minipage} & \begin{minipage}[b]{\linewidth}\raggedleft
q05, \(T=1\)
\end{minipage} & \begin{minipage}[b]{\linewidth}\raggedleft
\(T=\infty\)
\end{minipage} & \begin{minipage}[b]{\linewidth}\raggedleft
q10, \(T=1\)
\end{minipage} & \begin{minipage}[b]{\linewidth}\raggedleft
\(T=\infty\)
\end{minipage} \\
\midrule\noalign{}
\endhead
\bottomrule\noalign{}
\endlastfoot
BAU (5 kt) & 1025.5 & empty & 989.0 & 2219.6 & \textbf{884.6} &
\textbf{884.6} \\
flat 240 kt & 1233.5 & empty & 1196.4 & empty & 1014.0 & empty \\
flat 180 kt & 1180.0 & empty & 1143.1 & empty & 961.5 & 1637.8 \\
flat 120 kt & 1126.8 & empty & 1090.1 & empty & 909.3 & 1082.3 \\
flat 60 kt & 1073.8 & empty & 1037.2 & empty & \textbf{884.6} &
\textbf{884.6} \\
S1 / cascade & 1073.8 & empty & 1037.2 & empty & \textbf{884.6} &
\textbf{884.6} \\
flat 0 kt & 1021.1 & empty & 984.7 & 2070.9 & \textbf{884.6} &
\textbf{884.6} \\
Family A, \(\phi\)=0.25 & 1064.7 & empty & 1027.3 & empty &
\textbf{884.6} & \textbf{884.6} \\
Family A, \(\phi\)=0.50 & 1111.3 & empty & 1071.7 & empty &
\textbf{884.6} & \textbf{884.6} \\
Family A, \(\phi\)=0.60 & 1131.2 & empty & 1091.0 & empty & 895.2 &
1074.8 \\
Family A, \(\phi\)=0.75 & 1161.0 & empty & 1119.9 & empty & 920.2 &
empty \\
Family B, graded2 & 1073.8 & empty & 1037.2 & empty & \textbf{884.6} &
\textbf{884.6} \\
Family B, graded3 & 1073.8 & empty & 1010.9 & empty & \textbf{884.6} &
\textbf{884.6} \\
\end{longtable}

Under the 10th-percentile class the moratorium, zero catch, the flat
60-kt cap, the critical-zone rule, the cascade, both graded rules, and
the surplus-proportional family at \(\phi \le 0.50\) all hold the entire
safe set \([884.6, 10^4]\) kt at every horizon. Every higher nonzero cap
lifts the kernel's lower edge above the LRP, and the column is not a
list of accidents: by Proposition 2.2 the boundary is
\(\max(K^*, s_-(C))\) with \(s_-(C)\) the smaller root of
\(g(S) = C + |e_{q10}|\), so the catches divide at
\(C^* = 91.59\) kt (whole safe set), \(C_{\mathrm{vac}} = 215.2\) kt
(nonempty) and beyond it (empty) --- which is exactly what the last column
shows, to \(0.1\) kt, for all six flat caps. Under the perpetual-worst
floor the infinite-horizon kernel is empty for every policy, because the
maximum surplus \(g_{\max} = rK/4 = 296\) kt yr\(^{-1}\) lies below that
floor. Under the 5th-percentile class the infinite-horizon kernel is
nonempty for the moratorium (\(2219.6\) kt) and zero catch (\(2070.9\)
kt) but empty for every positive-cap rule.

The critical-floor axis makes the qualification explicit. At zero catch,
\(\bar e = g_{\max} = 296.1\) kt yr\(^{-1}\) separates vacuous from
informative classes. Only the perpetual-worst floor (\(-329.0\) kt) sits
beyond it and is vacuous. The 5th-percentile (\(-287.4\) kt) and
10th-percentile (\(-80.9\) kt) floors both lie on the informative side,
which is why the constructive boundary of Section 3.3 exists for the
10th-percentile class and the 5th-percentile class now carries
informative, non-vacuous content.

\begin{figure}[htbp]
\centering
\includegraphics[width=\linewidth]{figs_e2_v3/fig1_surplus.png}
\caption{Surplus production of the registered fit with the three persistent floors; only the perpetual-worst floor sits beyond \(g_{\max} = 296.1\) kt yr\(^{-1}\), so only that floor's emptiness is vacuous.}
\end{figure}

\begin{figure}[htbp]
\centering
\includegraphics[width=\linewidth]{figs_e2_v3/fig2_kernel_vs_catch.png}
\caption{Kernel lower boundary versus constant catch under the 10th-percentile class: the \(T=\infty\) boundary leaves the LRP at the constructive \(91.59\) kt.}
\end{figure}

\subsubsection{3.2 The no-dominance verdict and the vacuous-class
identity}\label{budget-the-no-dominance-verdict-and-the-vacuous-class-identity}

The first result concerns dominance; the second is the vacuous-class
identity, stated as a definitional note. Both are scoped to the declared
rule and disturbance classes.

\textbf{Result 3.1 (No dominance).} Under the dominance partial order of
Definition 2.6 (the frozen rule; clauses verbatim), no non-BAU policy
dominates BAU.

\emph{Reason.} The binding failure is clause (H1), read at its stated
scope: ``at every reading,'' with empty scored as worst. Under the
5th-percentile class at \(T=\infty\) every positive-catch rule --- the
flat caps, the critical-zone rule, the cascade, both graded members of
Family B, and every member of the surplus-proportional Family A --- has
an empty kernel, while BAU's kernel is nonempty there (\(2219.6\) kt).
Because empty is worst, none of them is at least as protective as BAU at
that reading, so clause (H1) fails for every positive-catch rule. The
informative 10th-percentile class is not where the verdict is decided:
under that class the flat 60-kt cap, the critical-zone rule, the
cascade, and both graded rules hold the LRP from itself at every horizon
(\(884.6\) kt at \(T=1\) and \(T=\infty\)), exactly matching BAU, and
the surplus-proportional family at \(\phi \le 0.50\) holds it too while
harvesting more than the moratorium --- so the informative class is
where the reactive family \emph{almost} passes, not where it fails.
Clause (H3) is therefore a second, independent failure rather than the
mechanism: the equal-protective flat-60 cap (mean allowed catch \(60\)
kt) exceeds the training means of the critical-zone rule (\(10.0\) kt),
the cascade (\(16.3\) kt), graded2 (\(10.0\) kt), and graded3 (\(5.0\)
kt). At \(\phi = 0.75\) Family A also fails on the informative class
itself, its \(T=\infty\) kernel empty (\(920.2\) kt at \(T=1\) against
\(884.6\); see Result 3.4). No non-BAU policy therefore dominates BAU,
and the mechanism is the empty-at-the-5th-percentile-class reading of
clause (H1), with the supply comparison as a second failure. The
verdict's status is selection-theoretic, not empirical: given the
declared classes, clause (H1) blocks every positive-catch rule by
construction (Definition 2.6's note), so the no-dominance outcome is a
property of the declared partial order and classes as much as of the
policies. The companion groundwater evaluation operates under a
different constraint structure --- not a single fatal floor --- so its
retained reactive rules and this no-dominance outcome are not comparable
selections. \ensuremath{\square}

Which governance architecture justifies its additional structure is
system-dependent. The framework's deliverable is the scored comparison,
not a universal architecture verdict.

\textbf{Definitional note 3.2 (The vacuous-class identity).} Under the
perpetual-worst floor, no catch policy --- zero catch included --- holds
the LRP. This statement is a definitional note rather than a numbered
result: it is an arithmetic identity of the declared class, not an
empirical finding, and it is not numbered on a par with the constructive
boundary or the stochastic layer. Under the 5th-percentile and
10th-percentile classes the statement is not vacuous: those floors lie
below the map's maximum surplus and the classes carry informative
content.

\emph{Statement.} The perpetual-worst floor (\(-329.0\) kt yr\textsuperscript{-1}) exceeds
the map's maximum surplus (\(g_{\max} = 296\) kt yr\textsuperscript{-1}). Every trajectory
therefore declines for every catch, zero included: \(|e| > g_{\max}\)
implies monotone decline for every \(C \ge 0\) --- an identity of the
map's algebra, not a measurement. The 5th-percentile class (\(-287.4\)
kt yr\textsuperscript{-1}) and the 10th-percentile class (\(-80.9\) kt yr\textsuperscript{-1}) both lie
below \(g_{\max}\) and are not vacuous: their kernels carry substantive
content (Table 1, Result 3.3), and the 5th-percentile moratorium kernel
is nonempty (\(2219.6\) kt). The correction of the residual convention
therefore reduces the vacuous family from two classes to one. The
analogy to the companion groundwater institutional certificate is one of
form only; that certificate is institutional, this one is a
floor-above-surplus identity, and the two are not pooled. \ensuremath{\square}

On this map the reference point is protected by good years rather than
by demand management --- but under the corrected class the margin that
good years must supply is smaller than the frozen reading implied.

\subsubsection{3.3 Constructive boundary}\label{budget-constructive-boundary}

The maximal robust flat catch is the largest constant catch whose
worst-case low equilibrium stays at or above the LRP. Under the
10th-percentile class this constructive boundary is \textbf{91.59 kt}.

\textbf{Result 3.3 (Constructive boundary).} Under the 10th-percentile
class, the maximal robust constant catch is \(91.59\) kt yr\textsuperscript{-1}. Under the
5th-percentile and perpetual-worst classes no positive catch is robust.

\emph{Reason.} The value is
\(g(K^*) - |e_{q10}| = 172.46 - 80.87 = 91.59\) kt yr\textsuperscript{-1}. It is
independent of the family scaling and rests only on the identified \(r\)
and the corrected 10th-percentile floor. Under the 5th-percentile floor
the corresponding quantity is \(172.46 - 287.36 = -114.9\) kt, so no
constant catch is robust there; under the perpetual-worst floor the same
holds (\(172.46 - 328.97 = -156.5\) kt). \ensuremath{\square}

Proposition 2.1 gives the bound a second reading: the margin by which
any policy holds the reference point is exactly \(91.59\) kt minus its
catch there, so the bound is the whole budget and harvest and protection
are the same quantity seen from two sides. This is certification
geometry at one declared shock class, not a harvest rule. The stochastic analogue of the bound is reported in
Section 3.8, and its sampling distribution in Section 3.10. The
corrected bound is read as order \(70\)--\(90\) kt, not the
\(40\)--\(60\) kt of the frozen convention.

\textbf{Result 3.4 (Reactive rules above the LRP).} The
surplus-proportional family is genuinely reactive: at \(\phi = 0.25\)
and \(0.50\) it holds the entire safe set under the 10th-percentile
class at every horizon, including \(T = \infty\), and its catch feeds
back on the stock (it rises with the surplus and vanishes below the
LRP). It is nonetheless empty at the 5th-percentile \(T=\infty\)
reading, so its protection is a property of the informative class only
--- a trade-off, not a strict improvement over BAU, as set out in Result
3.1. The graded family is protective on the informative class but not
supply-superior. The two families are compared against the frozen rules
in Table 1.

\emph{Reason.} For a surplus-proportional rule \(C = \phi\, g(S)\) the
closed loop is \(F(S) = S + (1-\phi)\,g(S) + e\), a single concave
quadratic, and Proposition 2.3 gives three regimes rather than two. The
safe set is held from itself iff the scaled surplus at the reference
point covers the floor, \((1-\phi)\,g(K^*) \ge |e_{q10}|\), i.e.
\(\phi \le 1 - 80.87/172.46 = 0.531\); at \(\phi = 0.25\) (\(0.50\))
the margin is \(48.5\) (\(5.4\)) kt, so both qualify, the second only
just. Nonemptiness is the weaker condition
\((1-\phi)\,g_{\max} > |e_{q10}|\), i.e.
\(\phi < 1 - 80.87/296.09 = 0.727\), and it fails at \(\phi = 0.75\)
(margin \(-37.8\) kt), so that member's \(T=\infty\) kernel is empty
(\(920.2\) kt at \(T=1\)). Between the two thresholds the rule is
viable but only from above the reference point: at the added member
\(\phi = 0.60\) the informative kernel is \([1074.8, 10^4]\) kt, not
the safe set, and it is this regime that an earlier version of this paper
missed by quoting \(0.727\) as the threshold for holding the safe set.
The three tabulated verdicts were unaffected, which is why the error
survived every numeric check. Under the 5th-percentile floor the
criterion is \(\phi < 0.029\), so no surplus-proportional rule with
positive \(\phi\) holds that class --- which is exactly why the same
family is empty there, and why it does not dominate BAU under the frozen
rule. Under the perpetual-worst floor the
loop declines for every \(\phi\), so the kernel is empty. The graded
members (Family B) are piecewise-constant and are scored exactly as the
frozen rules: under the 10th-percentile class both hold the LRP from
itself at every horizon, and both are empty at the 5th-percentile
\(T=\infty\) reading. The rejected-designs note following Definition 2.4
reports the same clause-H1 emptiness for the surplus-harvest and ramp
designs, which moreover do not improve supply. Supply replays (training
mean allowed catch over the observed states) give BAU \(5\) kt, S1
\(10.0\) kt, the cascade \(16.3\) kt, flat-60 \(60\) kt, flat-0 \(0\)
kt, graded2 \(10.0\) kt, graded3 \(5.0\) kt, and the
surplus-proportional family \(7.4\)/14.7/\(22.1\) kt at
\(\phi = 0.25/0.50/0.75\). The critical-zone rule's low mean is a
property of the window, not of the rule: the cut is active in 20 of the
24 replay years (the stock sat below the reference point in all but the
1985 and 1987--1989 states, \(83\%\) of the window), a fact about the
collapsed-era estimation window rather than about the rule's
post-recovery supply properties, and the two regimes are not mixed. \ensuremath{\square}

\begin{figure}[htbp]
\centering
\includegraphics[width=\linewidth]{figs_e2_v3/fig3_reactive_rules.png}
\caption{The two reactive families: (left) the catch schedules of the surplus-proportional Family A (\(C = \phi\,g(S)\), \(\phi = 0.25, 0.50, 0.75\)) against the flat 60-kt cap, showing that the rule harvests less than the cap near the LRP while scaling with the surplus; (right) the catch schedules of the graded Family B (graded2 and graded3) against the flat 60-kt cap. Both families vanish at or below the LRP. The dominated surplus-harvest (\(\alpha\cdot\max(0,g(S)-g(K^{\ast}))\)) and reactive-ramp designs are recorded in the rejected-designs note following Definition 2.4 and not plotted.}
\end{figure}

\subsubsection{3.4 Certified layer: the expansion
obstruction}\label{budget-certified-layer-the-expansion-obstruction}

The conversion of Definition 2.7 needs the closed loop's contraction
rate. The next result shows that no such contraction is available at the
declared safe set, and that the certified layer therefore has a finite
horizon: six years under the two harsher floors, seven under the
informative one.

\textbf{Result 3.5 (Expansion obstruction).} On the governed
surplus-production map of Definition 2.1 the certified kernel is empty
from \(T = 7\) under the perpetual-worst and 5th-percentile floors,
and from \(T = 8\) under the 10th-percentile floor, for every declared
policy, zero catch included. At the last nonempty horizon the certified
sets lie far above the observed range of the stock:
\([5132.9, 10^4]\) kt (perpetual worst, BAU, \(T = 6\)),
\([4593.2, 10^4]\) kt (5th percentile, BAU, \(T = 6\)), and
\([4560.3, 10^4]\) kt (10th percentile, BAU and zero catch,
\(T = 7\)); under the informative floor the 180- and 240-kt caps are
already empty at \(T = 7\).

\emph{Reason.} Here \(F'(S) = 1 + r(1-2S/K)\) is increasing as \(S\)
falls. At the LRP, \(F'(K^*) = 1.153 > 1\). The governed surplus map is
therefore expansive at the declared safe set, contracting only above
\(K/2 = 2500\) kt --- a stock level the series never approaches. The
contraction form of the conversion is therefore inapplicable. The
expansive form \(r_T = \varepsilon(a_{\max}^T - 1)/(a_{\max} - 1)\)
grows without bound; with the committed \(\varepsilon = 328.9725\) kt --- the
perpetual-worst class of Section 2.3, printed there to one decimal as
\(329.0\) --- and \(a_{\max} = 1.1531\): \(r_1 = 329.0\), \(r_2 = 708.3\),
\(r_3 = 1145.7\), \(r_4 = 1650.1\), \(r_5 = 2231.7\),
\(r_6 = 2902.4\), \(r_7 = 3675.7\), \(r_8 = 4567.4\) kt, so the
shifted threshold \(K^* + r_T\) runs \(1213.6\), \(1592.9\),
\(2030.3\), \(2534.7\), \(3116.3\), \(3787.0\), \(4560.3\),
\(5452.0\) kt at \(T = 1,\ldots,8\).

Emptiness is not simply the shifted threshold outgrowing the carrying
capacity: at \(T = 7\) the threshold (\(4560.3\) kt) is still below
\(K = 5000\) kt, and the certified kernel is empty under the two
harsher floors nevertheless. What fails is the trajectory, not the
arithmetic of the shift. Under the perpetual-worst floor the BAU closed
loop has no worst-case fixed point at all, so every trajectory declines;
started from the domain ceiling of \(10^4\) kt it reaches \(4689.6\)
kt after five steps and \(4424.5\) kt after six, which clears the
six-year threshold (\(3787.0\) kt) but not the seven-year one
(\(4560.3\) kt). Under the informative 10th-percentile floor the loop
does have an attracting worst-case fixed point (\(4606.5\) kt for BAU,
above the seven-year threshold and below the eight-year one), which is
why the horizon is one year longer there. The certified sets at those
horizons lie above the entire observed range of the stock. \ensuremath{\square}

Proposition 2.4 names the mechanism precisely: the certified horizon is
the crossing at which the geometrically growing margin
\(K^* + r_T\) overtakes the worst-case trajectory from the domain
ceiling, and it reproduces the computed horizons exactly
(\(T^* = 6\), 6 and 7). On this object the binding obstruction to
certified intervention claims is the expansion rate itself, not the
defect magnitude. This is a
failure mode qualitatively different from the companion groundwater
object, where the governed map is contracting and the certified horizon
is defect-bound to \(T \le 3\) years. The source-year convention
lengthens the horizon by one year at every class (\(T = 5 \to 6\)
under the two harsher floors, \(T = 6 \to 7\) under the informative
floor), because it replaces the registered defect magnitude
(\(\varepsilon = 460.0\) kt) by the source-year one (\(329.0\) kt)
while leaving \(a_{\max}\) unchanged.

The horizon is not a property of the policy, and that is the second
thing the crossing shows. Table 2 gives \(T^*\) against the constant catch.
It is 6 years under the perpetual-worst and 5th-percentile floors and 7 under
the informative one for every catch from zero to \(150\) kt --- the
moratorium included --- and for every declared reactive and graded rule as
well; the crossing agrees with the committed kernel computation at all
\(48\) (rule, class) pairs the campaign tries --- \(30\) constant
catches and \(18\) declared rules. Only catches outside the constructive bound
shorten it: \(6\) years under the informative floor from \(180\) kt and
\(5\) under the perpetual-worst floor from \(200\) kt. No declared policy
lengthens it. The reason is the one the formula exposes: the margin's
year-on-year increments are \(379\), \(437\), \(504\), \(582\),
\(671\) and \(773\) kt, so every one of them exceeds the entire
admissible catch range \([0, 91.59]\) kt --- the smallest by a factor
of four --- and no reduction inside that range can add a certified year.
Certified years
cannot be bought with conservation; the horizon is set by the expansion rate
at the reference point and by nothing the manager chooses.

\textbf{Table 2.} The certified horizon \(T^*\) by constant catch, under
each declared floor class.

\begin{longtable}[]{@{}
  >{\raggedright\arraybackslash}p{(\columnwidth - 6\tabcolsep) * \real{0.2500}}
  >{\raggedleft\arraybackslash}p{(\columnwidth - 6\tabcolsep) * \real{0.2500}}
  >{\raggedleft\arraybackslash}p{(\columnwidth - 6\tabcolsep) * \real{0.2500}}
  >{\raggedleft\arraybackslash}p{(\columnwidth - 6\tabcolsep) * \real{0.2500}}@{}}
\toprule\noalign{}
\begin{minipage}[b]{\linewidth}\raggedright
Constant catch (kt)
\end{minipage} & \begin{minipage}[b]{\linewidth}\raggedleft
perpetual-worst
\end{minipage} & \begin{minipage}[b]{\linewidth}\raggedleft
5th pct.
\end{minipage} & \begin{minipage}[b]{\linewidth}\raggedleft
10th pct.
\end{minipage} \\
\midrule\noalign{}
\endhead
\bottomrule\noalign{}
\endlastfoot
0 (moratorium) & 6 & 6 & 7 \\
5 (BAU) & 6 & 6 & 7 \\
60 (flat cap) & 6 & 6 & 7 \\
91.59 (\(C^*\)) & 6 & 6 & 7 \\
120 & 6 & 6 & 7 \\
150 & 6 & 6 & 7 \\
180 & 6 & 6 & 6 \\
200 & 5 & 6 & 6 \\
215.2 (\(C_{\mathrm{vac}}\)) & 5 & 6 & 6 \\
\end{longtable}

\begin{figure}[htbp]
\centering
\includegraphics[width=\linewidth]{figs_e2_v3/fig10_cadence.png}
\caption{The certified horizon is not a property of the policy. (a) The exact
crossing of Proposition 2.4 against the constant catch, under each declared
floor class, at 1-kt resolution: the horizon is \(6\) years under the
perpetual-worst and 5th-percentile floors and \(7\) under the informative one
for every catch in the shaded admissible range, a moratorium included, and it
shortens only outside the constructive bound (\(C_{\mathrm{vac}} = 215.2\)
kt). Every declared reactive and graded rule lies on the same plateau. Grey
band: the admissible range \([0, C^*]\). (b) Why: the erosion margin grows by
\(379\)--\(773\) kt a year, so a single year's growth exceeds the whole
admissible catch range (\(C^* = 91.59\) kt, dashed) by a factor of four at its
smallest. Bars are the growth from year \(T-1\) to year \(T\), computed from
the same closed form as (a); the horizon curve reproduces all \(30\) archived
(catch, floor) pairs exactly.}
\end{figure}

\begin{figure}[htbp]
\centering
\includegraphics[width=\linewidth]{figs_e2_v3/fig4_fprime.png}
\caption{The closed loop's slope \(F'(S)\); the map is expansive at the LRP and contracts only above \(K/2 = 2500\) kt. The expansion classification is not an artifact of the particular pinned value: the grid of Section 3.7 shows \(F' \ge 1.000\) at every admissible \(K \ge 2K^* = 1769.2\) kt. The classification is conditional on \(K \ge 2K^*\): contraction (\(F' = 0.61\)--\(0.93\)) is restored only below \(2K^*\) --- exactly where the informative certificates of Sections 3.2--3.3 collapse and the fit cost rises --- and the residual bootstrap of Section 3.10 gives \(F'\) a \(90\%\) interval of \([1.010, 1.179]\).}
\end{figure}

\subsubsection{3.5 Stress replay and
classification}\label{budget-stress-replay-and-classification}

The stress replay runs the closed loop from the observed 1990 SSB
(\(861.9\) kt --- already below the LRP, so the replay starts outside
the safe set and is uncontrolled shock accounting rather than a
kernel-membership test) with the observed 1991--1995 source-year
residuals. In 1991 every policy holds the stock above the LRP: flat-0
and S1 reach \(953.9\) kt (the S1 switch rule observes the sub-LRP 1990
stock and prescribes zero catch, so it coincides with flat-0 in that
year), business-as-usual \(948.9\) kt, the cascade \(923.9\) kt, and
flat-60 \(893.9\) kt. Yet by 1992 every policy has fallen below the LRP
--- flat-60 to \(678.9\) kt, S1 to \(747.8\) kt, the cascade to
\(713.4\) kt, and BAU to \(797.1\) kt --- because the 1992 transition
carries the perpetual-worst observed shock (\(-329.0\) kt). The sharpest
drawdown is under flat-60 (to \(434.3\) kt in 1993), while the switch
rules lose less (BAU \(669.7\) kt in 1994); no policy avoids the
post-1992 decline. The crash is a productivity event, not a catch event
--- exactly the catch-insufficiency certificate of the companion
forecast-evaluation study.

At the \(T=5\) classification under the 10th-percentile class, only the
1980s peak years of the 33 observed states lie inside the nominal
kernels --- \{1985, 1987, 1989\} for the 60-kt rules and additionally
1988 for BAU. The entire post-1990 history and most of the 1980s are
outside. Under the perpetual-worst class all 33 are outside; under the
5th-percentile class the moratorium's kernel is nonempty (its
\(T=\infty\) boundary is \(2219.6\) kt, above the observed range), so
the classification there follows the informative geometry rather than
universal emptiness.

\begin{figure}[htbp]
\centering
\includegraphics[width=\linewidth]{figs_e2_v3/fig5_replay.png}
\caption{Closed-loop replay from the observed 1990 stock with the observed residuals; BAU, flat-0, flat-60, the S1 switch rule, and the cascade are plotted separately (the S1 rule differs from the flat-60 cap below the LRP).}
\end{figure}

\subsubsection{3.6 Model-form comparison: the depensatory refit and the
Fox form as co-equal
specifications}\label{budget-model-form-comparison-the-depensatory-refit-and-the-fox-form-as-co-equal-specifications}

The primary kernels are Schaefer-form (Allee term off). The depensation
sensitivity refits the same object with the Allee term on, on the same
1983--2007 window with the same annual landings. The refit gives
\(r = 2.0\) (pinned at its optimization bound, the mirror image of the
registered \(K\)), \(K = 1671.7\) kt, and \(s_0 = 642.3\) kt --- \(242\)
kt below the LRP, inside the observed range --- with residual MSE
\(7690.1\) kt\(^2\) against \(12{,}772.2\) kt\(^2\) for the registered
Schaefer form. The bare fit parameters and the MSE are
convention-independent (the fit already pairs the year-\(t\) catch with
the \(t \to t+1\) transition), so they are unchanged from the frozen
reading. Only the fitted form's own floor classes and its kernels depend
on the residual convention, and those are reported here in the
source-year convention. The depensatory form is the better-fitting of
the two, under identification caveats of the same kind as the registered
fit.

The two forms are reported side by side as co-equal primary
specifications in Table 3. The identification caveats apply
symmetrically: the Schaefer form pins \(K\) at its bound, the
depensatory form pins \(r\) at its bound (\(2.0\)), and neither pin is a
biological finding. Table 3 reports the rows at the declared class
endpoints, held frozen across forms as the corrected source-year classes
(\(-287.4\), \(-80.9\), \(-329.0\) kt) so that only the form varies.

\textbf{Result 3.6 (Form sensitivity).} The certificate directions
survive the data-preferred depensatory refit. The Allee form removes the
vacuity of the perpetual-worst floor --- the one class vacuous on the
registered form --- because its maximum surplus (\(372.4\) kt) exceeds
both class magnitudes (\(-329.0\) and \(-287.4\) kt yr\textsuperscript{-1}).

\emph{Reason.} Under the 10th-percentile class the BAU and zero-catch
certificates are exactly unchanged --- lower boundary \(884.6\) kt at
every horizon, \(T=\infty\) included --- and the critical-zone and
cascade rules are protective of the entire safe set at that class
(\(884.6\) kt). The constructive boundary of Section 3.3 and the
dominance verdict of Section 3.2 are therefore unaffected --- indeed the
constructive bound rises, since at the declared 10th-percentile floor the
Allee form gives \(g(K^*) - |e_{q10}| = 196.06 - 80.87 = 115.2\) kt
(\(123.3\) kt on the declared-strength row), so the \(91.59\) kt
certificate holds a fortiori. The clause-(H1) block itself changes shape without changing the verdict:
under the Allee form the 60-kt rules are no longer empty at the
5th-percentile \(T=\infty\) reading (their boundary is \(1131.1\) kt
against BAU's \(1020.9\) kt), so the block operates through a margin
rather than through emptiness at that reading. Under the
Allee form, whose maximum surplus is \(372.4\) kt, both the
perpetual-worst floor (\(-329.0\) kt) and the 5th-percentile floor
(\(-287.4\) kt) lie below it, so the BAU kernel is nonempty at those
classes with infinite-horizon boundaries \(1098.7\) and \(1020.9\) kt
respectively (against the empty kernels of the registered Schaefer form
under the perpetual-worst floor). The expansion obstruction is not a
Schaefer-form artifact: the Allee refit raises the expansion rate at the
LRP from \(1.1531\) to \(1.7818\). A declared-strength alternative
(\(s_0 = 0.5K^* = 442.3\) kt) refits to \(r = 2.0\), \(K = 3223.7\) kt
(MSE \(9330.0\) kt\(^2\)), with maximum surplus \(883.6\) kt; under the
source-year classes its kernels are nonempty at every reported cell and
lie within \(12\) kt (about \(1\%\)) of the identified row's at all of
them, so the depensation row does not hinge on the identification of
\(s_0\).

\emph{Reproducibility note (two declared sensitivities).} First, the
manuscript's \(s_0 = 642.3\) kt is a rounding of the fitted
\(642.3296\) kt, and one printed cell moves by a single grid step between
them: the BAU worst-class \(T=\infty\) boundary is \(1098.75\) kt at the
rounded \(s_0\) and \(1098.8\) kt at the fitted one (the q05
\(T=\infty\) boundary moves \(1020.9 \to 1020.95\) kt, which prints
identically). Every other cell is unchanged, and the published cells are
those of the rounded \(s_0\). Second, the Allee and Fox rows are computed
on a \(0.05\) kt state grid, while the registered row comes from the
committed interval-arithmetic engine; the campaign reproduces the
registered row on that same grid, so the two engines can be read against
each other. At finite horizons they agree to within about one grid step
--- of the \(81\) cells the campaign checks, the \(72\) finite-horizon
ones differ by at most \(0.06\) kt --- but adjacent to the repelling
boundary, where \(F' > 1\), the grid's rounding compounds along the orbit:
at \(T = \infty\) five of the nine cells agree exactly and the other
four differ by \(0.1\), \(0.3\), \(0.6\) and \(0.9\) kt, the largest
being the BAU q05 boundary, \(2218.75\) against \(2219.649\) kt (and
\(2070.30\) against \(2070.884\) kt for zero catch). The two engines are
therefore not to be compared at \(0.1\) kt resolution, and the tolerance
the campaign declares is \(1.0\) kt.

The third co-equal form is the Fox surplus law \(g(S) = rS\ln(K/S)\),
refitted one-step on the same window with the same landings and the same
declared box. The refit gives \(r = 0.1044\), \(K = 5000\) kt pinned at
its bound, residual MSE \(13{,}873.1\) kt\(^2\) (Schaefer
\(12{,}772.2\); Allee \(7690.1\)). The Fox form is the worst-fitting of
the three, and the Allee form remains the data-preferred one. The
certificate directions survive the form. The expansion at the LRP is
preserved (\(F' = 1.0764 > 1\), weaker than the registered \(1.1531\)).
The constructive 10th-percentile bound is
\(g(K^*) - |e_{q10}| = 159.92 - 80.87 = 79.05\) kt, under the
source-year 10th-percentile floor (\(-80.87\) kt); the Fox form's own
refit-boundary floor is not used so that only the form varies across
Table 3. The BAU and zero-catch certificates under the informative class
are exactly unchanged (lower boundary \(884.6\) kt at every horizon,
\(T=\infty\) included). Three further Fox cells are computed by the campaign but not carried in
Table 3: the 60-kt rule's infinite-horizon lower boundary is empty under
the 5th-percentile floor, and the 120-kt rule is empty under the
5th-percentile and perpetual-worst floors and, under the informative
floor, at \(T=\infty\) (its \(T=1\) boundary there is \(922.7\) kt). Table 3 carries the Fox row.

\textbf{Table 3.} Depensation and form sensitivity: kernel lower
boundaries (kt) at the declared source-year class endpoints.

\begin{longtable}[]{@{}
  >{\raggedright\arraybackslash}p{(\columnwidth - 10\tabcolsep) * \real{0.1304}}
  >{\raggedleft\arraybackslash}p{(\columnwidth - 10\tabcolsep) * \real{0.1739}}
  >{\raggedleft\arraybackslash}p{(\columnwidth - 10\tabcolsep) * \real{0.1739}}
  >{\raggedleft\arraybackslash}p{(\columnwidth - 10\tabcolsep) * \real{0.1739}}
  >{\raggedleft\arraybackslash}p{(\columnwidth - 10\tabcolsep) * \real{0.1739}}
  >{\raggedleft\arraybackslash}p{(\columnwidth - 10\tabcolsep) * \real{0.1739}}@{}}
\toprule\noalign{}
\begin{minipage}[b]{\linewidth}\raggedright
Form
\end{minipage} & \begin{minipage}[b]{\linewidth}\raggedleft
BAU, q05 T=1
\end{minipage} & \begin{minipage}[b]{\linewidth}\raggedleft
BAU, q05 T=\(\infty\)
\end{minipage} & \begin{minipage}[b]{\linewidth}\raggedleft
BAU, q10 T=\(\infty\)
\end{minipage} & \begin{minipage}[b]{\linewidth}\raggedleft
BAU, worst T=\(\infty\)
\end{minipage} & \begin{minipage}[b]{\linewidth}\raggedleft
S1, q10 T=\(\infty\)
\end{minipage} \\
\midrule\noalign{}
\endhead
\bottomrule\noalign{}
\endlastfoot
Registered Schaefer & 989.0 & 2219.6 & \textbf{884.6} & empty &
\textbf{884.6} \\
Allee refit (\(s_0\) = 642.3) & 939.4 & 1020.9 & \textbf{884.6} & 1098.7
& \textbf{884.6} \\
Declared \(s_0\) = \(0.5K^{\ast}\) (442.3), refit & 939.2 & 1024.1 &
\textbf{884.6} & 1086.9 & \textbf{884.6} \\
Fox refit (\(r = 0.1044\), \(K\) pinned) & 1008.4 & empty &
\textbf{884.6} & empty & \textbf{884.6} \\
\end{longtable}

\subsubsection{3.7 Carrying-capacity
sensitivity}\label{budget-carrying-capacity-sensitivity}

The registered fit pins \(K\) at its optimization bound, and the
expansion classification of Section 3.4 inherits the question of how
much of the result is carried by that pin. The sensitivity refits \(r\)
by one-step least squares at each \(K\) on the same window, in the
source-year convention, with the kernel scored against each refit's own
source-year 10th-percentile floor.

\textbf{Table 4.} Carrying-capacity sensitivity in the source-year
convention: \(r\) refit at fixed \(K\), kernel scored against that
refit's own corrected 10th-percentile floor. Values are the BAU lower
boundary of the robust viability kernel at \(T=1\) and \(T=\infty\).

\begin{longtable}[]{@{}
  >{\raggedright\arraybackslash}p{(\columnwidth - 14\tabcolsep) * \real{0.1000}}
  >{\raggedleft\arraybackslash}p{(\columnwidth - 14\tabcolsep) * \real{0.1333}}
  >{\raggedleft\arraybackslash}p{(\columnwidth - 14\tabcolsep) * \real{0.1333}}
  >{\raggedleft\arraybackslash}p{(\columnwidth - 14\tabcolsep) * \real{0.1333}}
  >{\raggedleft\arraybackslash}p{(\columnwidth - 14\tabcolsep) * \real{0.1333}}
  >{\raggedleft\arraybackslash}p{(\columnwidth - 14\tabcolsep) * \real{0.1333}}
  >{\raggedleft\arraybackslash}p{(\columnwidth - 14\tabcolsep) * \real{0.1333}}
  >{\raggedright\arraybackslash}p{(\columnwidth - 14\tabcolsep) * \real{0.1000}}@{}}
\toprule\noalign{}
\begin{minipage}[b]{\linewidth}\raggedright
\(K\) (kt)
\end{minipage} & \begin{minipage}[b]{\linewidth}\raggedleft
\(r\)
\end{minipage} & \begin{minipage}[b]{\linewidth}\raggedleft
\(g_{\max}\)
\end{minipage} & \begin{minipage}[b]{\linewidth}\raggedleft
\(F'(K^*)\)
\end{minipage} & \begin{minipage}[b]{\linewidth}\raggedleft
Constr.\\bound (kt)
\end{minipage} & \begin{minipage}[b]{\linewidth}\raggedleft
BAU q10, \(T=1\)
\end{minipage} & \begin{minipage}[b]{\linewidth}\raggedleft
BAU q10, \(T=\infty\)
\end{minipage} & \begin{minipage}[b]{\linewidth}\raggedright
q05 vacuous
\end{minipage} \\
\midrule\noalign{}
\endhead
\bottomrule\noalign{}
\endlastfoot
1000 & 0.5094 & 127.4 & 0.608 & \(-28.87\) & 943.2 & empty & yes \\
1200 & 0.5248 & 157.4 & 0.751 & 41.14 & 884.6 & 884.6 & yes \\
1500 & 0.4099 & 153.7 & 0.926 & 67.9 & 884.6 & 884.6 & yes \\
1769.2 (\(=2K^*\)) & 0.3559 & 157.4 & 1.000 & 76.55 & 884.6 & 884.6 &
yes \\
2000 & 0.3273 & 163.6 & 1.038 & 80.6 & 884.6 & 884.6 & yes \\
2500 & 0.2908 & 181.7 & 1.085 & 85.33 & 884.6 & 884.6 & yes \\
3000 & 0.2704 & 202.8 & 1.111 & 87.79 & 884.6 & 884.6 & yes \\
4000 & 0.2485 & 248.5 & 1.139 & 90.31 & 884.6 & 884.6 & yes \\
5000 (registered) & 0.2369 & 296.1 & 1.153 & \textbf{91.59} & 884.6 &
884.6 & \textbf{no} \\
7000 (out of box) & 0.2248 & 393.5 & 1.168 & 92.89 & 884.6 & 884.6 &
no \\
\end{longtable}

\textbf{Result 3.7 (Carrying-capacity grid: three readings).} The grid
of Table 4 supports three readings.

\begin{enumerate}
\def\labelenumi{(\roman{enumi})}
\item
  The expansion obstruction holds for every admissible
  \(K \ge 2K^* = 1769.2\) kt and is the data-selected regime within that
  range --- not an artifact of the particular pinned value \(K = 5000\)
  --- while remaining conditional on \(K \ge 2K^*\): \(F' \ge 1.000\) at
  every \(K \ge 2K^*\) --- exactly \(1.0000\) at \(K = 2K^*\), rising to
  \(1.1531\) at the registered \(K\) --- and the informative
  constructive bound rises monotonically toward the registered end of
  the box (\(-28.87\) kt at \(K = 1000\) kt to \(91.59\) kt at
  \(K = 5000\) kt), with the fit cost rising as \(K\) falls below
  \(2K^*\).
\item
  Contraction is restored only below \(2K^*\) (\(F' = 0.61\)--\(0.93\)),
  exactly where the informative certificates collapse: the constructive
  bound is negative at \(K = 1000\) kt, and the BAU kernel under the
  10th-percentile class is empty at \(T = \infty\) with the \(T = 1\)
  boundary raised to \(943.2\) kt --- the moratorium itself no longer
  holds the reference point.
\item
  The vacuity structure is \(K\)-dependent in the source-year
  convention. The perpetual-worst floor is vacuous at every in-box
  \(K\), but the 5th-percentile class is vacuous only at \(K \le 4000\)
  kt and becomes informative at \(K \ge 5000\) kt (\(g_{\max} = 296.1\)
  kt at the registered \(K\)), which is itself the registered point.
\end{enumerate}

\begin{figure}[htbp]
\centering
\includegraphics[width=\linewidth]{figs_e2_v3/fig6_k_sensitivity.png}
\caption{The two panels of the carrying-capacity sensitivity of Section 3.7.}
\end{figure}

How much of that is carried by the pin? The question can be answered without
resolving it. Profiling \(K\) --- refitting \(r\) in closed form at each
fixed \(K\) on the same window, with the grid extended to \(50{,}000\)
kt precisely because the committed value sits on the box --- gives a
profile-likelihood \(95\%\) set running from \(1500\) kt to the top of the
grid: the series never approaches carrying capacity, so \(K\) is not
identified from above at all. The pin itself buys little: the sum of squares at
\(K = 5000\) kt is \(1.027\) times the profiled minimum. What is identified
is everything the results depend on. Across the profile set \(g(K^*)\) lies
in \([148.8, 176.1]\) kt and \(C^*\) in \([67.9, 95.2]\) kt, so the
constructive bound is pinned to within \(15\%\) by data that do not determine
\(K\) at all; and \(F'(K^*) > 1\) for every profile-set \(K\) at or above
\(2K^* = 1769.2\) kt --- the smallest expansive value on the grid is
\(1775\) kt --- which is exactly the condition reading (i) declares. The
expansion obstruction is conditional on \(K \ge 2K^*\); it is not conditional
on the pinned value.


\begin{figure}[htbp]
\centering
\includegraphics[width=\linewidth]{figs_e2_v3/fig9_identification.png}
\caption{Identification under the pinned carrying capacity. (a) The profile
sum of squares over \(K\), refitting \(r\) in closed form at each \(K\): the
profile is flat from about \(1500\) kt to the top of the grid, so \(K\) is not
identified from above and the committed value (dotted) buys \(2.7\%\) of the
criterion. (b) The functionals the results depend on are identified
nevertheless: across the profile \(95\%\) set (shaded), \(g(K^*)\) and
\(C^* = g(K^*) - |e_{q10}|\) vary by \(15\%\), and \(F'(K^*)\) (right axis)
exceeds 1 exactly above \(2K^*\).}
\end{figure}

\subsubsection{3.8 Stochastic viability}\label{budget-stochastic-viability}

Persistent floors are a worst-case layer. The empirical residual pool
supplies the stochastic counterpart, with the 1992 draw treated both as
a member of the pool and as a removed one-off. For each policy and
initial stock, \(20{,}000\) trajectories resample the 24 source-year
training residuals (i.i.d., and in moving blocks of four, respecting the
residual autocorrelation of \(0.55\)). Viability is scored as the
probability of remaining at or above the LRP.

\textbf{Table 5.} Probability of remaining at or above the LRP over 20
years (\(N = 20{,}000\) per cell; seed fixed; draws shared across
policies), source-year pool.

\begin{longtable}[]{@{}
  >{\raggedright\arraybackslash}p{(\columnwidth - 8\tabcolsep) * \real{0.1579}}
  >{\raggedleft\arraybackslash}p{(\columnwidth - 8\tabcolsep) * \real{0.2105}}
  >{\raggedleft\arraybackslash}p{(\columnwidth - 8\tabcolsep) * \real{0.2105}}
  >{\raggedleft\arraybackslash}p{(\columnwidth - 8\tabcolsep) * \real{0.2105}}
  >{\raggedleft\arraybackslash}p{(\columnwidth - 8\tabcolsep) * \real{0.2105}}@{}}
\toprule\noalign{}
\begin{minipage}[b]{\linewidth}\raggedright
Policy
\end{minipage} & \begin{minipage}[b]{\linewidth}\raggedleft
from the LRP, i.i.d.
\end{minipage} & \begin{minipage}[b]{\linewidth}\raggedleft
from the LRP, blocks-4
\end{minipage} & \begin{minipage}[b]{\linewidth}\raggedleft
from the LRP, i.i.d. without 1992
\end{minipage} & \begin{minipage}[b]{\linewidth}\raggedleft
from 1500 kt, i.i.d.
\end{minipage} \\
\midrule\noalign{}
\endhead
\bottomrule\noalign{}
\endlastfoot
zero catch & 0.906 & 0.852 & 0.954 & 1.000 \\
BAU (5 kt) & 0.903 & 0.852 & 0.954 & 1.000 \\
60 kt / S1 / cascade & 0.835 & 0.806 & 0.917 & 1.000 \\
120 kt & 0.647 & 0.650 & 0.769 & 0.999 \\
\end{longtable}

From 1980s-peak starts (\(1500\) kt and above) survival is near certain
under every scheme --- the peaks sit far above the LRP. The corrected
source-year pool is marginally less severe than the frozen
destination-year pool, so the survival probabilities are higher across
the board.

The constructive boundary of Section 3.3 carries a stochastic reading.
At the grid catch nearest the bound (\(92.5\) kt on the campaign's
\(2.5\) kt grid) the 20-year survival probability from the LRP is
\(0.74\) under i.i.d. draws, \(0.73\) under block resampling and
\(0.84\) when the 1992 draw is removed --- roughly three-quarters, not
the near-certainty that holds at moratorium-level removals.

The \(P \ge 0.9\) bar separates the two resampling schemes. Under i.i.d.
resampling it is attained, but only by near-moratorium catches: the
interpolated crossing is \(13.5\) kt against a ceiling of \(0.906\) at
zero catch (Table 5), and by \(60\) kt survival has already fallen to
\(0.835\). Under block resampling it is not attained at all, the ceiling
there being \(0.852\) (Table 5; \(0.849\) in the crossing sweep, which
carries its own fixed seed --- a Monte-Carlo difference of \(0.003\) at
\(N = 20{,}000\)); blocks of four preserve the residual autocorrelation
of Section 2.1, so a poor year tends to be followed by another. Removing
the 1992 draw lifts that crossing to \(78.9\) kt. The \(P \ge 0.8\)
crossings are \(81.2\) kt (i.i.d.), \(72.3\) kt (blocks) and \(105.2\)
kt (no-1992). The one-off treatment of 1992 is a declared sensitivity,
not an identified break.

Two properties of the pool explain the ceiling. First, failures are
concentrated at the start of the horizon: the chance of an immediate
breach is \(2/24 = 0.083\) against a total failure probability of
\(0.094\), and \(99\%\) of failures occur within the first three years.
The reason is that from the reference point the stock is one year's
growth from safety --- under zero catch a draw worse than \(-172.5\) kt
breaches the LRP at once, and two of the twenty-four residuals are worse
than that, the 1992 draw (\(-329.0\) kt) and one other (\(-323.5\) kt).
Second, a draw of that magnitude is harmless once the stock has grown
past about \(1020\) kt, roughly a year later, which is why the ceiling
is \(0.906\) and not the \(0.43\) that the \(1/24\) recurrence alone
would suggest: the probability of escaping the 1992 draw for twenty
years is \((23/24)^{20} = 0.427\), and that would be the ceiling only if
every 1992 draw were fatal.

The constructive bound of \(91.59\) kt sits above the 60-kt flat-cap
family, so those rules are robust in the worst case, while the stochastic
layer shows the survival probability falling steadily across that range
(\(0.91\) at zero-to-moratorium removals to \(0.65\) at \(120\) kt).

\begin{figure}[htbp]
\centering
\includegraphics[width=\linewidth]{figs_e2_v3/fig7_stochastic.png}
\caption{The three trade-off curves: 20-year survival probability from the LRP against constant catch, by resampling scheme (i.i.d. draws, blocks, no-1992).}
\end{figure}

\subsubsection{3.9 Finite-duration floors}\label{budget-finite-duration-floors}

A perpetual floor is a new stationary climate. The finite-duration
variant asks how long a poor-productivity episode must last before the
safe set contracts, with the floor active for \(n\) years and zero
residual thereafter. The boundaries are computed by exact backward
recursion on the monotone map, in the source-year convention (the
\(n=5\) worst-floor BAU boundary reproduces the registered \(T=5\)
boundary of the archived kernel table as a built-in check).

\textbf{Table 6.} Infinite-horizon lower boundary (kt) after a finite
floor of \(n\) years followed by zero residual (source-year); empty = no
state is robustly viable.

\begin{longtable}[]{@{}
  >{\raggedright\arraybackslash}p{(\columnwidth - 12\tabcolsep) * \real{0.1111}}
  >{\raggedleft\arraybackslash}p{(\columnwidth - 12\tabcolsep) * \real{0.1481}}
  >{\raggedleft\arraybackslash}p{(\columnwidth - 12\tabcolsep) * \real{0.1481}}
  >{\raggedleft\arraybackslash}p{(\columnwidth - 12\tabcolsep) * \real{0.1481}}
  >{\raggedleft\arraybackslash}p{(\columnwidth - 12\tabcolsep) * \real{0.1481}}
  >{\raggedleft\arraybackslash}p{(\columnwidth - 12\tabcolsep) * \real{0.1481}}
  >{\raggedleft\arraybackslash}p{(\columnwidth - 12\tabcolsep) * \real{0.1481}}@{}}
\toprule\noalign{}
\begin{minipage}[b]{\linewidth}\raggedright
Policy
\end{minipage} & \begin{minipage}[b]{\linewidth}\raggedleft
q05, \(n=5\)
\end{minipage} & \begin{minipage}[b]{\linewidth}\raggedleft
q05, \(n=10\)
\end{minipage} & \begin{minipage}[b]{\linewidth}\raggedleft
q05, \(n=15\)
\end{minipage} & \begin{minipage}[b]{\linewidth}\raggedleft
worst, \(n=5\)
\end{minipage} & \begin{minipage}[b]{\linewidth}\raggedleft
worst, \(n=10\)
\end{minipage} & \begin{minipage}[b]{\linewidth}\raggedleft
worst, \(n=15\)
\end{minipage} \\
\midrule\noalign{}
\endhead
\bottomrule\noalign{}
\endlastfoot
zero catch & 1280.8 & 1510.5 & 1657.6 & 1431.6 & 1770.9 & 2015.3 \\
BAU (5 kt) & 1298.7 & 1540.7 & 1697.8 & 1450.0 & 1803.7 & 2062.3 \\
60 kt / S1 / cascade & 1499.6 & 1893.2 & 2193.4 & 1656.7 & 2188.8 &
2657.5 \\
120 kt & 1727.7 & 2329.0 & 2898.1 & 1891.7 & 2671.8 & 3562.3 \\
180 kt & 1999.0 & 2869.7 & 3986.2 & 2172.3 & 3285.8 & 5180.7 \\
240 kt & 2623.1 & 4013.4 & 8790.9 & 2825.0 & 4687.8 & empty \\
\end{longtable}

\textbf{Result 3.8 (Finite-duration floors).} Two readings hold for the
finite-duration grid.

\begin{enumerate}
\def\labelenumi{(\roman{enumi})}
\item
  Duration is informative. A five-year 5th-percentile episode already
  lifts the BAU boundary from the LRP to \(1298.7\) kt, and a
  fifteen-year episode lifts it to \(1697.8\) kt. The safe set contracts
  as the episode lengthens, without the arithmetic vacuity of the
  perpetual classes. The source-year floors are milder than the frozen
  ones, so the boundaries sit below the earlier reading at every
  duration.
\item
  The ordering across policies is the same as in Table 1.
  Moratorium-level removals dominate every harvesting rule at every
  duration, and the reactive rules coincide with their 60-kt cap
  throughout because on the kernel domain their catch is exactly the
  cap's.
\end{enumerate}

\subsubsection{3.10 Uncertainty bands}\label{budget-uncertainty-bands}

The \(91.59\) kt constructive boundary is an analytic limit. Its sampling
distribution is reported because a point would overstate the precision
of the identified \(r\). A parametric residual bootstrap (\(B = 2000\),
seed fixed) generates synthetic one-step series on the registered map
with resampled source-year training residuals and refits \(r\) at
\(K = 5000\) kt on each. The refits give \(r\) a median of \(0.219\)
(\(90\%\) interval \([0.016, 0.277]\)), \(g(K^*)\) a median of \(159.6\)
kt yr\(^{-1}\) (\([11.6, 201.9]\)), and the constructive bound a median
of \(78.7\) kt with \(90\%\) interval \([0.0, 121.1]\) kt. Of these
refits, \(88.2\%\) retain a positive constructive bound. The expansion
classification survives the resampling: \(F'(K^*)\) has median \(1.142\)
and \(90\%\) interval \([1.010, 1.179]\). The constructive boundary is
therefore best read as order \(70\)--\(90\) kt under the declared
10th-percentile class, not as a quota figure.

Refitting \(r\) alone at \(K = 5000\) kt is the narrower of the two
available bands, because it conditions the pin away rather than propagating
it. Refitting both parameters on every replicate gives \(r\) a median of
\(0.261\) (\(90\%\) interval \([0.092, 0.500]\)), \(K\) a median of
\(4191\) kt with the pin active in \(41.7\%\) of replicates, and the
constructive bound a median of \(73.7\) kt with \(90\%\) interval
\([-89.4, 125.7]\) kt. The interval's negative tail is made of replicates in
which the refit places \(K\) below the reference point --- admissible under
the declared box, and excluded by construction when \(K\) is held fixed.
Restricted to the expansive regime \(K \ge 2K^*\) on which the results are
conditional --- \(72.9\%\) of the \(2000\) replicates --- the bound has
median \(88.1\) kt (\(90\%\) interval
\([-5.6, 130.6]\) kt) and \(F'(K^*)\) has median \(1.137\) with
\(90\%\) interval \([1.030, 1.178]\): expansion holds throughout.

Two intervals for one quantity, a factor of eight apart in width, need
the difference said out loud rather than left for a reader to reconcile.
The profile range of Section 3.7 is an \(F\)-based interval: it holds
the fit inside the \(95\%\) likelihood cut, lets \(K\) run the whole
ridge from \(1500\) to \(50{,}000\) kt with \(r\) reprofiled at each
step, and asks how far \(C^*\) can move. It answers an identification
question --- is the functional pinned, given the model --- and the answer
is yes, to within \(27\) kt, even though \(K\) is not identified at
all. The bootstrap answers a different question: how far the estimate
itself moves when the \(24\) transitions are resampled and both
parameters are refitted freely. It is wider for two separable reasons.
Its lower tail comes from replicates in which the refit places \(K\)
below the reference point (\(7.4\%\) of them) --- admissible under the
declared box, excluded by construction when \(K\) is held fixed, and
not admissible as a description of a limit reference point. Removing
those and every replicate outside the expansive regime still leaves
\([-5.6, 130.6]\) kt, so the width is not an artefact of degenerate
refits: the bootstrap is wider at both ends, \(35\) kt above the profile
ceiling and \(74\) kt below its floor. What remains is small-sample:
with \(24\) transitions the resampling dispersion of a two-parameter
refit exceeds what an \(F\)-based cut implies, and the profile range
should not be read as a statement of precision. Neither interval is in
error. They are evidence for different claims, and the claim this paper
makes is the identification one.

One further declared sensitivity is made quantitative here. The residual of
Section 2.1 conflates process noise and observation error, and only the process
part can act as a persistent floor. If a fraction \(\lambda\) of the residual
variance is observation error, rescaling the empirical pool about its mean by
\(\sqrt{1-\lambda}\) --- its shape preserved, only its scale changed ---
moves the constructive bound from \(91.6\) kt at \(\lambda = 0\) to
\(112.1\) kt at \(\lambda = 0.5\), and the vacuous class disappears from
\(\lambda \ge 0.2\) onward. The reported bound is the conservative end of
that range.

\subsubsection{3.11 Second specification (xteNCAM): a labelled
sensitivity
row}\label{budget-second-specification-xtencam-a-labelled-sensitivity-row}

The registered object is the NCAM series. The second, unpooled
specification is the xteNCAM series (Regular et al., 2025, Table 17;
1954--2024; LRP \(= 276\) kt; 2024 landings persisted from 2023),
refitted in the same source-year convention as the primary object
(Section 2.3) on 1954--2007, with the same box rule and the safe set
written against its own reference point. The floor
classes are not transferred from the NCAM fit (different series scale):
the row declares its own classes from its own training residuals, and no
verdict transfers between the specifications. The fit gives
\(r = 0.5023\), \(K = 4812.9\) kt (unpinned), MSE \(18{,}028\) kt\textsuperscript{2}, and
expansion at the LRP \(F' = 1.4447\) --- the expansion classification is
preserved and stronger than on the registered series.

The informative certificate does not survive. Under the fit's own
10th-percentile class the constructive bound is negative
(\(g(\mathrm{LRP}) = 130.7\) kt against \(|e_{q10}| = 178.7\) kt). The
zero-catch kernel does not hold the LRP from the LRP itself (\(T=1\)
lower boundary \(309.4\) kt against the \(276\)-kt reference point;
business-as-usual at \(5\) kt needs \(312.9\) kt). The infinite-horizon
zero-catch boundary under that class is \(386.9\) kt --- the reference
point is not robustly viable from itself on the second specification,
and every positive constant catch raises the \(T=1\) boundary further
(\(351.2\) kt at \(60\) kt, \(393.4\) kt at \(120\) kt). The harsh
classes are not vacuous there (\(g_{\max} = 604.4\) kt \(> 470.8\) kt),
and their \(T=1\) boundaries run \(515.7\) kt (zero catch) to \(602.3\)
kt (\(120\) kt).

The row is a labelled sensitivity with opposite structure at its own
reference point (Table 7). The two specifications agree on the expansion
classification and on the kernel ordering across the declared catches
(zero catch \ensuremath{\leq} business-as-usual \ensuremath{\leq} 60 kt \ensuremath{\leq} 120 kt at \(T=1\) on both).
They disagree on the reference point's self-viability. The 2024 xteNCAM
stock (\(342\) kt) sits between the second specification's \(T=1\)
(\(309\) kt) and \(T=5\) (\(368\) kt) 10th-percentile zero-catch
boundaries.

\textbf{Table 7.} The two specifications compared at their own reference
points (the registered NCAM row and the second, unpooled xteNCAM row of
Section 3.11).

\begin{longtable}[]{@{}
  >{\raggedright\arraybackslash}p{(\columnwidth - 12\tabcolsep) * \real{0.1111}}
  >{\raggedleft\arraybackslash}p{(\columnwidth - 12\tabcolsep) * \real{0.1481}}
  >{\raggedleft\arraybackslash}p{(\columnwidth - 12\tabcolsep) * \real{0.1481}}
  >{\raggedleft\arraybackslash}p{(\columnwidth - 12\tabcolsep) * \real{0.1481}}
  >{\raggedleft\arraybackslash}p{(\columnwidth - 12\tabcolsep) * \real{0.1481}}
  >{\raggedleft\arraybackslash}p{(\columnwidth - 12\tabcolsep) * \real{0.1481}}
  >{\raggedleft\arraybackslash}p{(\columnwidth - 12\tabcolsep) * \real{0.1481}}@{}}
\toprule\noalign{}
\begin{minipage}[b]{\linewidth}\raggedright
Object
\end{minipage} & \begin{minipage}[b]{\linewidth}\raggedleft
\(r\)
\end{minipage} & \begin{minipage}[b]{\linewidth}\raggedleft
\(K\) (kt)
\end{minipage} & \begin{minipage}[b]{\linewidth}\raggedleft
\(F'(\mathrm{LRP})\)
\end{minipage} & \begin{minipage}[b]{\linewidth}\raggedleft
Constr.\\bound (kt)
\end{minipage} & \begin{minipage}[b]{\linewidth}\raggedleft
zero-catch q10, \(T=1\)
\end{minipage} & \begin{minipage}[b]{\linewidth}\raggedleft
zero-catch q10, \(T=\infty\)
\end{minipage} \\
\midrule\noalign{}
\endhead
\bottomrule\noalign{}
\endlastfoot
NCAM (registered) & 0.2369 & 5000 (pinned) & 1.1531 & 91.59 & 884.6 &
884.6 \\
xteNCAM (this row) & 0.5023 & 4812.9 & 1.4447 & \ensuremath{-}48.0 & 309.4 & 386.9 \\
\end{longtable}


\subsubsection{3.12 Is the reference point self-viable under current
productivity?}\label{budget-is-the-reference-point-self-viable-under-current-productivity}

The fit window 1983--2007 spans the 1992 collapse, so the identified \(r\)
averages two productivity regimes and the floor classes average them with it.
Refitting the same object on post-moratorium windows --- same source-year
convention, same declared box, same construction throughout --- asks whether
the reference point is self-viable under recent productivity (Table 8). The
rows are labelled sensitivities in the sense of Section 3.11: no verdict
transfers between the two series, and none of these windows replaces the frozen
protocol.

\textbf{Table 8.} The same construction on post-moratorium windows. \(K\)
is reported with its bound status (ub: upper bound of the declared box,
lb: lower bound, i.e. just above the largest observed stock on that
window); \(C^*\) is the constructive bound at each window's own
10th-percentile floor. The same quantity at each window's perpetual-worst
floor is given in the text; the residual SD, which falls from \(114.9\) kt on
the frozen window to \(18.9\) kt on NCAM 1995--2015, is what drives the change.

\begin{longtable}[]{@{}
  >{\raggedright\arraybackslash}p{(\columnwidth - 16\tabcolsep) * \real{0.2200}}
  >{\raggedleft\arraybackslash}p{(\columnwidth - 16\tabcolsep) * \real{0.0700}}
  >{\raggedleft\arraybackslash}p{(\columnwidth - 16\tabcolsep) * \real{0.1000}}
  >{\raggedleft\arraybackslash}p{(\columnwidth - 16\tabcolsep) * \real{0.1500}}
  >{\raggedleft\arraybackslash}p{(\columnwidth - 16\tabcolsep) * \real{0.1300}}
  >{\raggedleft\arraybackslash}p{(\columnwidth - 16\tabcolsep) * \real{0.1100}}
  >{\raggedleft\arraybackslash}p{(\columnwidth - 16\tabcolsep) * \real{0.1100}}
  >{\raggedleft\arraybackslash}p{(\columnwidth - 16\tabcolsep) * \real{0.1100}}
@{}}
\toprule\noalign{}
\begin{minipage}[b]{\linewidth}\raggedright
Window
\end{minipage} & \begin{minipage}[b]{\linewidth}\raggedleft
\(n\)
\end{minipage} & \begin{minipage}[b]{\linewidth}\raggedleft
\(r\)
\end{minipage} & \begin{minipage}[b]{\linewidth}\raggedleft
\(K\) (kt)
\end{minipage} & \begin{minipage}[b]{\linewidth}\raggedleft
\(g(\mathrm{LRP})\)
\end{minipage} & \begin{minipage}[b]{\linewidth}\raggedleft
\(|e_{q10}|\)
\end{minipage} & \begin{minipage}[b]{\linewidth}\raggedleft
\(C^*\)
\end{minipage} & \begin{minipage}[b]{\linewidth}\raggedleft
\(F'(\mathrm{LRP})\)
\end{minipage} \\
\midrule\noalign{}
\endhead
\bottomrule\noalign{}
\endlastfoot
1983--2007 (frozen) & 24 & 0.2369 & 5000 (ub) & 172.5 & 80.9 & 91.6 & 1.153 \\
NCAM 1995--2015 & 20 & 0.2731 & 5000 (ub) & 198.9 & 27.4 & 171.4 & 1.176 \\
NCAM 1995--2007 & 12 & 0.3901 & 5000 (ub) & 284.0 & 10.0 & 274.0 & 1.252 \\
xteNCAM 1995--2024 & 29 & 0.4310 & 472 (lb) & 49.4 & 41.4 & 8.0 & 0.927 \\
xteNCAM 2005--2024 & 19 & 0.4439 & 472 (lb) & 50.9 & 58.5 & -7.7 & 0.925 \\
\end{longtable}

On the assessment series the answer is yes, and more comfortably than on the
collapse window. On 1995--2015 the fit gives \(r = 0.273\),
\(g(\mathrm{LRP}) = 198.9\) kt and a 10th-percentile floor of only
\(27.4\) kt --- the post-moratorium decline was smooth and therefore largely
predictable one step ahead, the residual SD falling from \(114.9\) kt to
\(18.9\) kt --- so the constructive bound rises to \(171.4\) kt and stays
positive under that window's perpetual-worst floor (\(166.4\) kt), with the
expansion at the LRP intact (\(F' = 1.176\)). On the shorter 1995--2007
window the bound is larger again (\(274.0\) kt).

On the second specification the answer is no. On xteNCAM 2005--2024 the fit
places \(K\) at its lower bound, \(472\) kt, barely above the largest
observed stock, so the fitted surplus at that specification's reference point
is \(50.9\) kt against a 10th-percentile floor of \(58.5\) kt: the
constructive bound is \(-7.7\) kt, and on 1995--2024 it is \(+8.0\) kt.
Under either window's perpetual-worst floor it is negative (\(-37.9\) and
\(-38.8\) kt). The classification reverses with it ---
\(F'(\mathrm{LRP}) = 0.925 < 1\) --- so on the modern series in its recent
state the closed loop contracts at the reference point and the contraction form
of the certified-layer conversion, inapplicable on the frozen object, is the
applicable one.

Two caveats are inseparable from these rows and are stated rather than buried.
The windows are short (12--29 transitions), and in three of the four the
carrying capacity is at a bound, so these are regime readings rather than
estimates of productivity. What they establish does not depend on reading them
as estimates: the reference point's self-viability is regime-dependent ---
comfortable on the assessment series after the moratorium, and at best marginal
on the modern series, at roughly \(0 \pm 8\) kt and negative under that
series' worst observed floor --- and the expansion classification moves with
it. That is consistent with Section 3.11, where the same specification's
1954--2007 fit also failed to make its reference point self-viable.

\subsection{4. Discussion}\label{budget-discussion}

The horizon has a management consequence that is easy to miss while it
is being read as a limitation of the method. A robustness certificate for a
reference point has a shelf life, and Table 2 and Figure 10 show the shelf life cannot be
extended by fishing less: it is six or seven years for every admissible
catch, including none at all. Any interval at which the reference point is
re-assessed therefore has to sit inside it, or the certificate on which the
reference point's defence rests has already expired when it is consulted.
Two real cadences bracket that. The IWC schedules an Implementation Review of
its Revised Management Procedure ``normally'' every six years, and its
aboriginal subsistence strike limits are set in six-year blocks; NOAA's
management-track assessments run on cycles of one to three years for many
stocks, but for others ``a decade may pass before an assessment is
conducted.'' A decadal cadence is outside the horizon computed here; a
six-year one sits exactly at it, and only a cycle shorter than six years is
comfortably inside. The same arithmetic says what would move the horizon:
it is set by the expansion rate \(F'(K^*)\), so neither better data about
the disturbance nor a more cautious catch extends it --- only a reference
point low enough that the governed map contracts there, or a policy that
changes the map's shape rather than its level.

Two layers of negative content must be kept distinct. The vacuous-class
identity (Section 3.2, definitional note) is a robust-layer statement:
under the perpetual-worst persistent floor, no catch policy --- zero
catch included --- holds the LRP. Under the 5th-percentile floor the
statement is no longer vacuous: in the source-year convention that floor
lies below the maximum surplus, and the moratorium's \(T=\infty\) kernel
is nonempty (\(2219.6\) kt). The certified-layer emptiness beyond
\(T = 6\) years under the two harsher floors and beyond \(T = 7\) under
the informative one (Section 3.4) is a different statement, about the
conversion's expansive form: the governed map's expansion rate
(\(F' = 1.153\) at the LRP) empties every certified kernel at those
horizons. Neither result paraphrases the other. The vacuous-class
identity's perpetual-worst reading is form-sensitive --- the
data-preferred depensatory refit restores a nonempty kernel at that
class, the one class vacuous on the registered form (Section 3.6) ---
and its 5th-percentile half is not, that class being informative on both
forms.

Three classes of findings must also be kept distinct --- the form and
specification sensitivities, the geometric findings, and the identified
findings. The form and specification sensitivities carry the same
division. The Fox form preserves every certificate direction at a higher
fit cost (Section 3.6), and the second specification preserves the
expansion classification and the policy ordering while reversing the
informative certificate at its own reference point (Section 3.11). The
geometric findings are form- and specification-independent; the
identified findings are not; the scope is exactly as stated.

\emph{Identified} findings are properties of the fitted map: the
constructive boundary, the expansion rate at the LRP, and the stochastic
survival probabilities of Section 3.8 all move with \(r\) and \(K\). The
boundary-harvesting rules' ``less protective by declared geometry''
reading under the frozen convention is not an identified finding but a
consequence of that convention's harsh floor classes: under the
corrected source-year 10th-percentile class the critical-zone rule, the
cascade, and the graded rules all hold the LRP from itself at every
horizon, so the geometry-based verdict of the frozen reading no longer
applies. What survives is the clause-(H1) structural block at the
5th-percentile \(T=\infty\) class: no positive-catch rule dominates BAU
because every such rule is empty there while BAU's kernel is nonempty
(\(2219.6\) kt). The supply comparison is a second failure, not the
mechanism. The surplus-proportional family is the closest to a
dominance, in that it is both reactive and protective at
\(\phi \le 0.50\) under the informative class while harvesting more than
the moratorium --- a genuine trade-off, not a strict improvement. It
fails of dominance because, on the 5th-percentile class and at
\(T=\infty\), it does not improve on BAU at all, and the
equal-protection flat cap yields a higher mean catch.
\emph{Form-sensitive} findings reverse under the data-preferred
depensatory refit, and the carrying-capacity grid of Section 3.7 shows
the identified findings survive \(K\) within the declared box while the
geometric ones are \(K\)-independent by construction.

The expansion obstruction is also the contribution to the
viability-methods record. The certified-layer machinery used in the
companion studies assumes a contracting closed loop, and the by-catch
fishery case that anchors kernel computation in the fisheries literature
(Krawczyk et al., 2013) is likewise a contracting setting. This cod
object is the first scored instance in which the contraction form is
provably inapplicable at the declared safe set --- the map's steepest
growth occurs exactly where the stock is scarcest, at the boundary the
governance question is about. For a collapsed stock below half its
estimated carrying capacity, every governance statement that survives
certification must therefore be time-bounded and expansion-bound, not
defect-bound.

Four consequences for the methods record follow from the new layers.
First, the carrying-capacity grid shows the expansion classification is
the data-selected regime, not an artifact of the pinned bound, so the
certified-layer emptiness is a property of the identified map rather
than of the conversion's tuning. Second, the finite-duration floors
separate the perpetual-class vacuity from a duration-respecting
robustness question that remains informative. Third, the stochastic
layer converts the worst-case certificates into survival probabilities
on the empirical residual pool --- the quantity a harvest-control
discussion can actually weigh. Fourth, the reactive families extend the
declared policy family with genuinely state-feedback rules, so the ``how
do reactive rules fare'' question is answered directly rather than by
proxy; the surplus-proportional family registers a real trade-off
(protective under the informative class, more harvesting than the
moratorium) while the surplus-harvest and ramp designs are dominated and
recorded as rejected.

The map is one-pool surplus production on annual means. There is no age
structure, migration, or survey catchability (the model-type limitations
of the companion forecast-evaluation study carry over). \(K\) is pinned
at its optimization bound. The expansive classification at the LRP
carries the \(K\)-pin's declared defect with it, in one reading
consistent with Section 3.7 and Figure 4:
\(F'(K^*) = 1 + r(1 - 2K^*/K)\) exceeds 1 exactly while
\(K > 2K^* = 1769.2\) kt. The expansion is conditional on
\(K \ge 2K^*\), and within that range it is not an artifact of the
particular pinned value --- it holds for every admissible
\(K \ge 2K^*\), and the carrying-capacity grid shows the data select
that range (the fit cost and the informative certificates both
deteriorate below \(2K^*\)). Any data-supported \(K\) below twice the
LRP would make the closed loop contract at the boundary and restore the
contraction form of the conversion; the condition is on \(K\), not on
the identified \(r\). The residual conflates productivity shock and
model error (no observation-model separation). The closed loop observes
the stock exactly at the decision instant; real governance operates
under assessment lags (the one-year delay module of the companion
forecast-evaluation study), so the reported kernels are upper bounds for
a perfect-observation controller, and delay-aware kernels would be
smaller or empty. The persistent-shock classes are deliberately harsh (a
perpetual floor, not an independent draw). The 10th-percentile class is
the mildest with non-vacuous content, and only the perpetual-worst floor
sits beyond the map's maximum surplus (Section 3.2), which is what makes
its emptiness vacuous as a productivity statement.

The dominance rule's protective clause is structurally conservative
toward the moratorium. Under the source-year informative class the
critical-zone, cascade, and graded rules match BAU's protection exactly,
and the surplus-proportional family at \(\phi \le 0.50\) holds the LRP
at every horizon while harvesting more than the moratorium; but under
the 5th-percentile class at \(T=\infty\) every positive-catch rule is
empty while BAU's kernel is nonempty (\(2219.6\) kt). Because clause
(H1) is read at every reading with empty scored as worst, that
5th-percentile emptiness is what blocks dominance, and only a rule that
is at least as protective as BAU at every reading and improves on the
moratorium somewhere --- none here --- could dominate BAU. The Allee
term is off in the primary specification; the depensation sensitivity of
Section 3.6 shows the certificate directions are stable under the
data-preferred Allee refit, with the expansion obstruction worsening and
class-vacuity narrowing. The safe set's upper edge (\(10^4\) kt, twice
\(K\)) is a computational grid cap, never approached on any reported
kernel path; kernels are written \([s, \infty)\) because the cap is
never approached, and the positive-part floor \([\cdot]_+\) never binds
on any reported kernel path. Sub-LRP cascade stages are declared
scenarios, not verified institutions. The certified layer is vacuous at
observed stock levels.

Nothing here promotes or demotes any forecast module, transfers numbers
from an interval-verified linear template (a companion methodological
study, under review), or pools the extended xteNCAM series. The \(91.59\)
kt boundary is not a quota recommendation; it is the analytic limit of
robustness at one declared shock class. For management, the deliverable
is the risk statement rather than the point: on this map the
moratorium-level removals and zero catch are the only policies that hold
the safe set under the informative class, the largest robust constant
catch is of order \(70\)--\(90\) kt, and the stochastic layer puts a
probability on holding the line (\(0.90\)--\(0.91\) at
zero-to-moratorium removals, falling with catch) that the persistent
floors cannot express. The result does not transfer to Northern cod
productivity or to harvest-control rules in general; it is a scored
comparison on one fitted map, and the LRP's protection in good years is
the property of that map and those classes.

\subsection{5. Conclusions}\label{budget-conclusions}

Scored intervention selection on the fitted Northern cod
surplus-production map yields five findings and two definitional notes,
stated at their actual strength.

\begin{enumerate}
\def\labelenumi{(\arabic{enumi})}
\item
  Under the informative 10th-percentile productivity class,
  moratorium-level removals and zero catch hold the entire safe set at
  every horizon, and the largest robust constant catch is \(91.59\) kt
  yr\(^{-1}\) --- an analytic limit whose bootstrap \(90\%\) interval is
  \([0, 121.1]\) kt and whose 20-year survival probability from the LRP
  is \(0.74\) under i.i.d. resampling (\(0.73\) under block
  resampling, \(0.84\) with the 1992 draw removed), so the operational
  reading is order
  \(70\)--\(90\) kt, not a quota.
\item
  No non-BAU policy dominates BAU, and the mechanism is the clause-(H1)
  reading at the 5th-percentile \(T=\infty\) class rather than boundary
  geometry or, alone, supply: every positive-catch rule is empty there
  while BAU's kernel is nonempty (\(2219.6\) kt). Under the informative
  10th-percentile class the critical-zone, cascade, and graded rules
  match BAU's protection exactly, and the surplus-proportional family at
  \(\phi \le 0.50\) holds the LRP at every horizon and harvests more
  than the moratorium --- so the reactive family registers a genuine
  trade-off, not a strict improvement. It still does not dominate BAU:
  it does not improve on BAU at the harsher-informative class, and an
  equally protective flat 60-kt cap supplies more.
\item
  The certified layer has a finite horizon --- six years under the two
  harsher floors, seven under the informative one --- because the map
  is expansive at the reference point --- for every admissible
  \(K \ge 2K^* = 1769.2\) kt, so not an artifact of the particular
  pinned value, though conditional on \(K \ge 2K^*\) --- and the
  carrying-capacity grid shows this is the data-selected regime within
  that range: contraction is restored only below twice the LRP, where
  the informative certificates themselves collapse.
\item
  The certificates survive the data-preferred depensatory refit, and the
  class-vacuity reading narrows under the Allee form.
\item
  The certificate has a shelf life, and it cannot be bought. The
  certified horizon is six or seven years for every declared policy and
  every constant catch from zero to \(150\) kt --- the
  moratorium included --- because the erosion margin's annual growth
  is \(379\)--\(773\) kt, four times the whole admissible catch
  range at its smallest. A review interval longer than that horizon is
  consulting an expired certificate.
\item
  The comparison is governed by an identity rather than by a tabulation:
  at the reference point a rule's protection margin plus its catch equal
  \(91.59\) kt exactly (Proposition 2.1), so no reactive rule can
  out-supply a flat cap it matches in protection and the no-dominance
  verdict of finding (2) is structural. The certified horizon is likewise
  a crossing, not a search: six or seven years is where the geometrically
  growing erosion margin overtakes the worst-case trajectory.
\item
  The results are convention-dependent in the residual convention only
  in the sense that the source-year convention sets the constructive
  bound at 91.59 kt, brings the 60-kt rules onto the reference point
  itself (their \(T=\infty\) boundary is \(884.6\) kt, against 900.3 kt
  under the registered convention), lengthens the certified horizon (to \(T=6\) under the two
  harsher floors and \(T=7\) under the informative one), and reduces
  the vacuous family to one class; every parameter and every certificate
  direction is otherwise stable.
\end{enumerate}

\emph{Definitional note A (identity).} Only the perpetual-worst floor is
vacuous, and that vacuity is an arithmetic identity of the declared
class --- the floor exceeds the map's maximum surplus, so every
trajectory declines for every catch, zero included --- not an empirical
finding about Northern cod productivity (Section 3.2's note). The
5th-percentile and 10th-percentile classes are informative, and the
5th-percentile moratorium kernel is nonempty (\(2219.6\) kt).

\emph{Definitional note B (rule-level).} The no-dominance outcome of
finding (2) is a selection property of the declared partial order as
much as of the policies: clause (H1), read at every reading with empty
scored as worst, blocks every positive-catch rule at the 5th-percentile
\(T=\infty\) class given BAU's nonempty kernel there, so within the
declared family the outcome could not have gone the other way
(Definition 2.6's note). It is recorded as a note so that it is not read
as an empirical selection finding on a par with the constructive bound
and the stochastic layer.

The methodological content is the protocol itself --- frozen scoring,
layered certificates, the separation of geometric from identified
findings, and the extension of the family to genuinely reactive rules
--- applied to a collapsed-stock reference point. The empirical content
is scoped to the map, the classes, and the rule declared above. On that
object the LRP is protected by good years, but the margin good years
must supply is smaller than the frozen convention implied, and reactive
rule design can protect the LRP but not out-supply a flat 60-kt cap.
Across assessment specifications the findings are
specification-conditional: on the second, unpooled xteNCAM specification
(LRP \(276\) kt) the expansion classification and the kernel ordering
survive while the reference point's self-viability does not --- the
current stock sits between that specification's one- and five-year
zero-catch boundaries (Section 3.11). Across productivity regimes they are
regime-conditional in the same way: the constructive bound is
\(171\) kt on the post-moratorium assessment window and \(0 \pm 8\)
kt on the modern series (Section 3.12), so the reference point's
self-viability is a property of the regime as much as of the rule.

\part*{Part II --- Regime viability on the Northern cod stock}
\label{part:regime}
\addcontentsline{toc}{part}{Part II --- Regime viability on Northern cod}

\maketitle

\begin{abstract}
The viability machinery of the obstruction calculus has been developed
and validated on deliberately small exact instances. This paper applies
it to a real assessment record --- the Northern cod stock (NAFO 2J3KL),
read from the published assessment input series, the 2025 reassessment
series, the official removals tables, the long-duration catch
reconstruction, the research-vessel survey index, and an independently
monitored forage-fish covariate --- and certifies, in exact rational
arithmetic, the obstruction structure that the historical collapse
exhibited. The certified object is a \emph{harvest-free multiplier
bracket}: for each annual step, the interval's per-capita non-fishing
growth factor lies in
\(\bigl[\rho,\ \max\{\rho+r,\ \rho/(1-r)\}\bigr]\), where \(\rho\) is
the realized net multiplier and \(r\) the removals fraction, under
either timing convention for the removals accounting. Both collapse
steps certify as harvest-free contractions under every convention
(upper bounds \(0.754\) and \(0.372\)); the reference window's worst
step and the two prelude steps do not certify (brackets straddle
unity) --- a typology, not a blanket verdict. Under the post-moratorium
removals rows (\(1{,}314\) and \(413\) t, at most \(1.4\%\) of the
stock), the brackets certify continued contraction, and the stock fell
by a factor of \(\tfrac{10105}{968} \approx 10.4\) in two years to its
series minimum. The \(276\) kt auxiliary floor separates the reference
window from the collapse era without exception, while the
\(884.6\) kt reference point --- the window's own mean --- has three of
the seven window readings below it in sample and therefore separates
nothing. Three extensions certify robustness and scope: on the
reassessment vintage's published 95\% bands, exactly the four
collapse-era steps 1991--1995 certify as declines and the 1993 floor
breach holds at the band's upper edge (\(115 < 276\)), while no
reference-decade step certifies; the reference decade's window profile
shows every subwindow of horizon four or more growing (the sustained-decline
reading holds only at horizons one to three); and the vintage's own
1954--1982 record has the window opening at \(28.3\%\) of the 1962 peak,
three years after a prior floor episode (1977--1979). The collapse is present in the raw survey index (a certified
factor of \(\tfrac{21797}{2127417}\) over 1989--1994, falling \(2.6\)
times further than the assessment series over their common window);
the forage-fish covariate collapses in the same annual step in which
the mortality diagnostic crosses \(M=1\), two steps before the floor
breach; and both assessment vintages re-cross the floor at the 2015
reading, with outcome years at most
\(\tfrac{44000}{88460} \approx 49.7\%\) of the reference point. Every
quantity is regenerated by an accompanying verification script in
exact rational arithmetic.
\end{abstract}

\textbf{Keywords:} viability; obstruction; regime shift; Northern cod;
exact computation; monitoring design; limit reference points

\section{Introduction}\label{regime-intro}

The viability and obstruction machinery of the calculus --- survivable
sets, robust multipliers, breach-time certificates, and probe design ---
is developed in Abaee (2026, An obstruction calculus) and validated on deliberately small exact
instances. This paper applies it to the Northern cod stock (NAFO
divisions 2J3KL) --- the running record of this paper --- and certifies
the obstruction structure of its collapse on the public assessment
record. The forecast layer of the
same record --- the scoring of competing forecast ladders and the
2022--2023 observation suspension as a review-timing event --- is
certified in a companion paper (Abaee, 2026, The Northern cod forecast ladder, Section 3.10); the
present paper certifies the viability layer. The two are companions by
construction: the forecast ladder prices what can be predicted, this
paper prices what could have been kept safe.

The claims below are of three kinds, and the label travels with the
claim. \emph{Arithmetic claims} are exact rational identities or
inequalities on the certified rows; they are regenerated by the
verification script and cannot drift. \emph{Provenance claims} state
which file, which vintage, and which column a number comes from, and
whether that column is a model input or a model output. \emph{Scope
claims} are conditional readings --- stress bounds, sustained
hypotheticals, design critiques --- and are labelled as such where
they appear. No claim of the first kind rests on an interpretation; no
claim of the third kind is stated as one of the first. The
contributions are: the record spine --- the certified 1960--2023
assessment-record reconstruction with its provenance ledger (Sections
\ref{regime-spine}--\ref{regime-windows}); the harvest-free multiplier brackets ---
the exact two-form accounting bracket isolating the non-fishing growth
factor (Section \ref{regime-brackets}); and the window profile and band
certificates --- the collapse's obstruction structure certified on the
public record, with the certified-record figure reading every annotated
value from the locked files (Section \ref{regime-profile}).

\section{Related work}
\label{regime-priorart}

This paper certifies that specific steps in the published Northern cod record were
contractions regardless of removals. Read carelessly, that contradicts the empirical
literature on this collapse, and the conflict is addressed here rather than left for a
referee to find.

\paragraph{The received attribution.} Hutchings and Myers (1994) concluded that the collapse
of northern cod could be attributed solely to overexploitation. Myers and Cadigan (1995)
tested the competing hypothesis that an increase in natural mortality in the first half of
1991 caused the collapse, did not support that hypothesis, and concluded that overfishing was
sufficiently high to cause the collapse of the population even allowing for the possibility
of elevated recent natural mortality. Myers, Hutchings and Barrowman (1997) rejected the
hypothesis that mortality of young cod is unrelated to fully recruited fishing mortality,
finding high juvenile mortality associated with high adult mortality in each of six stocks.

\paragraph{Why the two are consistent, and why the distinction matters.} Those papers ask
about the \emph{historical cause} of the collapse as an event. This paper asks a different
question of a different object: given the transitions recorded in the published assessment,
was the step a contraction even had removals been zero? That is a statement about the
\emph{arithmetic of the record}, not about the cause of the collapse. Removals superimposed
on an already-contractionary productivity term make the outcome worse, not better, and
Proposition~\ref{regime-prop:moratorium} says exactly that --- the loss exceeds the removals by
certified factors while the fishery was closing. Overfishing can be the cause of the collapse
and the individual annual steps can still be contractions at zero harvest. Both hold here.

\paragraph{Where a real tension remains, and is not resolved.} Myers and Cadigan (1995)
specifically declined to support elevated natural mortality as the explanation. To the extent
that the harvest-free contractions certified here are read as evidence of a productivity
collapse, they sit against that finding, and \textbf{this paper does not adjudicate it}: the
certificates are statements about the transitions as recorded, and they do not identify a
biological mechanism. A reader who takes the harvest-free result as evidence about
productivity is going beyond what is certified.

\paragraph{Which series the certificates are computed on.} The concern is concrete rather
than hypothetical. Myers, Hutchings and Barrowman (1997) documented that for these stocks the
VPA-based and survey-based abundance trends diverge --- from the early 1980s the VPA-based
trend shows a decline where none is apparent in the survey-based trend --- and attributed the
divergence plausibly to discarding of young fish. Any certificate computed on one series
inherits that series' artefacts, so the choice of series is load-bearing for this paper and
cannot be left implicit.

This paper therefore checks the raw survey index directly. The collapse is present there too
(Proposition~\ref{regime-prop:survey}), with a certified factor of $\tfrac{21797}{2127417}$ over
1989--1994, and it falls $2.6$ times further than the assessment series over their common
window. The contraction is accordingly robust to the choice of series, and is if anything
stronger in the raw survey than in the assessment record on which the main certificates are
computed. That result is the answer to the divergence concern, and it is stated here so that
the reader can see the concern was met rather than assumed away.

\paragraph{Position within viability theory.} That the governing question is which states
admit controls keeping every constraint satisfied, rather than which policy optimises, is the
founding move of viability theory (Aubin, 1991); the obstruction calculus used here is the
observation-constrained development of that programme (Abaee, 2026), and its constructive
counterpart on this same record is certified in the companion paper (Part~I). What this paper
contributes to that lineage is the \emph{direction} of the certificates --- necessity rather
than sufficiency --- and the fact that they are obtained in exact rational arithmetic with no
fitted model at all.

\section{Data and provenance}\label{regime-data}

Seven locked public files are used, each cross-verified against its
primary source at archival time (Abaee, 2026, Calibration data record): (i) the NCAM (Northern Cod Assessment Model) 2016
assessment input series (DFO, 2016, Table A2): spawning stock biomass
in kilotonnes, 1983--2015, cross-checked 33 of 33 years against the
primary report --- the \emph{input vintage} of the biomass record;
(ii) the xteNCAM 2025 reassessment SSB series (Regular et al., 2025,
Table 17), 1954--2024, with 95\% bands and SSB-over-\(B_{\mathrm{lim}}\)
ratios --- the \emph{reassessment vintage}; (iii) the official
removals table (Regular et al., 2025, Table 1), 1954--2023; (iv) the
catch reconstruction (Schijns et al., 2021, Table 1); (v) the fall
research-vessel survey abundance index (Schijns et al., 2021, Table 3,
citing DFO 2021b) --- a research-vessel input to the assessments, not
an output of them; (vi) the RAM Legacy database (Ricard et al., 2012)
per-year extract for this stock (the free 2024 release), 1850--2021 --- the \emph{database
vintage}, extending the series to 2021; and (vii) the capelin acoustic
biomass series (observed years only, no interpolation).

Two thresholds appear throughout, and their provenance is stated once
here. The \emph{auxiliary floor} \(B_{\mathrm{aux}} = 276\) kt is the
reassessment vintage's \(B_{\mathrm{lim}}\) at the file's rounding: the
published ratio column is rounded to two decimals, and the values it
implies for \(B_{\mathrm{lim}}\) range from \(250\) to \(289\) kt
across the file's years. Every floor comparison below is insensitive
across that range (the compared quantities clear it by large factors;
Remark~\ref{regime-rem:xtevintage}), so a single representative value is
fixed. The \emph{reference point} \(B_{\mathrm{ref}} = 884.6\) kt is
the input vintage era's limit reference point, whose documented
construction is the 1983--1989 mean of the assessment input series
(DFO, 2016); the exact mean is
\(\tfrac{44229}{50} = 884.58\) kt (Proposition~\ref{regime-prop:windows}),
and \(884.6\) is its published rounding.

Two discipline rules are respected throughout. First, the mortality
and fishing-mortality columns of the assessment input file are
diagnostics only (dynamics are never built on them). Second, every
annual multiplier certified below is a \emph{realized net} quantity ---
it conflates the mortality regime with all removals. Section~\ref{regime-brackets}
introduces the one decomposition this paper performs on it, which is
an accounting bracket, not a biological model.

\section{The record spine}\label{regime-spine}

Table~\ref{regime-tab:spine} tabulates the certified rows of the collapse era
and its surroundings: assessment input SSB (input vintage), official
removals, the survey index, the covariate (observed years only), and
the mortality diagnostic. Every proposition below reads its numbers
from this table's files.

\begin{table}[t]
\centering
\scriptsize
\setlength{\tabcolsep}{2.8pt}
\caption{The record spine, 1983--1995. SSB in kt (assessment input
vintage, Table A2); removals in t (official table; the reconstruction
gives identical rows for 1992--1995); survey index in the
research-vessel file's index units; covariate in kt, observed years
only; \(M\) is the assessment's mean natural-mortality diagnostic
(model-estimated). Computed entries are exact rational arithmetic on
the file rows; displayed values are rounded.}
\label{regime-tab:spine}
\begin{tabular}{@{}r r r r r r@{}}
\toprule
Year & SSB (kt) & Removals (t) & Survey & Covariate & \(M\) \\
\midrule
1983 & 841.08 & 232{,}345 & 2{,}088{,}958 & --- & 0.391 \\
1984 & 863.23 & 232{,}471 & 2{,}198{,}605 & --- & 0.377 \\
1985 & 902.94 & 231{,}293 & 1{,}288{,}360 & 3426 & 0.349 \\
1986 & 836.00 & 266{,}713 & 2{,}502{,}702 & 3697 & 0.277 \\
1987 & 940.75 & 239{,}924 & 1{,}020{,}462 & 2576 & 0.494 \\
1988 & 886.40 & 268{,}677 & 1{,}223{,}314 & 4285 & 0.335 \\
1989 & 921.66 & 253{,}990 & 2{,}127{,}417 & 3712 & 0.289 \\
1990 & 861.92 & 219{,}452 & 1{,}627{,}647 & 5783 & 0.403 \\
1991 & 734.51 & 172{,}012 & 1{,}117{,}670 & 138 & 1.002 \\
1992 & 381.95 & 40{,}956 & 239{,}740 & 138 & 2.214 \\
1993 & 101.05 & 11{,}392 & 90{,}709 & --- & 2.575 \\
1994 & 30.55 & 1{,}314 & 21{,}797 & --- & 2.331 \\
1995 & 9.68 & 413 & 43{,}240 & --- & 0.288 \\
\bottomrule
\end{tabular}
\end{table}

\section{The reference window and the collapse window}\label{regime-windows}

\begin{proposition}[seven readings, six transitions; the windows'
worst multipliers]
\label{regime-prop:windows}
From the Table A2 series: the reference window 1983--1989 contains
seven readings \((841.08, 863.23, 902.94, 836.00, 940.75, 886.40,
921.66)\) kt and six annual transitions; their minimum is the 1985
transition \(\rho_{\mathrm{g}} = 836.00/902.94 =
\tfrac{41800}{45147} \approx 0.9259 < 1\), and their exact mean is
\(\tfrac{44229}{50} = 884.58\) kt. The collapse window 1991--1993
contains two transitions,
\(\tfrac{38195}{73451} \approx 0.520\) (1992) and
\(\tfrac{2021}{7639} \approx 0.265\) (1993), each with numerator
strictly below denominator. The connecting 1990 reading is
\(861.92\) kt, with transitions
\(\tfrac{86192}{92166} \approx 0.935\) (1990) and
\(\tfrac{73451}{86192} \approx 0.852\) (1991).
\end{proposition}

\begin{proof}
Each multiplier is the exact ratio of consecutive Table A2 SSB entries,
computed in the script as a rational of the printed values; the
minimum of the six reference-window ratios and the two collapse-window
ratios are the displayed fractions by direct comparison. Each
displayed fraction is below \(1\) because its numerator is smaller
than its denominator (\(41800 < 45147\), \(38195 < 73451\),
\(2021 < 7639\)). The mean is the exact rational sum of the seven
readings over seven.
\end{proof}

\section{Harvest-free multiplier brackets}\label{regime-brackets}

The net multiplier \(\rho\) folds removals into the per-capita rate,
so a sub-unitary \(\rho\) alone does not certify a non-fishing
contraction: with a large enough harvest, even a growing stock
declines net of catch. This section isolates the harvest term by an
accounting bracket that is exact and robust to the one ambiguity the
rows leave open --- where in the interval the removals acted.

\medskip
\noindent\textbf{Shared symbols across the companion papers.} The
companion papers share this vocabulary with paper-local meanings; the
letters below carry meanings here that differ there:
\begin{center}\small
\begin{tabular}{@{}p{0.40\linewidth}p{0.52\linewidth}@{}}
\toprule
\textbf{Here} & \textbf{In the companions} \\
\midrule
\(\rho = B_{t+1}/B_t\): the net biomass multiplier &
the mesh relaxation parameter (certification); the contamination radius
\(\rho\) of \(\mathcal{D}_{\rho}\) (sufficiency; minimax duals) \\
\(r = C/B_t\): the catch share &
the intrinsic growth rates \(r\), \(r_{i}\) (forecast ladder; the
minimax benchmark) \\
\(B\): the biomass &
the belief/information sets \(B_t\), \(B\) (obstruction calculus;
certification; sufficiency) \\
\bottomrule
\end{tabular}
\end{center}

\begin{lemma}[harvest-free multiplier bracket]\label{regime-lem:bracket}
Fix a certified interval with readings \(B_t, B_{t+1} > 0\) and a
removals row \(C \ge 0\) assigned to the interval; write
\(\rho = B_{t+1}/B_t\) and \(r = C/B_t\). Suppose the interval's
accounting has one of the two extreme forms, with \(g_t > 0\) the
interval's per-capita non-fishing growth factor:
\(B_{t+1} = g_t B_t - C\) (removals deducted after growth), or
\(B_{t+1} = g_t (B_t - C)\) (removals deducted before growth). Then
respectively \(g_t = \rho + r\) or \(g_t = \rho/(1-r)\); in either
case
\[
\rho \;\le\; g_t \;\le\; \max\Bigl\{\rho + r,\ \frac{\rho}{1-r}\Bigr\},
\]
and the upper bound is strictly below \(1\) if and only if both
\(\rho + r < 1\) and \(\rho/(1-r) < 1\).
\end{lemma}

\begin{proof}
First form: \(g_t = (B_{t+1} + C)/B_t = \rho + r\). Second form:
\(g_t = B_{t+1}/(B_t - C) = \rho/(1 - r)\), well-defined since
\(C < B_t\) on every certified interval. Both expressions are
\(\ge \rho\) (as \(r \ge 0\), \(1 - r \le 1\)); the disjunction of the
two forms gives the displayed maximum, and the final equivalence is
direct algebra (\(\rho/(1-r) < 1 \iff \rho < 1 - r\)).
\end{proof}

The lemma's content is accounting, not biology: \(g_t\) still conflates
natural mortality with recruitment (the regime), but no placement of
the \emph{realized} removals within the interval's accounting --- early
or late --- can leave the non-fishing account non-contracting when the
bracket's upper bound is sub-unitary. The rows do not date removals
within the year, so both alignments of the two adjacent rows
\((C_t, C_{t+1})\) are carried; a verdict is certified only when it
holds under both conventions and both alignments.

\begin{proposition}[the collapse steps are harvest-free contractions;
the reference and prelude steps are not certifiable]
\label{regime-prop:brackets}
Applying Lemma~\ref{regime-lem:bracket} with both alignments of the removals
rows (Table~\ref{regime-tab:brackets}): both collapse-window steps certify
harvest-free contraction, the 1992 step under an upper bound of at
most \(0.7542\) and the 1993 step under an upper bound of at most
\(0.3718\), each computed in exact rational arithmetic. The reference
window's worst step (1985--1986, upper bound \(1.3140\)) and the two
prelude steps (1989--1990, upper bound \(1.2274\); 1990--1991, upper
bound \(1.0646\)) have brackets that straddle unity: for those steps
the certified rows do not decide between a fishing-dominated and a
non-fishing-dominated contraction.
\end{proposition}

\begin{proof}
Each row of Table~\ref{regime-tab:brackets} evaluates
\(\max\{\rho + r, \rho/(1-r)\}\) as an exact rational for the displayed
\((\rho, C)\) pair; the collapse-window upper bounds (at most
\(0.7542\)) are compared against \(1\) by
numerator--denominator comparison in the script. The prelude and
reference rows exceed \(1\) likewise.
\end{proof}

\begin{table}[t]
\centering
\scriptsize
\setlength{\tabcolsep}{3.2pt}
\caption{Harvest-free multiplier brackets (Lemma~\ref{regime-lem:bracket}),
both alignments of the removals rows. \(\rho\) exact to four decimals;
upper bound \(=\max\{\rho+r, \rho/(1-r)\}\) in exact rational
arithmetic, displayed rounded.}
\label{regime-tab:brackets}
\begin{tabular}{@{}l l r r l@{}}
\toprule
Step & \(C\) (t) & \(\rho\) & Upper bound & Verdict \\
\midrule
1985$\to$86 & 266{,}713 & 0.9259 & 1.3140 & straddles \\
1989$\to$90 & 219{,}452 & 0.9352 & 1.2274 & straddles \\
1990$\to$91 & 172{,}012 & 0.8522 & 1.0646 & straddles \\
1991$\to$92 & 172{,}012 & 0.5200 & 0.7542 & certified \\
1991$\to$92 & 40{,}956 & 0.5200 & 0.5758 & certified \\
1992$\to$93 & 40{,}956 & 0.2646 & 0.3718 & certified \\
1992$\to$93 & 11{,}392 & 0.2646 & 0.2944 & certified \\
1993$\to$94 & 1{,}314 & 0.3023 & 0.3153 & certified \\
1994$\to$95 & 413 & 0.3169 & 0.3304 & certified \\
\bottomrule
\end{tabular}
\end{table}

\section{The window profile and the band certificates}\label{regime-profile}

The reference window's profile answers the selection question at all
horizons, and the reassessment vintage's published bands answer the
robustness question the point estimates leave open.

\begin{proposition}[the reference window's profile: declines only at
horizons one to three]\label{regime-prop:profile}
Within the reference window 1983--1989, the minimum end-to-start ratio
of a \(k\)-year subwindow is: \(k=1\):
\(\tfrac{41800}{45147} \approx 0.9259\) (the worst step, at 2 of the 6
subwindows below unity); \(k=2\): \(\tfrac{83600}{86323} \approx
0.9685\) (2 of 5 below); \(k=3\): \(\tfrac{44320}{45147} \approx
0.9817\) (2 of 4 below); \(k = 4, 5, 6\): every subwindow grows ---
the minimum products are \(\tfrac{15361}{15049} \approx 1.0207\),
\(\tfrac{22160}{21027} \approx 1.0539\), and
\(\tfrac{15361}{14018} \approx 1.0958\). Fifteen of the twenty-one
subwindows are all-growth at the annual level. The sustained-decline
reading of the reference window therefore holds only at horizons one
to three; at every horizon of four or more years inside the window,
the record grew.
\end{proposition}

\begin{proof}
For each fixed \(k\), the subwindow products are exact rationals with
common denominator structure \(S_{t+k}/S_t\), compared by
cross-multiplication in the script; the \(k \ge 4\) growth statements
are the direct comparisons \(S_{t+k} > S_t\), which involve no roots.
The subwindow count is \(6+5+4+3+2+1 = 21\).
\end{proof}

\begin{proposition}[band certificates: exactly the collapse-era steps
certify, and the 1993 breach holds at the band's upper edge]
\label{regime-prop:bands}
The reassessment vintage publishes 95\% bands \((B_{\mathrm{lo}},
B_{\mathrm{hi}})\) per year. A band-certified decline is a step whose
band interval contains only values below unity:
\([\,\tfrac{B_{\mathrm{lo}}(t+1)}{B_{\mathrm{hi}}(t)},
\tfrac{B_{\mathrm{hi}}(t+1)}{B_{\mathrm{lo}}(t)}\,] \subset (0,1)\).
Over the vintage's 1954--2024 record, the band-certified declines are
exactly the four collapse-era steps: 1991--1992
(\([\tfrac{265}{754}, \tfrac{468}{478}]\), upper bound \(0.9791\)),
1992--1993 (\([\tfrac{54}{468}, \tfrac{115}{265}]\), upper bound
\(0.4340\)), 1993--1994, and 1994--1995 --- and no other step in
seventy years, and no reference-decade step. The 1993 floor breach is
band-certified (\(B_{\mathrm{hi}}(1993) = 115 < 276\)); the 1992
reading is indeterminate on the bands (\(265 \le 276 \le 468\)).
\end{proposition}

\begin{proof}
Each candidate step reduces to the single comparison
\(B_{\mathrm{hi}}(t+1) < B_{\mathrm{lo}}(t)\), certified by exact
rational comparison across all seventy pairs of consecutive years in
the script; the interval endpoints are the displayed integer ratios.
The floor comparisons are direct rational comparisons.
\end{proof}

\begin{remark}[the vintage dependence of the typology]
\label{regime-rem:vintagetypology}
The reassessment vintage reads the reference decade itself as
all-growth (its 1983--1989 multipliers are \(1.14, 1.22, 1.03, 1.28,
1.10, 1.01\), minimum \(1.01\)), with the decline beginning at
1989--1990 (\(738/864 \approx 0.854\)): the worst-step typology of
Proposition~\ref{regime-prop:brackets} is vintage-internal, and the
paper's decade statements are made in the input vintage, whose
worst step (\(0.9259\)) has no counterpart in the reassessment
vintage. A third vintage (the database extract's assessment, 2021)
reads 1990--1993 as \(765, 747, 354, 83\) kt; all three vintages
agree on the breach sign pattern against the floor (above, above,
below).
\end{remark}

\begin{remark}[the record before the window: 1954--1982]
\label{regime-rem:backward}
The reassessment vintage extends the record to 1954: a peak of
\(1508\) kt at 1962, with the 1983 window-opening reading at
\(427\) kt --- \(28.3\%\) of the peak, vintage-internally. The floor
\(276\) kt has a prior episode: the crossings of the point series are
exactly \(\{1977, 1980, 1993, 2016\}\) (below 1977--1979, recovered
1980--1992, below again 1993--2015, re-crossed 2016), and the 1978
reading is band-certified below the floor. The reference window thus
opens three years after a prior recovery, at \(28.3\%\) of the
record's peak scale: the window's start is a data artifact of the
input vintage's table, and the design critique of
Proposition~\ref{regime-prop:discrimination} applies to the window's
placement as well as its mean.
\end{remark}

\section{The moratorium obstruction and the post-moratorium
experiment}\label{regime-moratorium}

\begin{proposition}[the collapse proceeds with the fishery closing;
the loss exceeds the removals by certified factors]
\label{regime-prop:moratorium}
Across 1991--1992 the stock lost \(734.51 - 381.95 = 352.56\) kt while
removals were \(40{,}956\) t: the loss exceeds the removals by a
certified factor of \(\tfrac{352560}{40956} \approx 8.61\). Across
1992--1993 the loss was \(381.95 - 101.05 = 280.90\) kt while removals
fell to \(11{,}392\) t: a certified factor of
\(\tfrac{280900}{11392} \approx 24.66\); removals were \(4.06\%\) of
the loss. By Lemma~\ref{regime-lem:bracket} the two steps are harvest-free
contractions under every timing convention
(Proposition~\ref{regime-prop:brackets}): the decline over the collapse window
is not explainable by the realized removals under any accounting
placement of them.
\end{proposition}

\begin{proof}
The losses are differences of certified readings (exact decimals);
the factors are their exact ratios against the removals rows, computed
in the script as rationals (\(4.06\% = 11392/280900\)). The
harvest-free attribution is Proposition~\ref{regime-prop:brackets}'s
certified rows for 1991--1992 and 1992--1993. The statement is a
witness about the realized window, not a counterfactual claim about
other catch paths.
\end{proof}

\begin{proposition}[the reconstruction and the moratorium years]
\label{regime-prop:catchrecord}
The reconstruction's rows for 1992--1995 are
\(40{,}956, 11{,}392, 1{,}314, 413\) t --- identical to the official
removals table's rows for the same years (a cross-file equality) --- and
its 2015 row (\(4{,}436\) t) agrees with the official table and with
the assessment report, as recorded at archival time. The
moratorium-era removals are at least \(176\) and at most \(651\)
times below the 1983--1989 removals rows (all above \(130{,}000\) t in
the same file): a separation of between two and three orders of
magnitude (\(\tfrac{231293}{1314} \approx 176.0\) to
\(\tfrac{268677}{413} \approx 650.5\)).
\end{proposition}

\begin{proof}
The four 1992--1995 rows are read from both removals files and their
equalities are exact rational equalities in the script; the 2015
equality is \(4436 = 4436\). The scale comparison is the exact minimum
and maximum of the four ratios of the 1983--1989 reconstruction rows
against the two moratorium rows.
\end{proof}

\begin{proposition}[the post-moratorium experiment]
\label{regime-prop:experiment}
From the 1993 reading \(101.05\) kt, the certified removals rows are
\(1{,}314\) t (1994) and \(413\) t (1995): at most \(1.4\%\) of the
interval's start reading in both steps (\(r = 1314/101050 \approx
0.0130\) and \(413/30550 \approx 0.0135\)). Lemma~\ref{regime-lem:bracket}
certifies harvest-free contraction in both steps under both conventions
(upper bounds \(0.3153\) and \(0.3304\)): with the fishery effectively
closed, the stock fell \(101.05 \to 30.55 \to 9.68\) kt, a certified
factor of \(\tfrac{10105}{968} \approx 10.44\) in two years, and
\(9.68\) kt is the minimum of the entire 1983--2015 Table A2 series.
\end{proposition}

\begin{proof}
The removals fractions and the bracket upper bounds are exact
rationals in the script, compared against \(1\) by
numerator--denominator comparison. The two-year factor is the exact
ratio of the 1993 and 1995 readings; the minimum is the exact minimum
of the 33 certified readings.
\end{proof}

\section{Mechanized verification}
\label{regime-mech}

The multiplier-bracket algebra of Section~\ref{regime-brackets} is formalized in the project's
Lean~4 layer, in namespace \texttt{Formalizations.ARV} (\texttt{ARV\_RegimeViability}):
\texttt{bracket\_upper\_form1}, \texttt{bracket\_upper\_form2},
\texttt{bracket\_sandwich\_form1}, \texttt{bracket\_sandwich\_form2},
\texttt{bracket\_upper\_lt\_one\_iff}, \texttt{bracket\_subunitary\_form1},
\texttt{bracket\_subunitary\_form2}, \texttt{bracket\_subunitary\_of\_upper\_form1} and
\texttt{bracket\_subunitary\_of\_upper\_form2} --- thirteen theorems, together with the
order-maximum primitives (\texttt{omax} and its six lemmas) on which the interval reading
rests. The two \emph{forms} correspond to the two ways of writing the bracket, and the
subunitarity theorems are exactly the step that turns an upper bracket into a certified
contraction.

\noindent\textbf{No admitted gaps.} A comment-stripped source scan of the whole project finds
no \texttt{sorry}, \texttt{admit}, \texttt{axiom} or \texttt{constant} declaration anywhere;
declarations using classical reasoning depend only on \texttt{Classical.choice},
\texttt{propext} and \texttt{Quot.sound}.

\noindent\textbf{Scope, stated plainly.} Machine-checked here is the \emph{bracket algebra}:
the ordered-field reasoning that turns readings and a removals row into a certified
contraction, and the subunitarity step. Not machine-checked is everything upstream of it ---
the extraction of the certified intervals from the assessment record, the assignation of
removals rows to intervals, and the enumeration of the windows. Those are script computations
and are regenerated, not proved. The formalization therefore certifies that \emph{if} the
readings and the removals row are as stated, the bracket conclusion follows; it does not
certify the readings themselves.

\section{Breach certificates and the cross-vintage check}\label{regime-breach}

\begin{proposition}[breach certificates on the two thresholds]
\label{regime-prop:breach}
With \(B_{\mathrm{aux}} = 276\) kt and \(B_{\mathrm{ref}} = 884.6\)
kt: (i) \emph{conditional stress bound}: from the end-1991 reading
\(734.51\) kt, one step at the collapse window's worst multiplier
already lands below \(B_{\mathrm{aux}}\), since
\(\tfrac{2021}{7639} < \tfrac{27600}{73451}\); this is a hypothetical
sustained-rate bound, and the realized breach took two steps.
(ii) \emph{first below-floor reading}: the 1992 reading is above the
floor (\(381.95 \ge 276\)) and the 1993 reading is below it
(\(101.05 < 276\)): the realized first breach of \(B_{\mathrm{aux}}\)
is the 1993 reading, and the input vintage then stays below the floor
through 2014 (\(250.12\) kt), re-crossing at 2015
(Proposition~\ref{regime-prop:recovery}). (iii) \emph{reference point}:
\(734.51 < 884.6\): the 1991 reading is already below
\(B_{\mathrm{ref}}\), as is the 1990 reading before it
(\(861.92 < 884.6\)).
\end{proposition}

\begin{proof}
(i) is the exact cross-multiplication
\(2021 \cdot 73451 = 148{,}444{,}471 < 210{,}836{,}400 = 7639 \cdot
27600\) in the script. (ii) is the pair of direct rational comparisons
on the Table A2 rows, stated so that the intermediate year is visible:
the breach is \emph{first} at 1993 only because 1992 is certified
above the floor. (iii) is \(734.51 < 884.6\) and \(861.92 < 884.6\).
No step of the proof uses floating point.
\end{proof}

\begin{remark}[the breach frontier]\label{regime-rem:frontier}
On the realized path from the 1989 reading --- \(921.66 \to 861.92 \to
734.51 \to 381.95 \to 101.05 \to 30.55 \to 9.68\) kt --- the
first-below reading is a function of the threshold \(L\), certified on
intervals: \(L \in (861.92, 921.66]\) first falls below at the 1990
reading; \((734.51, 861.92]\) at 1991; \((381.95, 734.51]\) at 1992;
\((101.05, 381.95]\) at 1993; \((30.55, 101.05]\) at 1994;
\((9.68, 30.55]\) at 1995; and \(L \le 9.68\) is never breached within
1983--2015 (the series minimum is \(9.68\)). The realized threshold
\(B_{\mathrm{aux}} = 276\) sits in the 1993 interval
(Proposition~\ref{regime-prop:breach}(ii)); the reference point
\(B_{\mathrm{ref}} = 884.6\) sits in the 1990 interval, after breaching
repeatedly within its own construction window
(Proposition~\ref{regime-prop:discrimination}).
\end{remark}

\begin{remark}[the cross-vintage check]\label{regime-rem:xtevintage}
The reassessment vintage reads the same three years as
\(600, 352, 79\) kt, with floor ratios \(2.18, 1.28, 0.29\). The two
vintages agree on the sign pattern against the floor --- above, above,
below --- and disagree on levels by certified factors
\(\tfrac{73451}{60000} \approx 1.22\),
\(\tfrac{38195}{35200} \approx 1.09\), and
\(\tfrac{10105}{7900} = \tfrac{2021}{1580} \approx 1.28\): the
reassessment's 1993 level is \(21.8\%\) below the input vintage's, and
its \(0.29\) ratio is its own level (\(79\) kt) over its own
\(B_{\mathrm{lim}}\)-scale denominator, not the input vintage's
\(101.05\). Because the ratio column is rounded to two decimals, the
implied \(B_{\mathrm{lim}}\) ranges over \(250\)--\(289\) kt across
the file; the sign pattern would flip only if the true threshold were
at or below \(101.05\) kt or above \(352\) kt, far outside that range,
so every floor comparison here is insensitive to the rounding. The
cross-file statement certified is therefore agreement on the breach
sign pattern, not equality of levels --- and the level disagreement
itself (\(\le\) a factor of \(1.28\) between vintages six years
apart) is part of the certified record.
\end{remark}

\section{The instruments}\label{regime-instruments}

\begin{proposition}[the collapse is present in the raw survey index]
\label{regime-prop:survey}
The fall research-vessel survey index reads \(2{,}127{,}417\) (1989),
\(1{,}117{,}670\) (1991), \(239{,}740\) (1992), \(90{,}709\) (1993),
and \(21{,}797\) (1994): a certified factor of
\(\tfrac{21797}{2127417} \approx 0.0102\) over 1989--1994, and
\(\tfrac{239740}{1117670} \approx 0.214\) across 1991--1992. The index
is a research-vessel input to the assessments, not an output of them:
the collapse is visible in the raw instrument record, before any model
processing. Over the common window 1989--1993 the index falls by
\(\tfrac{90709}{2127417} \approx 0.0426\) while the assessment series
falls by \(\tfrac{10105}{92166} \approx 0.1096\): the instrument falls
\(2.57\) times further than the model output it feeds.
\end{proposition}

\begin{proof}
All five values are rows of the locked survey file; the three factors
are their exact ratios, and \(2.57\) is the exact ratio
\(\tfrac{10105 \cdot 2127417}{92166 \cdot 90709}\) of the two
common-window factors, evaluated in the script. The provenance
statement is a property of the file (survey rows are assessment
inputs, recorded at archival time). The index-to-biomass mapping
assumes constant catchability and coverage; the divergence factor is
certified as a difference between the two series' rates, not attributed.
\end{proof}

\begin{proposition}[the covariate crossing is contemporaneous with the
mortality crossing]\label{regime-prop:capelin}
In the observed capelin series, the acoustic estimate falls from
\(5783\) kt (1990) to \(138\) kt (1991) --- a certified factor of
\(\tfrac{138}{5783} \approx 0.024\) --- and reads \(138\) kt again in
1992. The Table A2 mean natural-mortality diagnostic crosses
\(M = 1\) across the same annual step (\(0.403\) in 1990 to \(1.002\)
in 1991; diagnostic column, model-estimated). At the data's annual
resolution the two crossings coincide; the covariate collapse precedes
the certified floor breach (the 1993 reading) by two annual steps and
the assessment series' steepest realized step (1991--1992) by one.
\end{proposition}

\begin{proof}
Both covariate readings are observed rows of the locked capelin file
(no interpolation, per the archival discipline); the factor is their
exact ratio. The mortality crossing is the diagnostic column's
consecutive values straddling \(1\) across the same pair of years as
the covariate collapse --- the rows display the coincidence directly,
and no sub-annual timing is claimed. Because the diagnostic is
model-estimated given assumed catch, its crossing corroborates the
regime reading but is not an independent certification of it. The
covariate is a single acoustic instrument; no comparability guarantee
across the early-1990s redistribution years is claimed, and the
certified statements are about the locked rows at annual resolution.
\end{proof}

\section{The thresholds as design objects}\label{regime-design}

\begin{proposition}[the auxiliary floor separates the eras; the
reference point does not]\label{regime-prop:discrimination}
Against \(B_{\mathrm{aux}} = 276\) kt: all seven reference-window
readings lie above the floor (minimum \(836.00\)), and the 1993--1995
readings (\(101.05, 30.55, 9.68\)) lie below it, with 1990--1992
above: the auxiliary floor separates the reference window from the
collapse era without exception. Against
\(B_{\mathrm{ref}} = 884.6\) kt: three of the seven reference-window
readings (\(841.08, 863.23, 836.00\)) already lie below, and within
1983--1995 the below-set is \(\{1983, 1984, 1986, 1990, \dots, 1995\}\)
--- nine of thirteen readings: the reference point does not separate
the eras. A threshold fixed at a window's own mean places roughly half
that window's readings below it by construction and retains no
headroom above it (the window's maximum is \(6.4\%\) above; its
minimum \(5.5\%\) below, outward-rounded).
\end{proposition}

\begin{proof}
All comparisons are direct rational comparisons of the certified
readings against \(276\) and \(\tfrac{8846}{10}\) in the script; the
below-sets are computed by exact membership. The mean-rule statement
is Proposition~\ref{regime-prop:windows}'s exact mean identity: a threshold
at the mean of a set of readings lies between the set's minimum and
maximum, so the set straddles it.
\end{proof}

\begin{proposition}[conditional: sustained worst-case contraction
breaches the floor in sixteen steps]\label{regime-prop:sustained}
Under the reference window's worst realized multiplier
\(\rho_{\mathrm{g}} = \tfrac{41800}{45147}\) sustained from the 1989
reading \(921.66\) kt --- a hypothetical rate, held fixed ---
the stock breaches \(B_{\mathrm{aux}} = 276\) kt within sixteen steps:
\(\rho_{\mathrm{g}}^{16} \le \tfrac{27600}{92166} < \rho_{\mathrm{g}}^{15}\).
The count is a ceiling convention (\(15.66\) steps rounds up); the
statement is conditional on the sustained hypothetical rate and
carries no harvest-free reading: the window's own bracket straddles
unity (Proposition~\ref{regime-prop:brackets}), so the sustained-rate
contraction is not attributable to the non-fishing account by the
certified method.
\end{proposition}

\begin{proof}
The sixteenth and fifteenth integer powers of
\(\tfrac{41800}{45147}\) are compared against
\(\tfrac{27600}{92166}\) by exact cross-multiplication in the script.
The ceiling convention is stated because the exact threshold step lies
strictly between the fifteenth and sixteenth powers.
\end{proof}

\section{The outcome years}\label{regime-outcome}

\begin{proposition}[both vintages re-cross the floor at the 2015
reading; outcome years stay below half the reference point]
\label{regime-prop:recovery}
Each vintage's own series re-crosses \(B_{\mathrm{aux}} = 276\) kt at
the 2015 reading: the input vintage goes \(250.12 \to 298.65\) kt
(2014--2015), the database vintage \(238 \to 277\) kt. The database
vintage then reads \(340, 433, 394, 419, 440, 411\) kt for
2016--2021 (six outcome years, all above the floor), with maximum
\(440\) kt at \(\tfrac{44000}{88460} \approx 49.7\%\) of
\(B_{\mathrm{ref}}\). The recovery's certified removals rows are
\(4.4\) to \(12.9\) kt per year (2015--2021, official table; maximum
\(12{,}881\) t in 2017) --- between \(2.3\%\) and \(3.0\%\) of the
same vintage's stock reading per year. The certificate does not
extend to the way up: on the 2020--2021 step (a \(29\) kt decline)
the removals were \(\tfrac{10977}{29000} \approx 37.9\%\) of that
year's decline, so no harvest-free attribution is certified on
recovery steps. The asymmetry the certificate thus documents --- rapid
certified decline against two decades below half the reference point on
the way up --- is the pattern the collapse literature records for this
stock (Hutchings and Myers, 1994).
\end{proposition}

\begin{proof}
The crossing pairs are direct rational comparisons in the two locked
files; the six outcome values are the database rows (exact integers in
kt); the maximum comparison is the reduced ratio
\(\tfrac{44000}{88460}\). The removals range is the exact minimum and
maximum of the seven official rows for 2015--2021.
\end{proof}

\begin{remark}[the monitoring consequence]\label{regime-rem:monitoring}
Read together: the certified obstruction is a property of the
dynamics under the realized catch (Propositions~\ref{regime-prop:brackets}
and~\ref{regime-prop:moratorium}), the raw instrument record shows the
collapse before the model output does (Proposition~\ref{regime-prop:survey}),
the covariate crossing is coincident with the mortality crossing and
two steps ahead of the floor breach (Proposition~\ref{regime-prop:capelin}),
and the mean-rule reference point separates nothing
(Proposition~\ref{regime-prop:discrimination}). The assessment machinery
itself was reconstructed around a monitoring indicator after 2016 (a
three-year review cycle with an interim survey indicator, the
indicator value available by early January) --- the same review-timing
reading the companion paper certifies for the 2022--2023 observation
suspension (Abaee, 2026, The Northern cod forecast ladder, Section 3.10).

\begin{figure}[t]
\centering
\includegraphics[width=0.96\linewidth]{figs_arv/fig_record.pdf}
\caption{The certified record, both vintages: the input vintage
(NCAM 2016, 1983--2015) and the database vintage (RAM v4.66, 2014--2021)
against the auxiliary floor ($276$ kt) and the reference point
($884.6$ kt); reference and collapse windows shaded; the 1993 breach
($101.05$ kt), the 1995 minimum ($9.68$ kt), the 22 consecutive
below-floor readings to 2014, and the 2015 re-crossing
($250.12 \to 298.65$; $238 \to 277$) as certified in the text.
Generated by the assertion script
\texttt{paper2\_arv\_record\_figure\_v1.py} (recompute-then-assert;
eleven exact checks against the locked files).}
\label{regime-fig:record}
\end{figure}
\end{remark}

\begin{table}[t]
\centering
\scriptsize
\setlength{\tabcolsep}{2.8pt}
\caption{The certified quantities and their files. Computed in exact
rational arithmetic on the locked rows; displayed values rounded.}
\label{regime-tab:summary}
\begin{tabular}{@{}l l l@{}}
\toprule
Quantity & Value & Source file \\
\midrule
Ref-window worst mult.\ (1986) & \(41800/45147\) & Table A2 \\
Collapse multipliers (1992, 1993) & \(0.520,\ 0.265\) & Table A2 \\
Collapse-step bracket bounds & \(0.7542,\ 0.3718\) & A2 $+$ T1 \\
Post-moratorium 2-yr factor & \(10105/968 \approx 10.44\) & Table A2 \\
Removals 1992, 1993 & \(40{,}956,\ 11{,}392\) & Table 1 (2025) \\
Removals 1994, 1995 & \(1{,}314,\ 413\) & T1 $=$ reconstr. \\
Survey index 1989, 1994 & \(2{,}127{,}417,\ 21{,}797\) & Survey T3 \\
Covariate 1990, 1991 & \(5783,\ 138\) & Capelin file \\
First below-floor reading & 1993 (\(101.05 < 276\)) & Table A2 \\
Cross-vintage 1993 levels & \(101.05\) vs \(79\) kt & A2 $+$ T17 \\
Outcome years 2016--2021 & \(340\)--\(440\) kt & DB extract \\
\bottomrule
\end{tabular}
\end{table}

\section{Delimitations}\label{regime-scope}

Every arithmetic claim is a certificate about the realized rows of the
certified files --- not a forecast, not a stochastic model, and not a
counterfactual catch analysis. The bracket of
Lemma~\ref{regime-lem:bracket} is an accounting decomposition: it removes the
realized removals from the per-capita rate under two extreme timing
conventions, and its harvest-free factor still conflates natural
mortality with recruitment --- the certified statement is that no
placement of the realized removals explains the decline, not that no
catch path could have. The biomass series are assessment outputs
(model-based), treated as the data of record with their published
uncertainty untouched; the exact arithmetic is exact on the published
point estimates, not a claim about the latent stock. The mortality
diagnostic is model-estimated given assumed catch, so the
mortality-crossing statement corroborates the regime reading without
independently certifying it; the survey index carries its own sampling
error and catchability assumptions, untouched here, and the
instrument--output divergence factor is certified as a difference of
rates, not attributed to a cause. The covariate series is
observed-years only; the contemporaneity statement is about certified
rows at annual resolution. Single stock; two vintages of one threshold
and one vintage of the reference point; no averaging beyond the
reference point's own construction.

\section{Verification methods}\label{regime-methods}

Every number, fraction, inequality, power, ratio, bracket, and
temporal claim in this paper is regenerated by the accompanying script
in exact integer and rational arithmetic, standard library only,
directly from the locked files: the seven reference-window readings
with their six multipliers, exact mean
(\(\tfrac{44229}{50}\)), and minimum; the collapse-window and
connecting multipliers; the harvest-free brackets of
Table~\ref{regime-tab:brackets} under both alignments, each upper bound
compared against unity by numerator--denominator comparison; the loss
factors \(\tfrac{352560}{40956}\) and \(\tfrac{280900}{11392}\) and
the removals share \(11392/280900\); the cross-file equality of the
1992--1995 removals rows and the 2015 row; the moratorium scale ratios
(\(\tfrac{231293}{1314}\) to \(\tfrac{268677}{413}\)); the
post-moratorium brackets and the two-year factor
\(\tfrac{10105}{968}\); the breach comparisons by integer
cross-multiplication (\(2021 \cdot 73451 < 7639 \cdot 27600\)) and the
first-below-floor pair with the 1992 intermediate; the cross-vintage
sign pattern, level factors, and implied-threshold range; the survey
factors and the divergence ratio; the covariate and mortality rows;
the threshold-discrimination below-sets; the sustained sixteenth-step
comparison; the two-vintage re-crossings, outcome values, and removals
range; and every headline number of this text asserted present
verbatim.

\section{Conclusion}\label{regime-conclusion}

Applied to the public record of the Northern cod, the calculus
certifies the following. The collapse window's two annual steps are
harvest-free contractions under every timing convention the rows admit
(upper bounds \(0.754\) and \(0.372\)): no placement of the realized
removals --- which the year of the floor breach puts at \(4.06\%\) of
the stock's loss --- explains the decline, and under the moratorium's
own removals rows the non-fishing account drove a further certified
factor-of-\(10.4\) contraction to the series minimum. The reference
window's worst step and the two prelude steps do not admit that
reading: their brackets straddle unity, so the certified record
supports a typology (two harvest-free collapse steps inside a
fishing-era decline), not a blanket verdict about the decade. The
\(276\) kt auxiliary floor separates the eras without exception,
while the \(884.6\) kt reference point --- fixed at the window's own
mean --- has three of the seven window readings below it in sample
and separates nothing: a design property of mean-rule thresholds,
certified here on the readings themselves. The collapse is present in
the raw survey index at a certified factor of
\(\tfrac{21797}{2127417}\) over 1989--1994, falling \(2.6\) times
further over the common window than the assessment series it feeds;
the covariate crossing coincides with the mortality crossing at annual
resolution and leads the floor breach by two steps. Both assessment
vintages re-cross the floor at the 2015 reading and the outcome years
stay below half the reference point. The certificates are robust
where the bands allow them to be: on the reassessment vintage's
published 95\% bands, exactly the four collapse-era steps certify as
declines and the 1993 breach holds at the band's upper edge; and the
reference window's own profile (declines at horizons one to three,
growth at every horizon of four or more) delimits the sustained-rate
reading the record will bear --- as does the vintage's pre-window
record, which has the window open at \(28.3\%\) of the 1962 peak,
three years after a prior floor episode. These are exact statements about
published point estimates; their scope is the realized record, and
their monitoring content --- which instruments moved, when, relative to
which threshold --- is the applied content the obstruction calculus
prices.

\part*{Part III --- Governance operators and viability kernels of the Edwards Aquifer}
\label{part:edwards}
\addcontentsline{toc}{part}{Part III --- Edwards Aquifer viability kernels}

{\centering\small \textbf{Article type: Research Paper}\par}

\begin{abstract}


\textbf{Problem.} Permitted groundwater systems manage drought with triggers; the protection competing rules deliver under adverse recharge is rarely scored against a fixed criterion.

\textbf{Approach.} We compute robust viability kernels of the J-17 annual-mean head for a policy family --- training-mean pumping (282.16 \ensuremath{\times} 10\textsuperscript{3} acre-ft yr\textsuperscript{-1}), flat caps of 0--90\%, a Stage-I reactive rule, and the critical-period-management cascade --- under persistent drought-floor recharge. A policy is retained only if at least as protective as training-mean pumping and more permissive than the smallest securing flat cap, under a protocol frozen before any score.

\textbf{Results.} At the 618-ft threshold under the drought-of-record floor, training-mean pumping's worst-case attractor (615.72 ft) lies below it, its kernel emptying beyond 13 years; every flat cut of 10\% or deeper and both reactive rules make the safe set invariant; the smallest securing cut is interpolated 7.2\% (31.5\% of current). The reactive rules are retained only \textbf{nominally} at 618 ft under the mildest floor class: their supply margins (+3.3\% Stage I, +0.4\% cascade) are hybrid, worst-case protection against wet-year replay entitlement the robust class excludes. Retention is \textbf{not} certified (positive-pumping certified kernels empty beyond T = 3 years), and \textbf{nothing is retained at 660 ft}, where the rules are invisible to the kernel of their own trigger --- a geometric property of trigger design.

\textbf{Implications.} Reactive rules match flat-cap protection at the physical threshold and seem generous only if wet-year pumping counts; a positive attractor-to-threshold margin defeats every non-negative pumping rule, so the 660-ft threshold is protected by wet years, not by policy; certified claims are limited to three years by the fitted defect.

\end{abstract}

\textbf{Keywords:} Edwards Aquifer; critical period management; J-17
index well; pumping policies; robust viability; viability kernel;
groundwater governance

\textbf{Article Impact Statement.} An intervention-selection protocol
--- governance operators, robust viability kernels under declared
recharge floors, declared-defect erosion, and a frozen retention rule
--- is demonstrated on the Edwards Aquifer. It resolves which pumping
rules protect a physical threshold under adverse recharge and at what
cost in permitted supply: on J-17's 618-ft level, reactive trigger rules
match flat-cap protection at 3.3\% (Stage I sketch) and 0.4\% (cascade)
more permitted supply than the kernel-matched flat-90\% cap --- nominal,
hybrid-supply results --- while the 660-ft institutional threshold is
protected by wet years, not by the pumping family. The protocol is the
generalizable content; the specific verdicts are scoped to this system.

\subsection{1. Introduction}\label{edw-introduction}

Drought management in permitted groundwater systems is built on
triggers. An index-well or springflow reading crosses a threshold, and
pumping reductions activate. Trigger design has a long evaluation
literature: drought indicators and triggers are properly assessed as
stochastic objects, with explicit attention to how indicator levels,
trigger thresholds, and drought stages interact (Steinemann 2003).
Robustness has been a recurring concern of water-resources planning
since the recognition that optimal solutions can be brittle to the
hydrologic assumptions that produce them (Watkins and McKinney 1997).
The operational question for a district that already operates triggers
is narrower, and comparative: given a declared family of pumping rules,
which rules deliver protection under adverse recharge, and at what cost
in permitted supply?

The Edwards (Balcones Fault Zone) Aquifer of south-central Texas is the
most fully institutionalized example of trigger-based groundwater
governance in North America. Its critical-period management program
stages pumping reductions at index-well thresholds. The reductions
protect springflow at Comal and San Marcos Springs, whose
endangered-species requirements are administered through the Edwards
Aquifer Habitat Conservation Plan, the subject of a standing scientific
review by the National Academies (National Research Council 2015). The
aquifer's regional behavior has long been represented --- and debated
--- through lumped and equivalent-porous-media models, from the Barton
Springs groundwater availability model (Scanlon et al.~2003) to its
recalibration by the Texas Water Development Board (Hutchison and Hill
2011). This study scores, rather than asserts, the protection supplied
by a declared family of pumping rules on the measured J-17 record.

The instrument is the robust viability kernel: the set of initial heads
from which a closed-loop pumping rule keeps the system inside a declared
safe set under every disturbance in a declared class. Viability methods
have an established lineage in natural-resource management, where
constraints rather than objectives define the policy question (Krawczyk
and Pharo 2013), and their computational tractability on
management-scale problems is documented by fishery case studies
(Krawczyk et al.~2013). The declared safe set here is a floor on the
productive store rather than on a year's yield. The head threshold marks
the level below which springflow at Comal and San Marcos approaches
cessation. The policy question is therefore the protection of the base
that regenerates the yield, and a trigger's function is to keep the
store from being drawn down to that floor at the cost of near-term
permitted supply. The governed object is the one-pool affine head map
fitted in a companion forecast-evaluation study on the same series
(under separate review), which found last-value persistence unbeaten by
the causal ladder (7.55-ft oracle RMSE against 13.23 ft persistence and
12.84 ft AR(1); Author et al., in review). This study asks the
governance question that the forecast study could not: whether a
declared pumping rule changes the viability kernel of the real system as
represented by the fitted map, and at what cost in permitted supply. A
companion intervention study on a marine fishery stock applies the same
design and reaches the mirror verdict --- none retained there, reactive
rules retained here (nominally, under the mildest floor class); the
mechanism contrast (fast wet-year recovery in this aquifer against a
fatal productivity floor at the cod reference point) is the only reason
to mention it, and the two systems' scores are never pooled (Author et
al., in review).

The complete comparison --- object, defect declaration, disturbance
classes, policy family, and retention rule --- was frozen in a dated
protocol (2026-08-26) before any score was computed.

\subsection{Related work}
\label{edw-priorart}

\paragraph{The regulatory object.} The J-17 index well is not merely a monitoring point: it
is the official regulatory trigger of the Edwards Aquifer Authority's Critical Period
Management plan, and crossing its thresholds automatically imposes pumping restrictions
(Edwards Aquifer Authority, 2024). Under the EAA Act, permitted withdrawals from the aquifer
are capped, and the critical-period stages impose graduated withdrawal reductions on the San
Antonio Pool as the J-17 level and Comal and San Marcos springflows decline; the Stage~I
trigger for the San Antonio Pool is a J-17 level below 660~ft MSL with a 20\% withdrawal
reduction, rising to a 40\% reduction at Stage~IV; the Uvalde Pool is staged separately on
J-27 (Edwards Aquifer Recovery Implementation Program, 2021). \textbf{The 660-ft threshold
studied in this paper is therefore not an arbitrary level chosen for the analysis but the
actual Stage~I trigger}, which is what makes the results below directly interpretable as a
statement about the instrument in force. The 618-ft level is the physical threshold at which
the analysis finds protection to be a different question.

\paragraph{Viability applied to groundwater.} The viability approach has been applied to
aquifer management repeatedly, and the comparison it invites is by now standard: in the
optimal-control framing an environmental objective enters as a monetised externality, while
in the viable framing it enters as a constraint to be satisfied, and the two produce
different trajectories and different long-run levels (Pereau, Pryet and Rambonilaza, 2019).
The viability kernel has been identified analytically in hydro-economic models of aquifers
under conflicting food-security and ecosystem objectives (Pereau, Mouysset and Doyen, 2018),
and the wider survey literature records reservoir and groundwater management as established
application areas (Oubraham and Zaccour, 2018).

What this paper adds to that line is the \emph{robust} element and the geometry of the
trigger. The kernels here are computed against a persistent drought-floor recharge rather
than a nominal one, and the central negative finding --- that \textbf{nothing is retained at
660~ft, because the rules are invisible to the kernel of their own trigger} --- is a
statement about the relationship between a trigger level and the reachable set under it, not
about the relative merits of optimal and viable control. That is a different question from
the one that literature asks, and to our knowledge it is not asked there.

\paragraph{From safe yield to sustainability, which this paper must acknowledge.} It is a
conspicuous feature of the groundwater literature that the classical management concept ---
safe yield, the balancing of annual withdrawal against annual recharge --- was shown to be
unsustainable in principle, precisely because it ignores the natural discharge that recharge
already supports (Sophocleous, 1997; Alley, Reilly and Franke, 1999; Sophocleous, 2000; Alley
and Leake, 2004). The persistence of the underlying water-budget reasoning has been
documented as well (Devlin and Sophocleous, 2005). \textbf{This paper does not use the term,
and should be read against that background.} Its finding that a positive attractor-to-threshold
margin defeats every non-negative pumping rule --- so that the 660-ft threshold is protected
by wet years rather than by policy --- is the robust-viability analogue of the safe-yield
critique: a level that looks protected under average or favourable conditions is not
protected under the adverse recharge the constraint set is required to survive.

\paragraph{Viability theory.} That the governing question is which states admit controls
keeping every constraint satisfied over time, rather than which policy optimises, is the
founding move of viability theory (Aubin, 1991); the obstruction calculus underlying the
certificates used here is the observation-constrained development of that programme.

\subsection{2. Methods}\label{edw-methods}

\subsubsection{2.1 Object and dynamics}\label{edw-object-and-dynamics}

The object is that of the companion forecast evaluation: J-17
annual-mean head H\_t (ft above mean sea level), San Antonio Pool,
1934--2023, drawn from the registered twenty-column analysis panel
(that evaluation's fixed analysis dataset --- the measured head, the
constructed recharge and pumpage series, and the service and derived
series declared there --- of which only the head, recharge, and pumpage
columns enter this analysis). The
head is a measured well level, not an assessment inversion. Recharge R
(USGS-estimated, Puente 1978 method; Umphres and Choi 2025) and pumpage
P (Edwards Aquifer Authority Table 1) are the fluxes. No new data enter
the analysis.

The dynamics are the causal stock-flow class of the companion forecast
study, with one pool and affine form:

\[\Delta H_t = \alpha + \beta R_t + \gamma P_t + \delta H_{t-1}.\]

\textbf{Definition 2.1 (One-pool affine head map).} Let H\_t denote the
annual-mean head (ft above msl) at J-17, R\_t the annual recharge, and
P\_t the annual pumpage. Under the hypotheses (F1) one-pool storage,
(F2) affine dependence on (R\_t, P\_t, H\_\{t\ensuremath{-}1\}), and (F3)
ordinary-least-squares fitting on 1934--1990 (56 transitions), the head
map is the relation above.

The fit is \ensuremath{\alpha} = 163.49, \ensuremath{\beta} = 0.0198 ft per 10\textsuperscript{3} acre-ft, \ensuremath{\gamma} = \ensuremath{-}0.02844 ft
per 10\textsuperscript{3} acre-ft, \ensuremath{\delta} = \ensuremath{-}0.2539, so a = 1 + \ensuremath{\delta} = 0.7461. The contraction a =
0.7461 corresponds to 25.4\% per year mean reversion (the companion's
rolling AR(1) on the same series gives 0.66 --- different window,
different object; the full-sample coefficients there are (0.017, \ensuremath{-}0.026)
against this window's (0.0198, \ensuremath{-}0.02844)). Residual SD is 5.60 ft;
training maximum is 15.41 ft. Out of sample (1991--2023, audit only, no
refitting), the SD is 8.40 ft and the maximum is 21.81 ft. The
out-of-sample maximum exceeds the declared defect. This fact is recorded
below and not repaired, per protocol. The coefficients are fitted on the
largely unregulated 1934--1990 era, so ordinary least squares cannot
separate the aquifer's physical response from the pumpers' historical
drought reactions, and imposing state-dependent rules under those
coefficients assumes the physics--behavior separation the estimator does
not guarantee (Lucas 1976); this bears on whether \ensuremath{\gamma} may be read as a
policy lever at all, and the companion found the sign of \ensuremath{\gamma} reversed on
the 1934--1950 subwindow.

The information pattern is the following: the manager ends year t
knowing (H\_t, R\_t, P\_t) and sets P\_\{t+1\} = \ensuremath{\pi}(H\_t). The recharge
R\_\{t+1\} is unknown at decision time and is treated adversarially
within a declared persistent floor. The kernels therefore assume
year-end J-17 is known exactly; the actual CPM triggers on 10-day means
and cuts within the year, so annual triggering overstates the lag and
understates the cut-active fraction.

\subsubsection{2.2 Uncertainty classes and safe
sets}\label{edw-uncertainty-classes-and-safe-sets}

The disturbance classes are persistent recharge floors computed on the
training window: UC-min = 43.7 (the 1956 drought-of-record year, held
perpetually), UC-q05 = 166.5, and UC-q10 = 179.1 \ensuremath{\times} 10\textsuperscript{3} acre-ft yr\textsuperscript{-1}.
UC-q05 and UC-q10 are the 3rd and 6th lowest recorded annual recharges
of the 57 training years (linear-interpolation percentile estimator);
the floors themselves have historical precedent, and only their
perpetuity is harsher than any recorded year. They are certification
geometry, not recharge forecasts.

The safe sets are the two declared normative thresholds, fixed in the
frozen protocol: \(K^{\ast}_{\mathrm{phys}}\) = 618 ft (Comal Springs
cessation proximity) and \(K^{\ast}_{\mathrm{inst}}\) = 660 ft (the
post-2007 Stage I trigger, not applied to pre-2007 history; a 10-day
mean in the actual institution, an annual mean here). The model domain
is {[}610, 710{]} ft. Upward exits above 710 ft are model-domain exits,
not threshold violations.

\subsubsection{2.3 Governance family}\label{edw-governance-family}

The declared governance family is the following. Training-mean pumping
(called BAU in the frozen protocol; the label is retained in tables and,
under this declaration, in running prose as the family's baseline label)
holds pumping at the 1934--1990 training mean, P \ensuremath{\equiv} \(\bar{P}\) = 282.16 \ensuremath{\times} 10\textsuperscript{3}
acre-ft yr\textsuperscript{-1} --- a flat 100\% cap on the unregulated mean, not the
historical pumpage path and not current use. Flat caps prescribe \ensuremath{\rho}\(\bar{P}\) with
\ensuremath{\rho} \ensuremath{\in} \{0.9, 0.8, 0.7, 0.6, 0.5, 0.0\}. The Stage-I reactive rule (S1)
cuts pumping 20\% when H \textless{} 660 ft; the reduction is a
20\%-at-660 sketch of the Authority's published Stage I (Edwards Aquifer
Authority, Critical Period / Drought Management; the actual rule keys on
10-day J-17 or Comal flow). The critical-period-management (CPM) cascade
prescribes cumulative stage totals of 20/30/35/40\% cuts at H
\textless{} 660/650/640/630 ft (cumulative, not stacked). The actual
program declares five stages (Stage V at 625 ft); only the four-stage
sketch is scored, Stage I is verified against the published 20\%/660
pair, and stages II--IV are declared scenarios, not verified
institutions.

The CPM supply figure below is the stage-weighted evaluation on the
observed occupancies of the training era (1934--1990, n = 57 years):
Stage I (H \textless{} 660 ft) 20 years (35.1\%), Stage II (H
\textless{} 650 ft) 12 (21.1\%), Stage III (H \textless{} 640 ft) 5
(8.8\%), and Stage IV (H \textless{} 630 ft) 1 (1.8\%; 1956). The
cumulative 20/30/35/40\% cascade evaluated on these occupancies
reproduces the replay's mean prescribed pumping exactly (254.93 \ensuremath{\times} 10\textsuperscript{3}
acre-ft yr\textsuperscript{-1}). Out of sample (1991--2023, n = 33 years) the occupancies
are 33.3\%, 15.2\%, 6.1\%, and 0.0\%. Stages II and III recur out of
sample (15.2\%, 6.1\%); Stage IV does not (0.0\%; 1956 alone is below
630 ft), so CPM's deepest cut is a 1956-only scenario in the historical
measure.

\textbf{Definition 2.2 (Robust viability kernel).} Fix a policy \ensuremath{\pi}, a
disturbance class \(\mathcal{D}\), a safe set \(K^{\ast}\), and a horizon T. The
robust T-step viability kernel Viab\_T(\ensuremath{\pi}, \(K^{\ast}\), \(\mathcal{D}\)) is the set of
initial heads from which the closed loop P\_\{t+1\} = \ensuremath{\pi}(H\_t) keeps H
inside \(K^{\ast}\) under every recharge sequence in \(\mathcal{D}\).

For each policy and disturbance class, the robust T-step viability
kernel of the safe set is computed by iterating the worst-case closed
loop. The nominal kernel is reported without erosion. The certified
layer applies the defect-to-margin conversion in the form the fitted map
admits. The term ``kernel'' is used in its closed-loop reading --- the
robust positively invariant set of a fixed declared policy --- not in
the classical viability reading of Aubin (1991), which is existential
over controls. No control choice enters the kernel computation. Only the
evaluation of declared rules under the disturbance floor enters. With
the declared policies all deterministic functions of the observed head,
each robust kernel reduces to the invariance statement of a
one-dimensional closed loop --- attractor and domain-top arithmetic
rather than control synthesis. The word is retained because the object
answers the same question (``from which states can the constraint be
held?'') in the policy-fixed form.

\subsubsection{2.4 Retention rule and certified
layer}\label{edw-retention-rule-and-certified-layer}

The retention rule decides which non-training-mean policies earn their
additional structure. The certified layer converts the declared defect
into an erosion margin on the kernel.

\textbf{Definition 2.3 (Retention rule).} A non-BAU policy is retained
only if (R1) its robust kernel is at least as protective as
training-mean pumping's at every reading, compared on the kernel lower
boundary with empty = worst, and (R2) it permits more pumping than the
smallest securing flat cap with matched protection --- on the declared
10\% grid this resolves as the largest-supply member of the matched
class (flat-90\%); the grid-resolution sensitivity is reported in
Section 3.4 (Table 4).

The scoring regimes are declared here rather than left implicit.
\textbf{Protection is scored under the disturbance classes} --- the
perpetual floors of Section 2.2. \textbf{Supply is scored as the
1934--1990 historical replay mean.} The retention verdict is therefore a
hybrid: worst-case protection at historical-mean entitlement, not a
within-scenario dominance. Under the floor classes alone, the reactive
rules are identically their matched flat caps (the cut is active every
year) and carry no supply margin. The retention decision is made at the
nominal level and re-checked at the certified level. The certified
re-check holds only over the horizons where the certified kernels are
nonempty --- a horizon-truncated object, stated as such whenever it is
invoked. At 618 ft under UC-min, R1 is close to automatic: every cut of
10\% or deeper in the family makes the safe set invariant, so retention
reduces to R2 (supply at matched protection).

\textbf{Calculation 2.1 (Certified erosion margin).} Let \ensuremath{\varepsilon} = 15.41 ft
denote the declared uniform defect (the training maximum residual) and
let a = 0.7461 denote the closed-loop contraction rate. Under the
hypothesis (F4) that the autonomous contraction rate applies uniformly
to the residual, the T-step erosion margin is

\[r_T = \varepsilon \frac{1 - a^T}{1 - a},\]

giving r\textsubscript{1} = 15.41, r\textsubscript{3} = 35.49, r\textsubscript{4} = 41.89, r\textsubscript{5} = 46.66, r\_\ensuremath{\infty} = 60.70 ft.

\emph{Derivation.} Under (F4) the closed-loop residual obeys e\_T =
a\^{}T e\_0. Summing the geometric series of perturbations of magnitude
at most \ensuremath{\varepsilon} yields r\_T = \ensuremath{\varepsilon}(1 + a + \ldots{} + a\^{}\{T\ensuremath{-}1\}) = \ensuremath{\varepsilon}(1 \ensuremath{-}
a\^{}T)/(1 \ensuremath{-} a); the result is the triangle-inequality bound for bounded
additive perturbations of a contracting scalar map. Numerical evaluation
at the listed horizons gives the listed values. \ensuremath{\square}

\textbf{Definition 2.4 (Certified kernel).} The certified kernel at
horizon T is the nominal kernel of \(K^{\ast}\) + r\_T, with r\_T as in
Calculation 2.1.

Two properties of the certified layer must be stated before the results,
because the audit correctly flagged the earlier account as imprecise.
First, in this affine family every policy enters the map only through
the intercept (\ensuremath{\gamma}P is a constant shift), so the contraction rate a is
identical across policies and (F4) is exact, not an approximation: the
certified kernel is the closed-loop nominal kernel of a raised
threshold, feedback included. Second, the mechanism by which the
reactive rules collapse at the certified level is trigger-band entry:
the eroded threshold \(K^{\ast}\) + r\_T crosses the 660-ft trigger
between T = 4 (618 + 41.89 = 659.89) and T = 5 (618 + 46.66 = 664.66),
so any threshold-erosion certificate makes every below-660 trigger rule
identical to training-mean pumping from T \ensuremath{\geq} 4 --- the same
trigger-on-boundary invisibility as Section 3.2, re-entering by
construction. The constant-r\_T form is also conservative by
construction: a per-step tube \(K^{\ast}\) + r\_t (t = 1..T) is tighter
than \(K^{\ast}\) + r\_T held at every step.

\subsection{3. Results}\label{edw-results}

\subsubsection{3.1 Worst-case attractors and the minimal
cut}\label{edw-worst-case-attractors-and-the-minimal-cut}

\textbf{Table 1.} Worst-case attractors of the closed loop (ft), by
policy and recharge floor.

\begin{longtable}[]{@{}lrrr@{}}
\toprule\noalign{}
Policy & UC-min & UC-q05 & UC-q10 \\
\midrule\noalign{}
\endhead
\bottomrule\noalign{}
\endlastfoot
training-mean pumping (BAU) & 615.72 & 625.31 & 626.29 \\
flat-90\% & 618.88 & 628.47 & 629.45 \\
flat-80\% & 622.04 & 631.63 & 632.61 \\
S1 (20\% cut \textless{} 660 ft) & 622.04 & 631.63 & 632.61 \\
flat-70\% & 625.20 & 634.79 & 635.77 \\
CPM cascade & 628.36 & 636.37 & 637.35 \\
flat-60\% & 628.36 & 637.95 & 638.93 \\
flat-50\% & 631.52 & 641.11 & 642.09 \\
flat-0 (zero pumping) & 647.32 & 656.91 & 657.90 \\
\end{longtable}

\begin{figure}[htbp]
\centering
\includegraphics[width=\linewidth]{figs_e4/fig1_attractors.png}
\caption{Worst-case attractor of the closed loop under the perpetual drought-of-record floor (UC-min) by policy, from the registered kernel computation (Table 1). Training-mean pumping's attractor (615.72 ft) sits below the 618-ft physical threshold. Every cut policy --- flat caps of 10\% and deeper, the Stage I reactive rule, and the critical-period-management cascade --- holds its attractor at or above the threshold. Under UC-min the reactive rules coincide with their attractor twins (S1 = flat-80\%, CPM = flat-60\%; under UC-q05 the CPM attractor sits between the Stage III and IV triggers and equals a \textasciitilde65\% cap); flat-90\% clears the threshold by 0.88 ft --- a margin the bootstrap of Section 3.7 shows is not resolved by the fit.}
\end{figure}

Under a perpetual 1956-recharge floor, training-mean pumping's attractor
(615.72 ft) sits below the physical threshold. The smallest flat cut
whose attractor clears 618 ft is an interpolated \textbf{7.2\%} of
training-mean pumping, which lies outside the declared family (on the
current-mean baseline of 382.16 \ensuremath{\times} 10\textsuperscript{3} acre-ft yr\textsuperscript{-1} the same arithmetic
gives a securing cut of about 31.5\%, and the Stage-I 20\% cut does not
secure --- Section 3.4). The Stage-I reactive rule's attractor equals
flat-80\%'s, because the cut is active on the entire attractor branch.

\subsubsection{3.2 Nominal kernels}\label{edw-nominal-kernels}

\textbf{Table 2.} Nominal kernel lower boundary (ft) at the 618-ft
threshold under UC-min, by horizon. Cut policies (flat-90\% and deeper,
S1, CPM) hold the whole safe set {[}618, 710{]} at every horizon,
including the infinite horizon.

\begin{longtable}[]{@{}lrrrrrl@{}}
\toprule\noalign{}
Horizon T & 1 & 5 & 10 & 12 & 12.7 (crossover) & \ensuremath{\infty} \\
\midrule\noalign{}
\endhead
\bottomrule\noalign{}
\endlastfoot
BAU boundary & 618.8 & 625.6 & 658.4 & 692.6 & 710 & empty \\
\end{longtable}

At the physical threshold (618 ft) under UC-min, training-mean pumping's
kernel boundary climbs 618.8 (T = 1) \ensuremath{\rightarrow} 625.6 (T = 5) \ensuremath{\rightarrow} 658.4 (T = 10),
and the kernel is empty beyond T \ensuremath{\approx} 13 years. The emptiness is a
domain-top event, and its arithmetic should be read exactly: the
boundary \(B(T) = H^{\ast} + (618 - H^{\ast})/a^{T}\) reaches the
declared safe-set ceiling when the \emph{required initial head} exceeds
710 ft --- trajectories fall toward the attractor, they do not exit
upward. The horizon is a function of the arbitrary ceiling:
\(T_{\mathrm{empty}}(C) = \mathrm{ln}\bigl((C - H^{\ast})/(K^{\ast} - H^{\ast})\bigr)/\mathrm{ln}(1/a)\),
so the 710-ft ceiling gives 12.7 years, the observed maximum annual mean
(691.96 ft, the 1992 mean) gives 12.0 years, and with no ceiling the
kernel is never empty in the domain sense --- every state converges to
the attractor, and ``empty beyond T'' means ``no initial head in the
domain lasts T years.'' With the ceiling stated, business-as-usual is
not robustly viable against a perpetual drought-of-record, and the
observed 1992 mean (691.96 ft) sits almost at the T = 12 boundary. Every
cut policy in the family --- a 10\% flat cap, S1, CPM, and deeper ---
makes {[}618, 710{]} robustly invariant: the whole declared safe set is
the kernel at every horizon, including the infinite horizon. Under
UC-q05/UC-q10 the safe set is already invariant at training-mean pumping
(attractors 625.3/626.3 ft); governance differentiates the kernel only
under the drought-floor class.

At the institutional threshold (660 ft) the result is a \textbf{design
observation about scoring a trigger against its own level}, and the
kernel statement splits by policy class. For the \textbf{reactive rules}
(S1, CPM) the equality with training-mean pumping is exact at every
horizon: the boundaries (675.1 at T = 1, 695.3 at T = 2 under UC-min;
empty from T = 3) lie strictly above the first-stage trigger (660 ft,
the highest of the four declared trigger levels), no reactive rule fires
in the viable region, and the kernel is policy-invariant for them. That
equality is a geometric identity --- a trigger placed on the boundary of
a constraint is invisible to the viability kernel of that same
constraint, and no rule that cuts only below 660 ft can improve the
invariance of \{H \ensuremath{\geq} 660\} --- and it is a category observation about the
scoring, not a deficiency of the rules: the cascade's institutional
purpose is springflow protection near 618 ft, not head invariance at 660
ft. For the \textbf{flat caps} the equality is false. They move the
finite-horizon boundaries (flat-80\% reads \ensuremath{\approx}673 at T = 1 against BAU's
675.1) and extend the viable horizon (from T = 3 at BAU through T \ensuremath{\approx} 6 at
flat-50\% to T \ensuremath{\approx} 6--11 for zero pumping by the same ceiling arithmetic;
zero pumping's attractor 647.3/657.9 ft still sits below 660 ft). Demand
management therefore extends the viable \textbf{horizon, not
invariance}, and the correct statement of the 660-ft reading is: no
policy in the declared family, zero pumping included, holds \{H \ensuremath{\geq} 660\}
invariant under the drought floors --- the threshold is protected by wet
years, not by the pumping family. This is the frequency-management
rationale the actual CPM rule implements, and that rationale is outside
the robust-kernel frame.

The negative reading is in fact stronger than ``the declared family''
suggests, and the stronger form is worth stating because it is what
makes the result portable. Write the fitted map as \(H_{t+1} = aH_t +
\alpha + \beta R_t + \gamma P_t\) with \(0 < a < 1\) and \(\gamma < 0\)
(here \(a\) = 1 + \(\delta\) = 0.7461 and \(\gamma\) = \ensuremath{-}0.02844, from the
Section 2.2 fit; both are fitted signs, reported rather than assumed). Because \(\gamma < 0\), any
non-negative pumping sequence \(P_t \geq 0\) gives \(H_{t+1} \leq aH_t +
\alpha + \beta F\) pointwise under a recharge floor \(F\); because \(0 < a
< 1\), that comparison map contracts to the zero-pumping attractor
\(H^{\ast}_0 = (\alpha + \beta F)/(1-a)\). So whenever the
\emph{attractor-to-threshold margin} \(m = K - H^{\ast}_0\) is positive,
every non-negative pumping rule --- not merely every member of the
declared family --- fails to hold \(\{H \geq K\}\) invariant. At 660 ft
the margin is positive under all three floors, which is why zero pumping
already fails there.

Two consequences follow. First, the argument quantifies over pumping
\emph{sequences}, so it is indifferent to a rule's functional form: it
covers rules that respond to several state variables at once, and rules
with hysteresis (a stage entered on one indicator and exited only on
another), neither of which the declared family here can express. A
district operating a richer rule set than the one scored here therefore
inherits the same negative verdict. Second, the converse does not hold.
A negative margin says only that zero pumping \emph{could} hold \(K\);
whether some member of a particular declared family does so is a
separate, family-relative question, and one that a larger family can
answer differently. Emptiness generalizes; construction does not.

\subsubsection{3.3 Certified kernels}\label{edw-certified-kernels}

With the erosion of Section 2.4 applied, every positive-pumping policy
in the family has an empty certified kernel beyond \textbf{T = 3 years}
at the physical threshold and beyond T = 1 year at the institutional
threshold. Zero pumping is the exception. Its certified
physical-threshold kernel is nonempty through T = 4 under UC-min (lower
boundary \ensuremath{\approx}687.9 ft from the certified-kernel algebra; the registered
horizon grid \{1, 2, 3, 5, 8, 10, 15, 20, \ensuremath{\infty}\} does not sample T = 4 ---
the algebra is the record there and the grid is annotated accordingly)
and through T = 5 under UC-q05/UC-q10. The certified boundaries at T =
3, UC-min, 618 ft: flat-0 662.2 \textless{} flat-80 697.8 \textless{}
BAU = S1 = CPM 706.7 ft. The certified-level inference is drawn
explicitly: at the one certified horizon tabulated, the reactive rules
inherit training-mean pumping's boundary (706.7) --- strictly above
flat-80\%'s (697.8) --- so certified dominance of S1 over its matched
flat cap fails. The +3.3\%/+16.2\% supply margins are nominal-level
comparisons. Certified retention is a different, horizon-truncated
object from nominal retention (Section 2.4).

The bound is optimistic, and the audit's sensitivity is recorded: at the
out-of-sample defect (\ensuremath{\varepsilon} = 21.81 ft) the erosion margins scale to r\textsubscript{3} =
50.22, r\textsubscript{5} = 66.04, r\_\ensuremath{\infty} = 85.90 ft, so certified emptiness beyond T = 3
is a \emph{lower} bound on how fast certification dies, and certified
horizons do not appear as findings beyond this section's scope
statement. The binding constraint on certified intervention claims is
the model defect, not the governance. The certified comparison is also
conservative in form against the state-dependent rules for the geometric
reason of Section 2.4 --- the eroded threshold enters the trigger band
from T = 4 --- which is a property of threshold-erosion certificates
generally, not of these rules.

\subsubsection{3.4 Supply and retention}\label{edw-supply-and-retention}

\textbf{Table 3.} Mean prescribed pumping: the frozen open-loop replay
(actual observed heads, 1934--1990) and the post-freeze closed-loop
replay (head simulated by the map under actual recharge, same policy in
force), 10\textsuperscript{3} acre-ft yr\textsuperscript{-1}.

\begin{longtable}[]{@{}
  >{\raggedright\arraybackslash}p{(\columnwidth - 8\tabcolsep) * \real{0.1579}}
  >{\raggedleft\arraybackslash}p{(\columnwidth - 8\tabcolsep) * \real{0.2105}}
  >{\raggedleft\arraybackslash}p{(\columnwidth - 8\tabcolsep) * \real{0.2105}}
  >{\raggedleft\arraybackslash}p{(\columnwidth - 8\tabcolsep) * \real{0.2105}}
  >{\raggedleft\arraybackslash}p{(\columnwidth - 8\tabcolsep) * \real{0.2105}}@{}}
\toprule\noalign{}
\begin{minipage}[b]{\linewidth}\raggedright
Policy
\end{minipage} & \begin{minipage}[b]{\linewidth}\raggedleft
Open-loop supply (frozen)
\end{minipage} & \begin{minipage}[b]{\linewidth}\raggedleft
Cut active
\end{minipage} & \begin{minipage}[b]{\linewidth}\raggedleft
Closed-loop supply (post-freeze)
\end{minipage} & \begin{minipage}[b]{\linewidth}\raggedleft
Closed-loop cut active
\end{minipage} \\
\midrule\noalign{}
\endhead
\bottomrule\noalign{}
\endlastfoot
BAU (training-mean) & 282.16 & 0\% & 282.16 & 0\% \\
flat-90\% & 253.94 & 100\% & 253.94 & 100\% \\
flat-80\% & 225.73 & 100\% & 225.73 & 100\% \\
flat-70\% & 197.51 & 100\% & 197.51 & 100\% \\
S1 & \textbf{262.36} & 35.1\% & 258.98 & 41.1\% \\
CPM & \textbf{254.93} & 35.1\% & 254.70 & 41.1\% (any stage) \\
flat-60\% & 169.29 & 100\% & 169.30 & 100\% \\
flat-50\% & 141.08 & 100\% & 141.08 & 100\% \\
flat-0\% & 0.00 & 100\% & 0.00 & 100\% \\
\end{longtable}

Out-of-sample open-loop replay (audit only, over the 1991\ensuremath{\rightarrow}2022
transitions): S1 264.5 and CPM 260.6 \ensuremath{\times} 10\textsuperscript{3} acre-ft yr\textsuperscript{-1}; closed-loop
from the observed 1990 head, S1 prescribes 268.5 with the cut active
24.2\% of years and ends 2023 at 643.3 ft (observed 635.7; training-mean
pumping ends 639.2). Every flat policy prescribes its cap throughout, so
its supplies coincide by construction. The closed loop is a model
object: run from the observed 1934 head it ends 1990 at 677.6 ft against
the observed 645.8 ft (+31.8 ft), so the closed-loop supplies carry the
map's own level bias and are reported as the audit requested --- as the
within-model check on the hybrid margin, not as a second retention leg.

\textbf{Table 4.} The retention margins under every comparator reading
(open-loop supplies; the closed-loop margins follow in the text).

\begin{longtable}[]{@{}
  >{\raggedright\arraybackslash}p{(\columnwidth - 6\tabcolsep) * \real{0.2000}}
  >{\raggedleft\arraybackslash}p{(\columnwidth - 6\tabcolsep) * \real{0.2667}}
  >{\raggedleft\arraybackslash}p{(\columnwidth - 6\tabcolsep) * \real{0.2667}}
  >{\raggedleft\arraybackslash}p{(\columnwidth - 6\tabcolsep) * \real{0.2667}}@{}}
\toprule\noalign{}
\begin{minipage}[b]{\linewidth}\raggedright
Comparator (matched protection at 618 ft / UC-min)
\end{minipage} & \begin{minipage}[b]{\linewidth}\raggedleft
Supply
\end{minipage} & \begin{minipage}[b]{\linewidth}\raggedleft
S1 margin
\end{minipage} & \begin{minipage}[b]{\linewidth}\raggedleft
CPM margin
\end{minipage} \\
\midrule\noalign{}
\endhead
\bottomrule\noalign{}
\endlastfoot
kernel-matched flat-90\% (declared grid) & 253.94 & +3.3\% & +0.4\% \\
attractor twin flat-80\% (S1) / flat-60\% (CPM) & 225.73 / 169.29 &
+16.2\% & +50.6\% \\
interpolated 7.2\% cut (outside family) & 261.84 & +0.2\% & \ensuremath{-}2.6\%
(fails) \\
flat-92\% (1\% grid) & 259.59 & +1.1\% & \ensuremath{-}1.8\% (fails) \\
flat-93\% (1\% grid) & 262.41 & \ensuremath{-}0.02\% (fails) & \ensuremath{-}2.9\% (fails) \\
\end{longtable}

\begin{itemize}
\tightlist
\item
  \textbf{S1: retained (nominal, UC-min, 618 ft; hybrid).} It matches
  the flat caps' robust invariance (kernel = whole safe set at all
  horizons) while supplying 262.36 versus flat-90\%'s 253.94 (+8.4 \ensuremath{\times} 10\textsuperscript{3}
  acre-ft yr\textsuperscript{-1}, +3.3\%) and flat-80\%'s 225.73 (+36.6, +16.2\%). The
  scoring regimes are the hybrid declared in Section 2.4 --- protection
  under the perpetual worst-year recharge floor, supply as the
  1934--1990 historical replay mean. The reading of the retention is
  exactly ``worst-case protection at historical-mean entitlement,'' not
  a within-scenario dominance. Under the floor itself S1 is identically
  flat-80\% (the cut is active every year, and its floor-attractor
  supply equals flat-80\%'s), so the supply margin exists only in wet
  years that the robust class excludes. Against flat-90\% the floor
  comparison is the sharper one --- S1 is always \ensuremath{-}20\% where flat-90\%
  is always \ensuremath{-}10\%, at identical invariance. The reactive architecture
  justifies its additional structure on the hybrid criterion, not on the
  robust criterion alone, and the comparator table shows how
  grid-dependent the margin is: against the 1\%-grid smallest securing
  cap it is +1.1\%, and against the interpolated 7.2\% cut it is +0.2\%.
\item
  \textbf{CPM: retained (nominal, same threshold and class).} Attractor
  628.36 ft (equal to flat-60\%'s) at supply 254.93 versus flat-60\%'s
  169.29 (+50.6\%) and flat-90\%'s 253.94 (+0.4\%). The +50.6\% is
  priced substantially off Stage IV, a stage occupied in 1 of 90 years.
\item
  \textbf{Closed-loop re-check (post-freeze layer).} With each policy in
  force and the head simulated by the map under actual recharge, S1's
  supply falls to 258.98 (\ensuremath{-}1.3\% versus the open-loop replay; the cut is
  active 41.1\% of years, not 35.1\%, because the map's trajectory dips
  below 660 ft more often than the observed record did) and CPM's to
  254.70 (\ensuremath{-}0.1\%). The hybrid margins attenuate but survive: S1 +2.0\%
  versus flat-90\% (+14.7\% versus flat-80\%), CPM +50.4\% versus
  flat-60\% (+0.3\% versus flat-90\%). The margin is a real property of
  the rule's wet-year behavior on this map, not an artefact of
  evaluating the rules on the observed (unregulated-era) heads alone.
\item
  \textbf{Certified level:} certified retention is horizon-truncated. At
  the one tabulated certified horizon the reactive rules inherit
  training-mean pumping's boundary (706.7 ft, Section 3.3) and fail the
  certified re-check. The supply figures are nominal-level comparisons,
  not certified retention.
\item
  \textbf{Institutional threshold: nothing retained.} The reactive rules
  are identically training-mean pumping at every horizon
  (trigger-on-boundary invisibility, Section 3.2). The flat caps are
  more protective at finite horizons and extend the viable horizon. No
  policy meets the retention test there.
\item
  \textbf{Current-baseline sensitivity (post-freeze, analytic).} The
  declared family and every supply figure are relative to the 1934--1990
  training mean (282.16). On the 1991--2023 mean actual pumpage (382.16
  \ensuremath{\times} 10\textsuperscript{3} acre-ft yr\textsuperscript{-1}), the same attractor arithmetic gives
  training-mean-pumping 604.52 ft, the securing cut \ensuremath{\approx}31.5\%, and the
  Stage-I/flat-80\% attractor 613.08 ft --- below 618 ft, so the 20\%
  sketch does not secure the physical threshold at current use.
  Re-expression of the family on the current baseline is a registered
  follow-up, not a scored object here.
\end{itemize}

\subsubsection{3.5 Classification and stress
replays}\label{edw-classification-and-stress-replays}

The T = 5 nominal kernel (UC-min, 618 ft) classifies the actual record.
Training-mean pumping excludes exactly one actual year from its viable
set --- 1956, the drought-of-record year (623.15 ft annual mean, above
618 ft and above the 615.72-ft attractor: 1956 is excluded as a
\emph{starting} head whose 5-year forward path under perpetual 1956
recharge falls through the threshold, not because the 1956 mean was
below 618; its daily minimum, 612.51 ft, was). S1 and CPM exclude none:
the entire 90-year actual record is robustly 5-year viable under the cut
rules. The T = 5 certified kernels are empty, so no actual year is
certified 5-year viable under any policy. This is the boundary of the
certified analysis.

A 1950s open-loop diagnostic (model driven by actual R, P versus actual
heads, 1951--1956) records the map's drought bias. The affine map
under-predicts the decline: model 659.5 \ensuremath{\rightarrow} 631.3 versus actual 659.5 \ensuremath{\rightarrow}
623.2 ft, maximum error 8.1 ft, biased high. The model-based policy
replays from the observed 1950 head keep all policies above 618 ft (BAU
minimum 629.7, S1 634.9, CPM 637.1 ft) --- note that this replay holds
pumping at the training mean (282.16) while actual 1956 pumping was \ensuremath{\approx}321
\ensuremath{\times} 10\textsuperscript{3} acre-ft, so it is not the decade's actual policy either. The
open-loop bias means the true margins are smaller than the replay
suggests. This is recorded, and no correction is applied.

\subsubsection{3.6 Finite-duration recharge floors and floor-class
supply (post-freeze
layers)}\label{edw-finite-duration-recharge-floors-and-floor-class-supply-post-freeze-layers}

Two declared post-freeze layers extend the persistent-floor object. The
first asks how long a drought-class episode must last before the safe
set contracts. The class floor holds for \(n\) years, recharge then
returns to its training mean, and the infinite-horizon lower boundary at
the 618-ft threshold is pulled back by exact backward recursion through
the \(n\) floor years. The second computes the floor-class supply ---
the closed loop from the observed 1934 head with recharge held at each
class floor --- supplying the declared-scoring half that the historical
replay cannot.

\textbf{Table 5.} Finite-duration floors (post-freeze): infinite-horizon
lower boundary (ft) at the 618-ft threshold after \(n\) years at the
class floor followed by training-mean recharge. The BAU row at \(n=5\)
and \(n=10\) reproduces the registered \(T=5\) and \(T=10\) nominal
boundaries (625.6 and 658.4 ft). The \(n=15\) emptiness reproduces the
registered 13-year horizon bound.

\begin{longtable}[]{@{}
  >{\raggedright\arraybackslash}p{(\columnwidth - 8\tabcolsep) * \real{0.1579}}
  >{\raggedleft\arraybackslash}p{(\columnwidth - 8\tabcolsep) * \real{0.2105}}
  >{\raggedleft\arraybackslash}p{(\columnwidth - 8\tabcolsep) * \real{0.2105}}
  >{\raggedleft\arraybackslash}p{(\columnwidth - 8\tabcolsep) * \real{0.2105}}
  >{\raggedleft\arraybackslash}p{(\columnwidth - 8\tabcolsep) * \real{0.2105}}@{}}
\toprule\noalign{}
\begin{minipage}[b]{\linewidth}\raggedright
Policy
\end{minipage} & \begin{minipage}[b]{\linewidth}\raggedleft
UC\_min, \(n{=}5\)
\end{minipage} & \begin{minipage}[b]{\linewidth}\raggedleft
UC\_min, \(n{=}10\)
\end{minipage} & \begin{minipage}[b]{\linewidth}\raggedleft
UC\_min, \(n{=}15\)
\end{minipage} & \begin{minipage}[b]{\linewidth}\raggedleft
UC\_q05/q10, any \(n\)
\end{minipage} \\
\midrule\noalign{}
\endhead
\bottomrule\noalign{}
\endlastfoot
BAU & 625.6 & 658.4 & empty & 618.0 \\
flat-90\% and deeper; S1; CPM & 618.0 & 618.0 & 618.0 & 618.0 \\
\end{longtable}

Duration differentiates exactly one policy. Under the drought-class
floor, training-mean pumping is the only policy whose boundary moves
with episode length. Every cut policy holds the whole safe set at every
duration, because the cuts' floor attractors sit at or above the
threshold (618.9--647.3 ft, Section 3.1, Table 1). A finite drought
episode followed by normal recharge pulls nothing back beyond 618 ft.
Under the two milder classes the safe set is already invariant at BAU
(Section 3.2), and duration is inert for every policy.

\textbf{Table 6.} Floor-class supply (post-freeze): closed loop from the
observed 1934 head (670.4 ft), recharge held at the class floor, mean
prescribed pumping (10\textsuperscript{3} acre-ft yr\textsuperscript{-1}) over the training span
(1934--1990) and the full span (1934--2023); end head at 2023.

\begin{longtable}[]{@{}
  >{\raggedright\arraybackslash}p{(\columnwidth - 8\tabcolsep) * \real{0.1579}}
  >{\raggedleft\arraybackslash}p{(\columnwidth - 8\tabcolsep) * \real{0.2105}}
  >{\raggedleft\arraybackslash}p{(\columnwidth - 8\tabcolsep) * \real{0.2105}}
  >{\raggedleft\arraybackslash}p{(\columnwidth - 8\tabcolsep) * \real{0.2105}}
  >{\raggedleft\arraybackslash}p{(\columnwidth - 8\tabcolsep) * \real{0.2105}}@{}}
\toprule\noalign{}
\begin{minipage}[b]{\linewidth}\raggedright
Policy
\end{minipage} & \begin{minipage}[b]{\linewidth}\raggedleft
UC\_min supply (train/full)
\end{minipage} & \begin{minipage}[b]{\linewidth}\raggedleft
UC\_min end head
\end{minipage} & \begin{minipage}[b]{\linewidth}\raggedleft
UC\_q05 supply (train)
\end{minipage} & \begin{minipage}[b]{\linewidth}\raggedleft
UC\_q10 supply (train)
\end{minipage} \\
\midrule\noalign{}
\endhead
\bottomrule\noalign{}
\endlastfoot
BAU & 282.16 / 282.16 & 615.7 & 282.16 & 282.16 \\
flat-90\% & 253.94 / 253.94 & 618.9 & 253.94 & 253.94 \\
flat-80\% & 225.73 / 225.73 & 622.0 & 225.73 & 225.73 \\
S1 & 226.72 / 226.36 & 622.0 & 226.72 & 226.72 \\
flat-60\% & 169.30 / 169.30 & 628.4 & 169.30 & 169.30 \\
CPM & 174.49 / 172.59 & 628.4 & 187.36 & 187.61 \\
flat-0 & 0 / 0 & 647.3 & 0 & 0 \\
\end{longtable}

The table quantifies the scoring-regime reading of Section 2.4. Under
the floors the reactive rules converge to their matched flat-cap
attractors (S1 to flat-80\%'s 622.04 ft, CPM to flat-60\%'s 628.36 ft
--- the cut active every year on the attractor branch). Their
closed-loop span-mean supply exceeds the matched caps only by the
trigger-lag margin of the descent from the observed 1934 head --- S1
+0.4\% over flat-80\% (226.7 versus 225.7 \ensuremath{\times} 10\textsuperscript{3} acre-ft yr\textsuperscript{-1}, a
floor-class trigger lag), CPM +3.1\% over flat-60\% (174.5 versus
169.3). These are different objects from the hybrid margins of Table 4:
S1's +0.4\% here is the within-floor trigger lag, CPM's +0.4\% there is
the historical-replay margin against flat-90\% --- same number,
different quantities, and neither is a robust-supply statement. The
hybrid criterion's margins exist only in the wet years that the robust
classes exclude; the floor-class supply leg is the margin's own measure
of how little it would buy inside the class.

\subsubsection{3.7 Attractor uncertainty (post-freeze
layer)}\label{edw-attractor-uncertainty-post-freeze-layer}

\textbf{Table 7.} Residual bootstrap of the 1934--1990 fit (5,000
replications, seeded; no replication was non-contractive): worst-case
attractor under UC-min, by policy.

\begin{longtable}[]{@{}
  >{\raggedright\arraybackslash}p{(\columnwidth - 10\tabcolsep) * \real{0.1364}}
  >{\raggedleft\arraybackslash}p{(\columnwidth - 10\tabcolsep) * \real{0.1818}}
  >{\raggedleft\arraybackslash}p{(\columnwidth - 10\tabcolsep) * \real{0.1818}}
  >{\raggedleft\arraybackslash}p{(\columnwidth - 10\tabcolsep) * \real{0.1818}}
  >{\raggedright\arraybackslash}p{(\columnwidth - 10\tabcolsep) * \real{0.1364}}
  >{\raggedleft\arraybackslash}p{(\columnwidth - 10\tabcolsep) * \real{0.1818}}@{}}
\toprule\noalign{}
\begin{minipage}[b]{\linewidth}\raggedright
Policy
\end{minipage} & \begin{minipage}[b]{\linewidth}\raggedleft
Point attractor (ft)
\end{minipage} & \begin{minipage}[b]{\linewidth}\raggedleft
Bootstrap mean (ft)
\end{minipage} & \begin{minipage}[b]{\linewidth}\raggedleft
Bootstrap SD (ft)
\end{minipage} & \begin{minipage}[b]{\linewidth}\raggedright
5th--95th percentile (ft)
\end{minipage} & \begin{minipage}[b]{\linewidth}\raggedleft
P(attractor \ensuremath{\geq} 618 ft)
\end{minipage} \\
\midrule\noalign{}
\endhead
\bottomrule\noalign{}
\endlastfoot
BAU (training-mean) & 615.72 & 613.72 & 13.09 & 590.0 -- 630.6 & 0.41 \\
flat-90\% & 618.88 & 617.00 & 12.46 & 594.5 -- 633.0 & 0.53 \\
flat-80\% / S1 & 622.04 & 620.27 & 11.87 & 598.7 -- 635.7 & 0.65 \\
flat-70\% & 625.20 & 623.55 & 11.35 & 603.1 -- 638.3 & 0.75 \\
flat-60\% / CPM & 628.36 & 626.83 & 10.89 & 607.3 -- 641.2 & 0.83 \\
flat-50\% & 631.52 & 630.11 & 10.51 & 611.3 -- 644.2 & 0.89 \\
flat-0 & 647.32 & 646.49 & 10.03 & 629.0 -- 661.2 & 0.99 \\
\end{longtable}

The audit's resolvability question has a quantitative answer. The
\emph{ordering} of the attractors is resolved: the flat-90\%-minus-BAU
gap has a 90\% bootstrap band of {[}1.9, 5.0{]} ft, excluding zero. The
\emph{threshold clearance} is not: whether training-mean pumping fails
618 ft (P = 0.41) and whether the marginal securing caps (flat-90\%, P =
0.53; S1/flat-80\%, P = 0.65) clear it are coin flips at this fit; only
cuts of 60\% and deeper reach probabilities above 0.8, and zero pumping
reaches 0.99. The nominal-level statement that a 10\% cut secures the
threshold rests on a 2.3-ft gap (615.72 ft to the 618-ft threshold; the
marginal flat-90\% attractor sits at 618.88 ft) on a map with 5.6-ft
residual SD and an acknowledged 8.1-ft high bias in the drought it is
scored on --- the fitted map cannot resolve whether the marginal members
of the family protect 618 ft. The deep cuts and the failure of
training-mean pumping are the resolved parts of the result.

\subsection{4. Discussion}\label{edw-discussion}

The complete evaluation loop --- measured state, calibrated stock-flow
map, declared governance operators, declared uncertainty classes,
viability kernels with explicit declared-defect erosion, a held-out
defect audit, and a fixed retention rule --- yields a \textbf{nominal
selection at the physical threshold under the mildest floor class}. The
reactive architecture matches flat-cap protection at 3.3\% (S1) and
0.4\% (cascade) more permitted supply than the kernel-matched flat-90\%
cap --- 16.2\% and 50.6\% against their attractor twins, +1.1\% and
\ensuremath{-}1.8\% against the 1\%-grid smallest securing cap --- under a hybrid
criterion whose wet-year entitlement the closed-loop re-check attenuates
but does not remove. An interpolated 7.2\% mean cut (outside the
declared family; about 31.5\% of current use) secures the physical
threshold against a perpetual drought-of-record where training-mean
pumping fails within roughly 13 years (a ceiling-relative statement,
Section 3.2). Two boundaries are equally part of the result: the
institutional threshold is not demand-manageable to invariance under the
declared classes --- and, by the margin argument of Section 3.2, not
under any non-negative pumping rule whatever, whatever its functional
form --- and the certified content is defect-bound to T \ensuremath{\leq} 3
years.

The reactive result is system-dependent, not architectural. Its
mechanism is the aquifer's physics. High transmissivity and rapid karst
recharge make wet-year recovery fast, so a flat cap permanently taxes
recoveries that a reactive rule harvests --- the supply margin exists
only in the wet years that the robust classes exclude (Section 3.4). A
companion intervention study on Northern cod applies the identical
design and retains nothing: there the reactive rules cut catch exactly
where the moratorium already protects. The framework's deliverable is
the scored comparison itself; which governance architecture earns its
complexity is a property of the system, and the two scored systems
occupy opposite ends of that design space without bounding it.

Two layers of results must be kept distinct. The robust-kernel findings
are statements about the declared pumping family under the declared
floor classes. The certified-kernel findings are a different layer,
produced by applying the erosion conversion to the declared training
defect. With \ensuremath{\varepsilon} = 15.41 ft --- exceeded out of sample at 21.81 ft ---
every positive-pumping policy's certified kernel is empty beyond T = 3
years, and the bound is optimistic rather than conservative on this
object (Section 3.3). The certified emptiness is a bound on what the
fitted map can certify, not a robustness statement about the physical
system. The mechanism of the reactive rules' certified collapse is
geometric --- the eroded threshold enters the trigger band (Section 2.4)
--- and the certified comparison is conservative in form against
state-dependent rules for that reason.

The 660-ft design observation also reads correctly against the
institution it describes. The actual CPM rule manages \textbf{frequency}
--- the fraction of time the head spends in critical stages --- not
robust invariance. The finding that no pumping rule can hold the
institutional threshold under a perpetual drought floor is consistent
with that design intent: the rule cannot make wet years, it prices dry
ones. The National Academies' review of the Habitat Conservation Plan
(National Research Council 2015) was likewise organized around
monitoring and modeling adequacy under drought. The
horizon-not-invariance distinction is the robust-kernel expression of
the same boundary.

The karst honesty of the object is inherited from its forecast-study
companion and from the modeling lineage it descends from. The one-pool
affine map is the simplest member of the lumped family whose regional
adequacy was tested on the Barton Springs segment (Scanlon et al.~2003;
Hutchison and Hill 2011). Conduits, the Uvalde--San Antonio divide, and
unconfined recharge-zone storage remain in the residual. San Antonio and
Uvalde recharge and pumpage are lumped. The actual CPM triggers are
10-day averages, so the annual-mean rule is a coarse relative of the
real institution. Stages II--IV are declared scenarios, not verified
institutions. Nominal kernels carry no defect margin, and the certified
kernels use a training defect that the out-of-sample audit exceeds. The
counterfactual swap from the largely unregulated fitting era to
state-dependent rules carries the declared econometric boundary of
Section 2.1 (Lucas 1976). The UC floors are certification geometry, not
forecasts. The 1950s replay is biased high by 8.1 ft. Nothing in this
leg promotes or demotes any forecast module, and no two-pool, karst, or
solute claim is made.

\emph{Cross-system reading.} The protocol is the deliverable, and its
value is shown by contrast rather than by a single result. The identical
intervention-selection design applied to Northern cod (Abaee 2026)
retains nothing --- there the reactive rules cut catch exactly where the
moratorium already protects, so no policy earns its complexity. Here, on
the Edwards Aquifer, the reactive architecture earns a \emph{nominal}
supply margin at the physical threshold (+3.3\% Stage I, +0.4\% cascade)
that is real but wet-year-contingent. The two scored systems therefore
occupy opposite ends of the design space --- one where reactive
complexity is unrewarded, one where it is marginally rewarded --- and the
framework does not prescribe which end a new system lands on; it shows
how to find out. That is the general, transferable content: the method
scores protection against the robust floor class and separates the
geometric (kernel/ threshold) from the identified (defect) findings,
independent of the specific system.

\emph{Practical reading.} For a district that already operates triggers,
the result gives two usable bounds. First, at the physical threshold,
deep cuts are robust and the smallest securing cut is interpolated at
7.2\% of training-mean pumping (about 31.5\% of current use) against a
perpetual drought-of-record --- but a cap's nominal supply advantage over
a reactive rule is not robust under the declared floor classes, so a
manager should not credit wet-year replay supply when selecting among
equally protective rules. Second, at the institutional (660-ft) threshold
nothing in the pumping family is protective: the line is priced by
frequency, not held by invariance, so a manager should not expect the
pumping family to substitute for wet years when the cap is evaluated
there. Where certified claims are wanted (a formal finite-horizon
guarantee), the statement is limited to $T \le 3$ years by the fitted
defect (train $\varepsilon = 15.4$ ft, exceeded out of sample at 21.8 ft),
so any longer certified claim should be read as a bound, not a
certificate.

\emph{Generalizability boundary.} The protocol generalizes; the verdicts
do not. Every numerical certificate is scoped to the J-17 annual-mean
one-pool affine map and the declared governance operators and floor
classes, and no transfer to the Uvalde Pool (J-27) is claimed,
consistent with the programme's no-pooling discipline. Conduits, the
Uvalde--San Antonio divide, and unconfined recharge-zone storage remain
in the residual, and the annual-mean trigger is a coarse relative of the
real 10-day rule.

Two senses of generality should be kept apart here, because the margin
argument of Section 3.2 is general in one of them and not the other.
Within \emph{this} calibrated map, the negative result quantifies over
\emph{rules}: a positive margin defeats every non-negative pumping
sequence, so no richer policy vocabulary --- multi-variable triggers,
hysteresis, finer stages --- can recover invariance at 660 ft. That is a
theorem about the map, and it holds whatever family a manager declares.
It says nothing about \emph{other aquifers}: whether any particular pool
has a positive margin at its own threshold depends on that pool's fitted
\((a, \alpha, \beta, \gamma)\) and its own floor, which must be
estimated there and reported with the hypotheses \(0 < a < 1\) and
\(\gamma < 0\) checked rather than assumed. Generality across rules is
established; generality across pools is not claimed.

\subsection{5. Conclusions}\label{edw-conclusions}

On the measured J-17 record, three verdicts, not one. (1)
\textbf{Nominal, physical threshold, mildest floor class:} reactive
trigger rules match flat-cap robust protection of the 618-ft threshold
at 3.3\% (Stage-I sketch) and 0.4\% (cascade) more permitted supply than
the kernel-matched flat-90\% cap --- 16.2\% and 50.6\% against their
attractor twins --- under a hybrid criterion (worst-case protection,
historical-replay entitlement) whose wet-year margin survives a
closed-loop re-check attenuated (S1 +2.0\% versus flat-90\%). An
interpolated 7.2\% mean cut of training-mean pumping (about 31.5\% of
current use) secures the threshold against a perpetual
drought-of-record, where training-mean pumping fails within roughly 13
years (ceiling-relative). (2) \textbf{Certified:} nothing --- every
positive-pumping certified kernel is empty beyond three years, an
optimistic bound on a defect the out-of-sample audit exceeds. (3)
\textbf{Institutional threshold:} nothing --- the 660-ft line is
protected by wet years rather than by the pumping family, a geometric
property of scoring triggers against their own level. The fitted map
resolves the deep-cut protections and the failure of training-mean
pumping, and cannot resolve the marginal caps (P = 0.41--0.65); that is
the honest resolution limit of the deliverable --- the scored comparison
on one measured system.

\section{Cross-system conclusion}
\label{conclusion}

The three parts above certify rather than simulate, on two real assessment records with no hydrology,
no biology and no institutional history in common. Read together they support one claim about the
method and one about its limits, and the second is the more useful.

\emph{What certification buys: the verdict becomes structural.} Part I's central result is that the
no-dominance verdict is \emph{structural, not empirical}. Because the largest catch holding the
reference point from itself is $C^* = g(\mathrm{LRP}) - |e| = 91.59$ kt, and any rule's protection
margin is exactly $C^*$ minus its catch there, harvest and protection are \emph{one budget} at that
point --- and a reactive rule therefore cannot out-supply a flat cap it matches in protection. That
is not a fact about a particular simulation run; it is a fact about two constants, and it holds for
every rule in the class. Part II supplies the same discipline on a harder object: the certified
harvest-free multiplier bracket certifies both collapse steps as harvest-free contractions under
\emph{every} timing convention for the removals accounting, while the reference window's worst step
and the two prelude steps fail to certify. The output is a \emph{typology, not a blanket verdict} ---
which is the honest form of a certification, and the form a simulation cannot produce.

\emph{What it costs: the horizon, and where it comes from.} This is the finding the three parts
establish jointly and none establishes alone. \textbf{Certification has a horizon, and the horizon
is set by the fitted map, not by the method.} In Part I the map is expansive at the reference point
for every admissible carrying capacity $K \ge 2K^*$, so the certified layer has a finite horizon:
six years under the two harsher floors, seven under the informative one. In Part III every
positive-pumping certified kernel is empty beyond three years --- described there as an optimistic
bound on a defect the out-of-sample audit exceeds. The two numbers differ by a factor of two and
come from unrelated systems, and that is the point: the limitation is not an artefact of one
specification. \emph{A certified verdict is valid over a stated horizon, and quoting it without that
horizon is quoting a different claim.} This is the practical discipline certification imposes, and
it is one simulation does not impose --- a simulation shows a trajectory of some length and the
reader is left to supply the horizon.

\emph{What the protocol is for.} All three parts fix the scoring criterion before any score is
computed, and all three report negative results. That combination is what makes a negative result
informative: a rule that fails a criterion fixed in advance has been \emph{tested}, whereas a rule
that fails a criterion chosen afterwards has merely been described. Part III is explicit that its
reactive rules are retained only \emph{nominally}, on a hybrid criterion whose wet-year margin is
exactly what the robust class excludes --- so the nominal retention is reported together with the
reason it is nominal. Part II likewise reports which steps do \emph{not} certify rather than
reporting only the two collapse steps that do. A certification method that only ever produced
positive verdicts would be indistinguishable from advocacy.

\emph{The design consequence, stated generally.} Two results generalise beyond their systems. First,
\emph{where a threshold is protected by the weather rather than by policy, no rule in the scored
family can be credited with protecting it}: Part III finds that the 660-ft line is protected by wet
years, not by the pumping family --- a geometric property of scoring triggers against their own
level, and one that would be invisible to any method that only asked whether the threshold was
breached. Second, \emph{a positive attractor-to-threshold margin defeats every non-negative pumping
rule}, so the binding constraint is geometric rather than institutional. Part I's corresponding
general statement is that harvest and protection are one budget at the reference point, so a gain in
permitted catch is exactly a loss in protection margin and cannot be argued into existence.

\emph{Scope, stated once.} These are scored comparisons on measured systems, and each part states
its own resolution limit. Part III cannot resolve the marginal caps ($P = 0.41$--$0.65$) and says so;
Part II certifies some steps and not others rather than issuing a verdict on the collapse as a whole;
Part I's certified layer is finite-horizon by construction. The general claim is therefore not that
certification resolves more than simulation --- it is that \emph{certification states what it
resolves and over what horizon}, which is the property that makes a verdict usable.

\subsection*{References}
\label{references}
Abaee, A. (2026).

Abaee, A. (2026). An obstruction calculus for viability under
incomplete observation. Submitted manuscript.

Abaee, A. 2026. Does a one-pool water-balance model improve forecasts of
Edwards Aquifer head? A scored test at J-17. Zenodo.
https://doi.org/10.5281/zenodo.22552680. Companion forecast-evaluation
study.

Abaee, A., 2026. Does a surplus-production ladder improve forecasts of
Northern cod? A scored test on NAFO 2J3KL. Zenodo.
https://doi.org/10.5281/zenodo.22553609. Companion forecast-evaluation
study.

Abaee, A., 2026. Governance operators and viability kernels of the
Edwards Aquifer: an intervention-selection test at J-17. Zenodo.
https://doi.org/10.5281/zenodo.22553311. Companion intervention study
(groundwater object).

Abaee, A. 2026. Robust viability of the 2J3KL limit reference point
under a surplus-production map: policy scoring, expansion, and when
catch cannot help. Zenodo. https://doi.org/10.5281/zenodo.22552060.
Companion intervention study (Northern cod, NAFO 2J3KL).

Abaee, A. (2026). The Northern cod forecast ladder under a frozen
specification. Preprint.

Alley, W.M. and Leake, S.A., 2004. The journey from safe yield to sustainability.
\emph{Ground Water} \textbf{42}(1), 12--16.

Alley, W.M., Reilly, T.E. and Franke, O.L., 1999. \emph{Sustainability of Ground-Water
Resources}. U.S.

Aubin, J.-P., 1991. \emph{Viability Theory}.

Aubin, J.-P., 1991. Viability Theory.

Author, C., et al., in preparation. Interval-verified bounds in linear
management templates. Companion methodological study.

B\'en\'e, C., Doyen, L., Gabay, D., 2001. A viability analysis for a
bio-economic model. Ecol. Econ. 36, 385--396.

Birkh\"auser, Boston.

Birkh\"auser, Boston.

Edwards Aquifer Authority. 2024/25. \emph{2023 Groundwater Discharge and
Usage}, Table 1 (after USGS letter report, 5 April 2024).
https://www.edwardsaquifer.org/wp-content/uploads/2025/05/2023-Groundwater-Discharge-and-Usage.pdf

Edwards Aquifer Authority. Critical Period / Drought Management.
https://www.edwardsaquifer.org/groundwater-users/critical-period-drought-management/

Hutchison, W.R., and Hill, M.E. 2011. \emph{Recalibration of the Edwards
BFZ (Barton Springs Segment) Aquifer Groundwater Flow Model}. Texas
Water Development Board, Austin, Texas, 121 p.

Butterworth, D.S., 2007. Why a management procedure approach? Some positives and negatives.
ICES J. Mar. Sci. 64, 613--617.

Calibration data record, Northern cod (NAFO 2J3KL):
sources, cross-checks, and locked extracts. Preprint.

Devlin, J.F. and Sophocleous, M., 2005. The persistence of the water budget myth and its
relationship to sustainability. \emph{Hydrogeology Journal} \textbf{13}, 549--554.

DFO, 2009. A fishery decision-making framework incorporating the
Precautionary Approach.

DFO, 2016. Stock Assessment of Northern cod (NAFO Divs. 2J3KL) in 2016.
DFO Can. Sci. Advis. Sec. Sci. Advis. Rep.~2016/026.
DFO (2016). Stock assessment of Northern cod (NAFO 2J3KL). \emph{Can.
Sci. Advis. Sec. Sci. Advis. Rep.} 2016/026.

Doyen, L., Th\'ebaud, O., B\'en\'e, C., Martinet, V., Gourguet, S., Bertignac,
M., Fifas, S., Blanchard, F., 2012. A stochastic viability approach to
ecosystem-based fisheries management. Ecol. Econ. 75, 32--42.

Edwards Aquifer Authority, 2024. Managing the Edwards Aquifer: critical period management and
index wells.

Edwards Aquifer Authority, San Antonio, TX.

Edwards Aquifer Recovery Implementation Program, 2021. \emph{Habitat Conservation Plan}.

Edwards Aquifer Authority, San Antonio, TX.

Fisheries and Oceans Canada, Ottawa.

Geological Survey Circular 1186, Denver, CO.

Hutchings, J.A., Myers, R.A. (1994). What can be learned from the
collapse of a renewable resource? Atlantic cod, Gadus morhua, of
Newfoundland and Labrador. \emph{Canadian Journal of Fisheries and
Aquatic Sciences} 51, 2126--2146.

Krawczyk, J.B., and Pharo, A. 2013. Viability theory: An applied
mathematics tool for achieving dynamic systems' sustainability.
\emph{Mathematica Applicanda} 41: 97--126.

Krawczyk, J.B., Pharo, A., Serea, O.S., and Sinclair, S. 2013.
Computation of viability kernels: A case study of by-catch fisheries.
\emph{Computational Management Science} 10: 365--396.

Krawczyk, J.B., Pharo, A., Serea, O.S., Sinclair, S., 2013. Computation
of viability kernels: a case study of by-catch fisheries. Comput. Manag.
Sci. 10, 365--396.

Krawczyk, J.B., Pharo, A., 2013. Viability theory: an applied
mathematics tool for achieving dynamic systems' sustainability. Math.
Appl. 41, 97--126.

Lucas, R.E., Jr., 1976. Econometric policy evaluation: A critique.
\emph{Carnegie-Rochester Conference Series on Public Policy} 1: 19--46.

National Research Council. 2015. \emph{Review of the Edwards Aquifer
Habitat Conservation Plan: Report 1}. The National Academies Press,
Washington, DC. https://doi.org/10.17226/21699

Puente, C. 1978. \emph{Method of estimating natural recharge to the
Edwards Aquifer in the San Antonio area, Texas.} U.S. Geological Survey,
Austin, Texas.

Martinet, V., Th\'ebaud, O., Doyen, L., 2007. Defining viable recovery
paths toward sustainable fisheries. Ecol. Econ. 64, 411--422.

Myers, R.A. and Cadigan, N.G. (1995). Was an increase in natural mortality responsible for
the collapse of northern cod? \emph{Can. J. Fish. Aquat. Sci.} \textbf{52}, 1274--1285.

Myers, R.A., Hutchings, J.A. and Barrowman, N.J. (1997). Why do fish stocks collapse? The
example of cod in Atlantic Canada. \emph{Ecol. Appl.} \textbf{7}, 91--106.

Oubraham, A. and Zaccour, G., 2018. A survey of applications of viability theory to the
sustainable exploitation of renewable resources. \emph{Ecological Economics} \textbf{145},
346--367.

Pereau, J.-C., Mouysset, L. and Doyen, L., 2018. Groundwater management in a food security
context. \emph{Environmental and Resource Economics} \textbf{71}(2), 319--336.

Pereau, J.-C., Pryet, A. and Rambonilaza, T., 2019. Optimality versus viability in groundwater
management with environmental flows. \emph{Ecological Economics} \textbf{161}, 109--120.

Punt, A.E., Butterworth, D.S., de Moor, C.L., De Oliveira, J.A.A., Haddon, M., 2016.
Management strategy evaluation: best practices. Fish Fish. 17, 303--334.

Regnier, E., De Lara, M., 2015. Robust viable analysis of a harvested ecosystem model.
Environ. Model. Assess. 20, 687--698.

Regular, P. et al. (2025). Northern cod (NAFO 2J3KL) reassessment.
\emph{DFO Can. Sci. Advis. Sec. Res. Doc.} 2025/048.

Regular, P.M., et al., 2025. Assessment of the Northern Cod stock in
NAFO Divisions 2J3KL in 2024. DFO Can. Sci. Advis. Sec. Res. Doc.
2025/048.

Ricard, D. et al. (2012). Examining the knowledge base and status of
commercially exploited marine species with the RAM Legacy Stock
Assessment Database. \emph{Fish and Fisheries} 13, 380--398.

Scanlon, B.R., Mace, R.E., Barrett, M.E., and Smith, B. 2003. Can we
simulate regional groundwater flow in a karst system using equivalent
porous media models? Case study, Barton Springs Edwards aquifer, USA.
\emph{Journal of Hydrology} 276: 137--158.

Schijns, R. et al. (2021). Reconstructing three centuries of fish and
fisheries catch in the Northwest Atlantic. \emph{ICES Journal of Marine
Science} 78, 2675--2684.

Schijns, R., Froese, R., Hutchings, J.A., Pauly, D., 2021.
Five centuries of cod catches in Eastern Canada. ICES J. Mar.~Sci. 78,
2675--2683.

Sophocleous, M., 2000. From safe yield to sustainable development of water resources --- the
Kansas experience. \emph{Journal of Hydrology} \textbf{235}, 27--43.

Sophocleous, M., 1997. Managing water resources systems: why ``safe yield'' is not
sustainable. \emph{Ground Water} \textbf{35}(4), 561.

Steinemann, A.C. 2003. Drought indicators and triggers: A stochastic
approach to evaluation. \emph{Journal of the American Water Resources
Association} 39: 1217--1233.
https://doi.org/10.1111/j.1752-1688.2003.tb03704.x

Texas Water Development Board.

Water Data for Texas, well 6837203
(J-17). https://waterdatafortexas.org/groundwater/well/6837203

Umphres, G.D., and Choi, N.J. 2025. Estimated Annual Recharge to the
Edwards Aquifer in the San Antonio Area, by Stream Basin or Ungaged
Area, 1934--2024. U.S. Geological Survey data release.
https://doi.org/10.5066/P1BI62NY

Watkins, D.W., Jr., and McKinney, D.C. 1997. Finding robust solutions to
water resources problems. \emph{Journal of Water Resources Planning and
Management} 123: 49--58.

\section*{Declarations}

\subsection*{Data availability}

 The analysis is reproducible: the deterministic layers are closed-form or
fixed-point computations and every stochastic layer uses a fixed seed and a
fixed draw budget (Section 3.8 uses \(N = 20{,}000\) trajectories per cell,
Section 3.10 \(B = 2000\) bootstrap refits). All input data, analysis
scripts, and result files are archived in the public repository at
https://github.com/MIKEAA2020/general-sustainability. The primary kernel
table (Table 1, Sections 3.1--3.5) is produced by
\texttt{wave\_e\_cod/src/run\_intervention\_v3.py}, which writes
\texttt{intervention\_results\_v3.json} and
\texttt{intervention\_boundaries\_v3.csv} in \texttt{results/}; the two post-freeze
reactive families of Definition 2.4 (Table 1) by
\texttt{wave\_e\_cod/src/\allowbreak{}run\_families\_v3.py}, which
writes
\texttt{results/e2\_families\_v3.csv}. Re-executing either runner
regenerates its output files; a
verification re-execution in a fresh environment reproduced every
reported value to nine significant figures; the trailing digits of the
floating-point fields differ by about one part in \(10^{9}\) across BLAS
builds and are not byte-identical. In the source-year convention the flat-180-kt policy's
\(T=\infty\) boundary (Table 1) is empty under the 5th-percentile class,
so no converged fixed point is reported for it; the runner raises the
fixpoint iteration cap to \(20{,}000\) and refuses to return an
unconverged iterate, which an earlier \(300\)-iteration cap did silently
on exactly this row. The critical-zone
rule and cascade vocabulary follows the DFO precautionary-approach
framework (DFO, 2009); the SSB series and LRP are DFO (2016) Table A2;
the catch series is Schijns et al.~(2021). The elevation layers of
Sections 3.7--3.10 (carrying-capacity grid, stochastic viability,
finite-duration floors, bootstrap bands) are produced by
\texttt{wave\_e\_cod/src/campaign\_e2\_elevation\_v3.py} with fixed
seeds, its outputs archived in \texttt{src/results\_srcyear\_v3/}, and
re-execution regenerates them exactly. Figures 1--9 are produced by
\texttt{wave\_e\_cod/src/make\_figs\_v18.py}, which adds the two
structural panels to the seven of
\texttt{wave\_e\_cod/src/make\_figs\_v17.py} without changing them;
\texttt{v17} superseded \texttt{make\_figs\_v16.py} only in
provenance --- it imports the v3 runner and reads the v3 elevation
outputs --- and its output is
identical. Figure 10 is produced by
\texttt{wave\_e\_cod/src/make\_figs\_v19.py}; it re-derives the
horizon curve from the same closed form and refuses to write the panel
unless every tabulated catch reproduces the archived cadence campaign,
so the figure cannot silently drift from the table. The structural
constants of Section 2.4 and the corrected reactive
criterion of Section 3.4 are verified cell by cell by
\texttt{wave\_e\_cod/src/\allowbreak{}campaign\_e2\_structure\_v3.py} (23 checks,
outputs in \texttt{src/results\_struct\_v3/}), and the profile
likelihood, joint bootstrap, observation-error sensitivity and
post-moratorium windows of Sections 3.7, 3.10 and 3.12 by
\texttt{wave\_e\_cod/src/\allowbreak{}campaign\_e2\_identification\_v3.py}
(outputs in \texttt{src/results\_ident\_v3/}); the horizon-in-catch invariance of Table 2 and the form-generality of Section 2.4 by
\texttt{wave\_e\_cod/src/\allowbreak{}campaign\_e2\_cadence\_v3.py} (outputs in \texttt{src/results\_cadence\_v3/}). The three Section 3.6
form rows are produced by
\texttt{campaign\_e2\_depensation\_v3.py} (free-\(s_0\) Allee refit),
\texttt{campaign\_e2\_allee\_declared\_v3.py} (declared-strength refit)
and \texttt{campaign\_e2\_fox\_form\_v3.py} (Fox form), with outputs in
\texttt{src/results\_forms\_v3/}; the Section 3.11 xteNCAM row by
\texttt{arena agent 1/other documents/rerun\_campaigns/\allowbreak{}campaign\_e2\_xteNCAM\_row.py};
and the Section 3.8 one-off treatment of the 1992 draw by
\texttt{e2\_breakpoint\_1992.py} in the same directory. Every script
named above implements the source-year convention of Section 2.3. The
registered-convention counterparts are retained rather than overwritten
(\texttt{run\_intervention\_v2.py} and its outputs, and
\texttt{wave\_e\_cod/src/superseded\_v2/}), so the convention sensitivity
listed in Section 2.3 can be recomputed from the archive rather than
taken on trust.

\subsection*{CRediT authorship contribution statement}

{[}To be completed at submission.{]}

\subsection*{Declaration of competing interest}

The authors declare that they have no known competing financial
interests or personal relationships that could have appeared to
influence the work reported in this paper.

\subsection*{Funding}

This research received no external funding.

\subsection*{Code availability}

The intervention record, the kernel-boundary table and every auxiliary
campaign are produced by the scripts named above, archived with their
outputs. The verification script
\texttt{paperE2\_cod\_intervention\_v29\_verification.py} recomputes the
fit, the floor classes, the kernel tables, the constructive bound, the
certified horizon and the form-comparison rows from the archived campaign
outputs and compares them against every value this manuscript prints; it
is archived alongside this manuscript together with the mutation harness
\texttt{paperE2\_cod\_intervention\_v29\_sabotage.py} and the
numeric-provenance audit
\texttt{paperE2\_cod\_intervention\_v29\_basis\_audit.py}.

\subsection*{AI declaration}
GLM (Z.ai), Qwen (Alibaba Cloud) and DeepSeek AI assisted with drafting
and iterative review. The author reviewed and edited outputs and takes
responsibility for the final work.
\section*{Declarations}

\section*{Funding}
None.

\section*{Competing interests}
None.

\section*{Data availability}
All inputs are the archived public files of the calibration record
(wave-e-cod extracts and the record's CSV files), with primary sources
cited at archival time.

\section*{Code availability}
The verification script
\url{applied_regime_viability_v9_verification.py}
is publicly available at
\url{https://github.com/MIKEAA2020/general-sustainability} (folder
\texttt{arena agent 1/paper rewrites/latex}); it regenerates every
fraction, inequality, bracket, power, ratio, and temporal claim
verbatim from the locked files. The record figure is generated by
\texttt{paper2\_arv\_record\_figure\_v1.py} under the same discipline:
every plotted series value, threshold, crossing, and annotation is read
from the locked files and asserted exactly before the figure is written
(eleven checks).

\section*{AI declaration}
GLM (Z.ai), Qwen (Alibaba Cloud) and DeepSeek AI assisted with drafting
and iterative review. The author reviewed and edited outputs and takes
responsibility for the final work.
\section*{Declarations}

\subsection*{Data Availability Statement}

All input data, analysis scripts, and result files are archived in the
public repository at
https://github.com/MIKEAA2020/general-sustainability. J-17 daily highs:
Texas Water Development Board well 6837203. Recharge: Umphres and Choi
(2025), USGS data release, https://doi.org/10.5066/P1BI62NY. Pumpage:
Edwards Aquifer Authority (2024/25), Table 1. The full intervention
protocol --- object, defect declaration, disturbance classes, policy
family, and retention rule --- was frozen and dated 2026-08-26 before
any kernel, boundary, replay, or retention score was computed. It is
archived with the analysis code as the frozen-protocol record, alongside
the companion forecast-evaluation protocols dated 2026-08-25. The
analysis is fully deterministic: re-executing the registered runner
regenerates both output files, and a verification re-execution in a
fresh environment reproduced both byte for byte. The machine-readable
outputs include the nominal and certified retention fields and the
certified-horizon record. The Section 3.6 layers (finite-duration
floors; floor-class supply) are produced by
\texttt{rerun\_campaigns/campaign\_e4\_elevation.py}, archived alongside
their outputs, and regenerate them exactly. The post-freeze audit layers
of Sections 3.4 (closed-loop supply replay), 3.4's current-baseline
sensitivity, and 3.7 (attractor bootstrap) are produced by
\texttt{wave\_e\_edwards/src/e4\_audit\_layer.py} (seeded,
deterministic, on the registered panel and the paper's declared
coefficients), with outputs archived as
\texttt{wave\_e\_edwards/results/e4\_audit\_layer.json}.

\subsection*{Funding}
None.

\subsection*{Competing interests}
None.

\subsection*{Code availability}
The kernel computation and the policy scoring are produced by the scripts
named above, archived with their outputs. The verification script
\texttt{paperE4\_edwards\_intervention\_v16\_verification.py} re-runs
the intervention computation from the registered panel and recompares
every attractor of Table~1, the smallest securing cut, the certified
horizons, and the fit diagnostics against the recomputed values; it is
archived alongside this manuscript.

\subsection*{AI declaration}
GLM (Z.ai), Qwen (Alibaba Cloud) and DeepSeek AI assisted with drafting
and iterative review. The author reviewed and edited outputs and takes
responsibility for the final work.

\end{document}
